\documentclass{article}

\usepackage{iclr2027_conference,times}
\usepackage{amsmath,amssymb,amsthm,mathtools}
\usepackage{microtype}
\usepackage{float}
\usepackage{xcolor}
\setlength{\marginparwidth}{1.2in}
\setlength{\marginparsep}{0.15in}
\usepackage{todonotes}
\usepackage{aliascnt}
\usepackage[colorlinks=true,linkcolor=blue!55!black,citecolor=blue!55!black,urlcolor=blue!55!black,hypertexnames=false]{hyperref}
\usepackage[capitalise,noabbrev]{cleveref}

\newtheorem{theorem}{Theorem}
\newaliascnt{proposition}{theorem}
\newtheorem{proposition}[proposition]{Proposition}
\aliascntresetthe{proposition}
\newaliascnt{lemma}{theorem}
\newtheorem{lemma}[lemma]{Lemma}
\aliascntresetthe{lemma}
\newaliascnt{corollary}{theorem}
\newtheorem{corollary}[corollary]{Corollary}
\aliascntresetthe{corollary}

\usepackage{thm-restate} % load after the aliascnt theorem declarations
\usepackage{autonum} % load last

\crefname{equality}{equality}{equalities}
\Crefname{equality}{Equality}{Equalities}
\crefname{inequality}{inequality}{inequalities}
\Crefname{inequality}{Inequality}{Inequalities}
\newcommand{\eqlabel}[1]{\label[equality]{#1}}
\newcommand{\ineqlabel}[1]{\label[inequality]{#1}}
\newcommand{\proofstep}[2]{%
    \par\medskip
    \noindent
    \textcolor{blue!55!black}{\rule{2pt}{1.2em}}%
    \hspace{0.6em}%
    {\sffamily\bfseries\color{blue!55!black}Step #1}%
    \quad
    {\sffamily #2}%
    \par\smallskip
}
\newcommand{\proofstepinline}[1]{%
    {\normalfont\normalsize\bfseries\color{blue!55!black}Step #1.}\enspace
}
\newcommand{\reason}[1]{\texttt{(#1)}}
\newcommand{\explan}[2]{\stackrel{\reason{#2}}{#1}}
\newcommand{\R}{\mathbb R}
\newcommand{\E}{\mathbb E}
\newcommand{\Pp}{\mathbb P}
\newcommand{\K}{\mathcal K}
\newcommand{\cP}{\mathcal P}
\newcommand{\F}{\mathcal F}
\newcommand{\N}{\mathcal N}
\newcommand{\T}{\mathcal T}
\newcommand{\tr}{\operatorname{tr}}
\newcommand{\Lip}{\operatorname{Lip}}
\newcommand{\op}{\mathrm{op}}
\newcommand{\argmin}{\operatornamewithlimits{arg\,min}}
\newcommand{\1}{\mathbf 1}
\newcommand{\eps}{\varepsilon}
\renewcommand{\epsilon}{\varepsilon}
\newcommand{\ip}[2]{\left\langle #1,#2\right\rangle}
\newcommand{\norm}[1]{\left\lVert #1\right\rVert}
\newcommand{\normm}[2]{\left\lVert #1\right\rVert_{#2}}
\newcommand{\alireza}[1]{\todo[color=blue!20,size=\tiny]{Alireza: #1}}
\newcommand{\kevin}[1]{\todo[color=green!20,size=\tiny]{Kevin: #1}}
\newcommand{\lillian}[1]{\todo[color=purple!20,size=\tiny]{Lillian: #1}}

\floatstyle{ruled}
\newfloat{algorithm}{tbp}{loa}
\floatname{algorithm}{Algorithm}

\title{Direct Quadratic Estimation for Stochastic Bandit Convex Optimization}
\author{\qquad}
\date{}

\begin{document}
\maketitle

\begin{abstract}
    We consider stationary stochastic bandit convex optimization with a bounded
    convex loss and subgaussian observation noise.  Gaussian online Newton step
    methods are usually analyzed through an optimistic Gaussian surrogate whose
    derivatives are estimated using a likelihood-ratio correction.  We show that
    neither device is needed in the stationary problem.  The algorithm estimates
    the first two Gaussian Hermite coefficients of a queryable convex extension
    directly from the observed loss and applies online Newton step to the resulting
    (damped) quadratic estimates, without projecting the mean.
    For a rounded convex body $\K$, the regret is
    \begin{align}
        O\!\left(d^2 M\sqrt n\,
        \log^8\!\frac{nd}{\delta}\right)
    \end{align}
    with probability at least $1-\delta$, where $M$ is $\tilde{O}(d^{-1/2})$
    when $\frac1d \K$ is in L\"owner position,
    which recovers the stationary rate of \citet{fokkema2024online}
    of $\widetilde O(d^{3/2}\sqrt n)$.
\end{abstract}

\section{Introduction}

The learner in bandit convex optimization chooses a point in a convex body
and observes only a noisy value of the loss at that point.  A useful way to
extract second-order information is to sample from a Gaussian and estimate the
first two coefficients in the corresponding Hermite expansion.  The first
coefficient behaves like a gradient and the second like an average Hessian, so
the natural update is online Newton step.

There is one difficulty with this direct idea.  The online Newton step
potential contains the curvature of the quadratic at the comparator, while
convexity controls regret through a Bregman divergence of the original loss.
These quantities need not be comparable globally.  The usual Gaussian ONS
analysis resolves this by first replacing the loss with an optimistic
Gaussian surrogate and then estimating its derivatives with a change of
measure; see \citet{lattimore2023second} and Chapters~10--12 of
\citet{lattimore2026bandit}.  Here we take a
different route for the stationary stochastic problem.  We compare the direct
Hermite curvature with the Bregman divergence on an ellipsoidal region around
the Gaussian mean.  The comparison loses only a small part of the usual
log-determinant trace budget.

The resulting algorithm has one Gaussian sample and one function evaluation
per round, with no projection of the mean.
The proof still has to address two features of bandit feedback:
the one-sample Hessian estimate is indefinite, and the queryable extension
amplifies noise outside the constraint set.
We handle the first by comparing the random covariance with its
predictable counterpart, and the second by bounding the extension value at the means and closing
a feedback-moment budget through extension regret.

The paper is organized as follows.  \Cref{sec:setting} introduces the problem,
the extension and the algorithm.
The stability and regret analysis, including the trace, moment and
concentration estimates, appears in \cref{sec:stability}, with the
supporting proofs deferred to the appendices.

\section{Problem formulation and algorithm}
\label{sec:setting}

For $r>0$, let $\mathbb{B}_r=\{x\in\R^d:\norm{x}\leq r\}$
and $\mathbb{S}_r=\{x\in\R^d:\norm{x}=r\}$ be the Euclidean ball and sphere of radius $r$.
Let $\K\subset\R^d$ be a convex body.  An affine preprocessing step puts the body
in the position $\mathbb{B}_1\subseteq\K\subseteq \mathbb{B}_{2d}$.
The loss $f:\K\to[0,1]$ is convex and does not change with time.  In round
$t$, the learner chooses $A_t\in\K$ and observes $f(A_t)+\xi_t$
where $\xi_t$ is the noise in round $t$.
Let $\zeta_t$ be the randomness of the algorithm in round $t$ which it uses to make decisions, and let
$\F_t=\sigma(\zeta_1,\xi_1,\ldots,\zeta_t,\xi_t)$ be the filtration at the end of round $t$,
and $\E_t[\cdot]=\E[\cdot\mid\F_{t-1}]$ and $\Pp_t(\cdot)=\Pp(\cdot\mid\F_{t-1})$ be the conditional expectation and probability given the history.
The noise is conditionally subgaussian, i.e., for all $t$,
\begin{align}
    \E_t[\xi_t|\zeta_t]=0,
    \qquad
    \E_t[\exp(\xi_t^2)|\zeta_t]\leq2.
    \label{eq:noise}
\end{align}
The quantity of interest is the realized (\textit{pseudo-})regret
\begin{align}
    \operatorname{Reg}_n
    =\sum_{t=1}^n\left(f(A_t)-\inf_{x\in\K}f(x)\right).
    \label{eq:regret}
\end{align}

Let $\pi(x)=\inf\{r>0:x\in r\K\}$ be the Minkowski functional of $\K$ and let
$U$ be uniform on the Euclidean unit sphere.
The only geometric quantity in the result is
\begin{align}
    M=\max\left\{\frac1{\sqrt d},\E[\pi(U)]\right\}.
    \label{eq:M}
\end{align}
The rounding  $\mathbb{B}_1\subseteq\K$ implies that $\pi$ is $1$-Lipschitz.
Write $\gamma_{\mu,\Sigma}$ for the density of
$\N(\mu,\Sigma)$, for $\mu\in\R^d$ and $\Sigma\in\R^{d\times d}$ symmetric positive semidefinite.

We write $Dh$ and $D^2h$ for the first and second distributional derivatives of $h$;
see Appendix~\ref{app:distributional-derivatives} for their definitions.

\subsection{A queryable convex extension}

Fix $\eps\in(0,1/2]$ and write $\K_\eps=(1-\eps)\K$ and
\begin{align}
    \pi_\eps(x)
     & =\max\left\{1,\frac{\pi(x)}{1-\eps}\right\}
    =\begin{cases}
         1               & \text{if }x\in\K_\eps, \\
         \pi(x)/(1-\eps) & \text{otherwise}\,,
     \end{cases}
    \quad\text{and}\quad v(x)=\pi_\eps(x)-1\,.
    \label{eq:extension-notation}
\end{align}
We use the queryable extension from
\citet[Eq.~(10.6)]{lattimore2026bandit} which is defined for all $x\in\R^d$ by
\begin{align}
    e(x)
    =\pi_\eps(x)f\!\left(\frac{x}{\pi_\eps(x)}\right)
    +\frac{2}{\eps}v(x)\,,
    \quad\text{and satisfies}
    \Lip(e)\leq\frac{8}{\eps}\,.
    \label{eq:extension-lipschitz}
\end{align}
Moreover, the extension $e$ is convex and agrees with $f$ on $\K_\eps$.
These properties are established in \cref{app:extension}.

% \subsection{Gaussian $L^2$ projection and Hermite basis}
% \label{sec:hermite-projection}
% The Gaussian space $L^2(\N(\mu,\Sigma))$ consists of the functions that are
% square-integrable under the law of $X \sim \N(\mu, \Sigma)$, i.e.
% $L^2(\N(\mu,\Sigma))
%     =
%     \left\{
%     r:\R^d\to\R:
%     \E[r(X)^2]<\infty
%     \right\}$.
% Since $e$ is Lipschitz, we have
% \begin{align}
%     \E\left[
%         e(X)^2
%         \right]
%     \leq
%     \E\left[
%         (e(0)+\Lip(e)\norm{X})^2
%         \right]
%     \leq
%     2e(0)^2 + 2\Lip(e)^2\E[\norm{X}^2] < \infty,
% \end{align}
% which implies that $e\in L^2(\N(\mu,\Sigma))$.
% Next, define
% \begin{align}
%     \mathcal P_{\mu,\Sigma}e(x)
%     =
%     \E[e(X)]
%     +\ip{\bar{g}}{x-\mu}
%     +\frac12
%     \left(
%     \|x-\mu\|^2_{\bar{H}}
%     -\tr(\bar{H}\Sigma)
%     \right).
%     \label{eq:hermite-best-approximation}
% \end{align}
\section{The algorithm}
On an active round $t$, the learner samples
$X_t\sim \N(\mu_t,\Sigma_t)$ and plays its radial contraction
$A_t=X_t/\pi_\eps(X_t)\in\K$.
The resulting feedback is $f(A_t) + \xi_t$ which is converted into the feedback from the queryable extension $e$ as
\begin{align}
    Y_t
     & =
    \pi_\eps(X_t)\bigl(f(A_t)+\xi_t\bigr)
    +\frac2\eps\bigl(\pi_\eps(X_t)-1\bigr)
    =e(X_t)+\pi_\eps(X_t)\xi_t\,,
    \label{eq:Y}
\end{align}
so $\E_t[Y_t\mid X_t]=e(X_t)$.
Let $\lambda \in (0,1)$ be a damping parameter and define the quadratic surrogate
\begin{align}
    \bar{q}_t(x)
    =
    \ip{\bar{g}_t}{x-\mu_t}
    +\frac\lambda2\normm{x-\mu_t}{\bar{H}_t}^2,
    \qquad
    \bar{g}_t\triangleq \langle De,\gamma_{\mu_t,\Sigma_t}\rangle,
    \qquad
    \bar{H}_t\triangleq \langle D^2e,\gamma_{\mu_t,\Sigma_t}\rangle.
\end{align}
The main observation is that $\bar{g}_t$ and $\bar{H}_t$ can be estimated from the single pair $(X_t, Y_t)$.
More precisely, define
\begin{align}
    q_t(x)
    \triangleq
    \ip{g_t}{x-\mu_t}
    +
    \frac\lambda2(x-\mu_t)^\top H_t(x-\mu_t)\,,
\end{align}
where for $Z_t\triangleq\Sigma_t^{-1/2}(X_t-\mu_t)$, we define
\begin{align}
    g_t
     & \triangleq - \frac{Y_t \nabla \gamma_{\mu_t,\Sigma_t}(X_t)}{\gamma_{\mu_t,\Sigma_t}(X_t)}
    =
    Y_t\Sigma_t^{-1/2}Z_t,
     &
    H_t \triangleq
    \frac{Y_t \nabla^2 \gamma_{\mu_t,\Sigma_t}(X_t)}{\gamma_{\mu_t,\Sigma_t}(X_t)}
     & =Y_t\Sigma_t^{-1/2}(Z_tZ_t^\top-I_d)\Sigma_t^{-1/2}\,,
    \label{eq:sampled-hermite-coefficients}
\end{align}
then from distributional integration by parts we get
\begin{align}
    \bar{g}_t & = \langle De, \gamma_{\mu_t,\Sigma_t}\rangle
    = - \langle e, \nabla \gamma_{\mu_t,\Sigma_t}\rangle
    = \E_t[g_t],
    \quad\text{and}\quad\nonumber                              \\
    \bar{H}_t & = \langle D^2e, \gamma_{\mu_t,\Sigma_t}\rangle
    = \langle e, \nabla^2 \gamma_{\mu_t,\Sigma_t}\rangle
    = \E_t[H_t]\,,
    \label{eq:sampled-hermite-unbiasedness}
\end{align}
and more importantly, $\E_t[q_t]=\bar{q}_t$.

The algorithm is just unprojected online newton steps (ONS) applied to quadratic estimates $q_t$.
If $\Sigma_t^{-1}+\lambda\eta H_t\succ0$,
then completing the square gives the equivalent forms
\begin{align}
    \mu_{t+1}
     & =\argmin_{\mu\in\R^d}
    \left\{
    \frac12\normm{\mu-\mu_t}{\Sigma_t^{-1}}^2
    +\eta q_t(\mu)
    \right\}
    =\mu_t-\eta\Sigma_{t+1}g_t.
    \label{eq:projected-ons-update}
\end{align}
\begin{algorithm}[!t]
    \caption{Direct Hermite--ONS algorithm}
    \label{alg:direct-hermite}
    \small
    \begin{tabbing}
        \qquad\=\qquad\=\kill
        \textbf{input:} $\eta,\lambda,\sigma^2,\eps$\\
        $\mu_1=0$, $\Sigma_1=\sigma^2I_d$ and $\mathsf{safe}=\mathsf{true}$\\[1mm]
        \textbf{for} $t=1,\ldots,n$ \textbf{do}\\
        \>Sample $X_t\sim\N(\mu_t,\Sigma_t)$\\
        \>Play $A_t=X_t/\pi_\eps(X_t)$ and observe $f(A_t)+\xi_t$\\
        \>Form $Y_t=e(X_t) + \pi_\eps(X_t) \xi_t$\\
        \>Form $g_t = Y_t\Sigma_t^{-1/2}Z_t$ and $H_t = Y_t\Sigma_t^{-1/2}(Z_tZ_t^\top-I_d)\Sigma_t^{-1/2}$\\
        %  and set $V_{t+1}=\Sigma_t^{-1}+\lambda\eta H_t$\\
        % \>\textbf{if} $\Sigma_t^{-1}+\lambda\eta H_t\succ0$:\\
        \>Set $\Sigma_{t+1}^{-1}=\Sigma_t^{-1}+\lambda\eta H_t$ and $\mu_{t+1} =\mu_t-\eta\Sigma_{t+1}g_t$\\
        % \>\textbf{else}:\\
        % \>\>Set $\Sigma_{t+1}=\Sigma_t$ and $\mu_{t+1}=\mu_{t}$ \qquad\textit{\(\triangleright\) The algorithm has failed.}\\
        \textbf{end for}
    \end{tabbing}
\end{algorithm}

\subsection{Main result}

\begin{theorem}
    \label{thm:main}
    For $\delta\in(0,1/2)$, let
    \begin{align}
        \ell    & =128\log\frac{128nd}{\delta},
                &
        \sigma  & =\frac1{M\sqrt{2d}},
                &
        \lambda & =\frac{1}{1000d^{3/2}M\ell^4},
                &
        \eta    & =\frac{8dM}{\sqrt n},
                &
        \eps_0  & =\frac{100d^2M\ell^8}{\sqrt n}.
        \label{eq:parameters-one}
        \label{eq:parameters-two}
    \end{align}
    If $\eps_0\leq1/2$, run \cref{alg:direct-hermite} with $\eps=\eps_0$.
    Otherwise play any fixed point in $\K$.  In the setting above,
    \begin{align}
        \operatorname{Reg}_n
        \leq 200d^2M\ell^8\sqrt n
        \label{eq:main-regret}
    \end{align}
    with probability at least $1-\delta$.
\end{theorem}

\begin{corollary}
    \label{cor:loewner}
    If $\K/d$ is in L\"owner position, then
    \begin{align}
        \operatorname{Reg}_n
        =O\!\left(
        d^{3/2}\sqrt{\log(ed)}\,\ell^8\sqrt n
        \right)
        =\widetilde O(d^{3/2}\sqrt n)
    \end{align}
    with probability at least $1-\delta$.
\end{corollary}

\begin{proof}
    In L\"owner position,
    $\E[\pi(U)]=O(\sqrt{\log(ed)/d})$; see
    \citet[Table~10.1(a)]{lattimore2026bandit}.  Substitute this bound in
    \eqref{eq:main-regret}.
\end{proof}

\section{Regret switching}
\label{sec:analysis}
We fix a time step $t$ and drop the subscript $t$ from the remaining notation in this subsection.
The proof rests on relating the regret of the quadratic surrogate $\bar{q}$ to the regret of the original loss $e$.
For a convex Lipschitz function $e$, and a differentiable point $z$, the Bregman divergence is
\begin{align}
    D_e(x\,\|\,z)
    =e(x)-e(z)-\ip{\nabla e(z)}{x-z}.
    \label{eq:bregman}
\end{align}

\begin{corollary}[Regret switching]
    \label{cor:regret-switch}
    For
    $\bar{q}(x)=\ip{\bar{g}}{x-\mu}+\frac\lambda2\normm{x-\mu}{\bar{H}}^2$,
    \begin{align}
        \E[e(X)]-e(x)
        =
        \bigl(\bar{q}(\mu) - \bar{q}(x)\bigr)
        +\tr(\bar{H}\Sigma)
        +\frac\lambda2\normm{x-\mu}{\bar{H}}^2
        -\E[D_e(x\,\|\,X)].
        \label{eq:q-regret-switch}
    \end{align}
\end{corollary}
This identity follows from \cref{lem:regret-switch} in Appendix~\ref{app:regret-switch} and the definition of $\bar{q}$.
To transfer a bound on quadratic regret to the expected loss regret
$\E[e(X)]-e(x)$, it remains to upper bound
$\lambda\normm{\mu-x}{\bar{H}}^2/2-\E[D_e(x\,\|\,X)]$.
The following corollary controls this difference on the focus region.
It is obtained from the general comparison in \cref{prop:bracket} using the parameters in \cref{thm:main}.
The general result and this specialization are proved in Appendix~\ref{app:bracket}.

\begin{corollary}[Curvature comparison]
    \label{cor:bracket-ell}
    % Let $\ell$, $\lambda$ and $\sigma$ be as in \cref{thm:main}, and let $\eps>0$.
    Let $e:\R^d\to\R$ be convex and
    $(8/\eps)$-Lipschitz, and let $X\sim\N(\mu,\Sigma)$ with
    $0\prec\Sigma\preceq2\sigma^2I_d$.
    For every $x\in\R^d$ satisfying
    \begin{align}
        \frac12\normm{\mu-x}{\Sigma^{-1}}^2
        \leq\frac{1}{128^2\lambda^2\ell},
        \label{eq:bracket-ell-focus}
    \end{align}
    it holds that
    \begin{align}
        \frac\lambda2\normm{\mu-x}{\bar{H}}^2
        -\E[D_e(x\,\|\,X)]
        \leq\frac1{16}\tr(\bar{H}\Sigma)
        +\frac{n^{-8}}{\eps}.
        \label{eq:bracket-ell}
    \end{align}
\end{corollary}

Together, \cref{cor:regret-switch,cor:bracket-ell} bound convex regret by
quadratic regret, a trace term, and a negligible error, provided the comparator
stays in the focus ellipsoid. Restoring the time indices, if
$0\prec\Sigma_t\preceq2\sigma^2I_d$ and
$\frac12\normm{\mu_t-x}{\Sigma_t^{-1}}^2\leq\frac{1}{128^2\lambda^2\ell}$, then
\begin{align}
    \E_t[e(X_t)]-e(x)
    \leq \bar{q}_t(\mu_t)-\bar{q}_t(x)
    +\frac{17}{16}\tr(\bar{H}_t\Sigma_t)
    +\frac{n^{-8}}{\eps}.
    \label{eq:time-indexed-regret-switch}
\end{align}

\section{Population-level Analysis}
\label{sec:proof-sketch}
Let's first worry about the expected convex regret of the extension,
even though this is not quite our final target. Let $x_\star$ minimize
$f$ on $\K_\eps$, and imagine that, in every round:
\begin{center}
    \textcolor{blue!55!black}{\textbf{(1)}} We observe $\bar{g}_t,\bar{H}_t$ exactly
    \qquad
    \textcolor{blue!55!black}{\textbf{(2)}}
    $\normm{\mu_t-x_\star}{\bar{\Sigma}_t^{-1}}<\sqrt{2B}$.
\end{center}

With these two conveniences, \eqref{eq:time-indexed-regret-switch}
already gets us started:
\begin{align}
    \operatorname{eReg}_n \triangleq \sum_{t=1}^n\bigl(\E_t[e(X_t)]-e(x_\star)\bigr)
     & \leq\sum_{t=1}^n\bigl(\bar{q}_t(\mu_t)-\bar{q}_t(x_\star)\bigr)
    +\frac{17}{16}\sum_{t=1}^n\tr(\bar{H}_t\bar{\Sigma}_t)+\frac{n^{-7}}\eps.
\end{align}

We have traded a general convex loss for quadratic losses, which makes
life easier. The catch is that $\bar{q}_t$ now changes from round to round,
even though $e$ is fixed. Thankfully, this is where ONS shines: its
regret identity works even for adversarial sequences of quadratics.
Starting from $\mu_1=0$ and $\bar{\Sigma}_1=\sigma^2I_d$, apply ONS to $\bar{q}_t$:
\begin{align}
    \bar{\Sigma}_{t+1}^{-1}
    =\bar{\Sigma}_t^{-1}+\lambda\eta \bar{H}_t,\qquad
    \mu_{t+1}
    =\argmin_{x\in\R^d}
    \left\{\frac12\normm{x-\mu_t}{\bar{\Sigma}_t^{-1}}^2+\eta \bar{q}_t(x)\right\}
    =\mu_t-\eta\bar{\Sigma}_{t+1}\bar{g}_t.
\end{align}
To see what this buys us, track the squared distance to $x_\star$.
Expanding the update, substituting the covariance update, and then
recognizing $\bar{q}_t$ gives
\begin{align}
    \frac12\normm{\mu_{t+1}-x_\star}{\bar{\Sigma}_{t+1}^{-1}}^2
     & =\frac12\normm{\mu_t-x_\star-\eta\bar{\Sigma}_{t+1}\bar{g}_t}
    {\bar{\Sigma}_{t+1}^{-1}}^2                                      \\
     & =\frac12\normm{\mu_t-x_\star}{\bar{\Sigma}_{t+1}^{-1}}^2
    +\eta\ip{\bar{g}_t}{x_\star-\mu_t}
    +\frac{\eta^2}{2}\normm{\bar{g}_t}{\bar{\Sigma}_{t+1}}^2         \\
     & =\frac12\normm{\mu_t-x_\star}{\bar{\Sigma}_t^{-1}}^2
    +\frac{\lambda\eta}{2}\normm{\mu_t-x_\star}{\bar{H}_t}^2
    +\eta\ip{\bar{g}_t}{x_\star-\mu_t}
    +\frac{\eta^2}{2}\normm{\bar{g}_t}{\bar{\Sigma}_{t+1}}^2         \\
     & =\frac12\normm{\mu_t-x_\star}{\bar{\Sigma}_t^{-1}}^2
    +\eta\bigl(\bar{q}_t(x_\star)-\bar{q}_t(\mu_t)\bigr)
    +\frac{\eta^2}{2}\normm{\bar{g}_t}{\bar{\Sigma}_{t+1}}^2.
    \label{eq:ons-regret-identity}
\end{align}
Rearrange to isolate $\bar{q}_t(\mu_t)-\bar{q}_t(x_\star)$ and sum over $t$.
The successive squared distances telescope; dropping the negative
final-distance term leaves exactly the quadratic-regret bound we need:
\begin{align}
    \sum_{t=1}^n\bigl(\bar{q}_t(\mu_t)-\bar{q}_t(x_\star)\bigr)
    \leq\frac{\normm{\mu_1-x_\star}{\bar{\Sigma}_1^{-1}}^2}{2\eta}
    +\frac\eta2\sum_{t=1}^n\normm{\bar{g}_t}{\bar{\Sigma}_{t+1}}^2.
\end{align}

\section{Stability and regret}
\label{sec:stability}
\label{sec:stochastic-control}
Let $x_\star \in\argmin_{x\in\K_\eps}f(x)$ be a minimizer of $f$ on the contracted set $\K_\eps$.
Also in all that follows, we assume that $\eps=\eps_0\leq1/2$ and that the parameters $\sigma,\lambda,\eta$ are as in \cref{thm:main}.
If $\eps>1/2$, then ....

\proofstepinline{0}
We first bound realized regret in terms of extension regret, as in \citet{lattimore2026bandit}[Section~10.6]. The following lemma suggests that bounding the extension regret is sufficient to bound the realized regret.
\begin{lemma}[{Adapted from \citet[Section~10.6, Eq.~(10.14)]{lattimore2026bandit}}]
    \label{lem:actual-to-extension}
    Let $\mathcal E_{\mathrm{reg}}$ be the event that
    \begin{align}
        \operatorname{Reg}_n
        \leq
        \operatorname{eReg}_n
        +n\eps+\sqrt{2n\ell}.
        \label{eq:actual-to-extension}
    \end{align}
    Then $\Pp(\mathcal E_{\mathrm{reg}})\geq1-\delta/4$.
\end{lemma}
\proofstepinline{1}
To apply \cref{cor:bracket-ell}, we need to ensure that $\Sigma_t$ remains positive definite and that $x_\star$ remains in the focus ellipsoid around $\mu_t$.
We handle both by introducing stopping times that track the first round when either of these conditions is violated.
To this end, define the predictable (expected) counterpart of $\Sigma_t^{-1}$ as
\begin{align}
    \bar{\Sigma}_1^{-1}=\Sigma_1^{-1}=\sigma^{-2}I_d,
    \qquad
    \bar{\Sigma}_{t+1}^{-1}
    =\bar{\Sigma}_t^{-1}+\lambda\eta \bar{H}_t.
    \label{eq:predictable-covariance}
\end{align}
Then, define the stopping time $\tau_{\mathrm{cov}}$ when the empirical covariance $\Sigma_t$ deviates from its predictable counterpart $\bar{\Sigma}_t$ by more than a constant factor as
\begin{align}
    \tau_{\mathrm{cov}}
    =
    \min\left\{
    t \in [n]:
    \frac12 \bar{\Sigma}_{t+1}^{-1}
    \preceq
    \Sigma_{t+1}^{-1}
    \preceq
    \frac32 \bar{\Sigma}_{t+1}^{-1}
    \text{ fails}
    \right\}\wedge (n + 1).
    \label{eq:covariance-stopping-time}
\end{align}
Both $\Sigma_{t+1}^{-1} = \Sigma_t^{-1} + \lambda\eta H_t$ and $\bar{\Sigma}_{t+1}^{-1}$ are $\F_t$-measurable, so $\tau_{\mathrm{cov}}$ is an $\F_t$-adapted stopping time.

For $t \leq \tau_{\mathrm{cov}}$, the empirical covariance $\Sigma_t$ remains positive definite, which let us define the focus radius around $\mu_t$ as
\begin{align}
    R_t = \left\|\mu_t - x_\star\right\|_{\Sigma_t^{-1}}\,, \quad t \leq \tau_{\mathrm{cov}}.
\end{align}
Given a parameter $0 < F \leq \frac{1}{128^2\lambda^2 \ell}$ to be fixed later, let
\begin{align}
    \tau_{\mathrm{focus}}
    =
    \min\left\{
    t \in [n]:
    R_t^2 > F
    \quad\text{or}\quad
    t \geq \tau_{\mathrm{cov}}
    \right\}
    \wedge (n + 1).
\end{align}
Also define the two additional stopping times
\begin{align}
    \tau_{\mathrm{shell}}
    =\min\{t\in[n]:v(\mu_{t})>\ell\}\wedge (n+1)\,,
    \label{eq:energy-shell-stopping-time}
\end{align}
and for some parameter $S > 0$ to be fixed later, define the stopping time
\begin{align}
    \tau_{\mathrm{moment}}
    =\min\left\{
    t \in [n]: \sum_{i=1}^{t} \E_i[Y_i^2] > S
    \right\}\wedge (n+1)\,.
\end{align}
Finally, let $\tau = \tau_{\mathrm{cov}} \wedge \tau_{\mathrm{focus}} \wedge \tau_{\mathrm{shell}} \wedge \tau_{\mathrm{moment}}.$

Moreover, in various steps in the proof we use high-probability events to control the behavior of the random quantities.
It is convenient to define these events upfront and refer to them throughout the proof.
First, we define the high-probability event controlling the norm of the Gaussian samples
\begin{align}
    \mathcal E_Z
    =\left\{\norm{Z_t}^2\leq4d\ell\,\,
    \text{ for all }t\in [1, n \wedge \tau]\right\}\,.
    % \text{ for all }1\leq t\leq n\right\}.
    \label{eq:gaussian-norm-event}
\end{align}
For the noise, the sub-Gaussian assumption in \eqref{eq:noise} gives the high-probability event
\begin{align}
    \mathcal E_{\mathrm{noise}}
    =\left\{|\xi_t|\leq\sqrt\ell\,\,
    \text{ for all }t\in [1, n \wedge \tau]\right\}\,.
    % \text{ for all }1\leq t\leq n\right\}.
    \label{eq:noise-cutoff-event}
\end{align}
The displacement $X_t - \mu_t$ is Gaussian with bounded covariance $\Sigma_t\preceq 2\bar{\Sigma}_t \preceq 2 \bar{\Sigma}_1 =2\sigma^2I_d$, and $\pi$ is 1-Lipschitz, so we can control $\pi(X_t - \mu_t)$ with high probability by defining the event
\begin{align}
    \mathcal E_{\mathrm{gauge}}
    =\left\{\pi(X_t - \mu_t)\leq3+2\sqrt\ell\,\,
    \text{ for all }t\in [1, n \wedge \tau]\right\}\,.
    \label{eq:gauge-cutoff-event}
\end{align}
Let $\mathcal E_{\mathrm{global}} = \mathcal E_Z \cap \mathcal E_{\mathrm{noise}} \cap \mathcal E_{\mathrm{gauge}}$.
\begin{lemma}
    The event $\mathcal E_{\mathrm{global}}$ holds with high probability.
\end{lemma}

\proofstepinline{2}
The comparison in
\eqref{eq:time-indexed-regret-switch} requires $x_\star$ to remain
in an ellipsoid around $\mu_t$.
For $t < \tau$ both requirements of \cref{cor:bracket-ell} are satisfied:
we have $\frac12R_t^2 \leq \frac{1}{128^2\lambda^2\ell}$, and
$0 \prec \Sigma_t \preceq 2\bar{\Sigma}_t \preceq 2\sigma^2 I_d$.
Therefore, by \cref{eq:time-indexed-regret-switch} we obtain
\begin{align}
    \operatorname{eReg}_{1:\tau -1}(x_{\star})
    \leq
    % \overline{\operatorname{qReg}}_{1:\tau_1}(x_{\circ})
    \sum_{t=1}^{\tau -1}
    \left(\bar{q}_t(\mu_t)-\bar{q}_t(x_{\star})\right)
    +
    2\sum_{t=1}^{\tau -1}
    \tr(\bar{H}_t\Sigma_t)
    +\frac{(\tau -1) n^{-7}}{\eps}\,.
    \label{eq:regret-switch-focus}
\end{align}

\proofstepinline{3.1}
Bounding the quadratic regret $
    \overline{\operatorname{qReg}}_{1:\tau -1}(x_{\star})
    \triangleq
    \sum_{t=1}^{\tau -1}
    \left(\bar{q}_t(\mu_t)-\bar{q}_t(x_{\star})\right)
$ in \eqref{eq:regret-switch-focus} is quite delicate.
The main difficulty is that the quadratic $\bar{q}_t$ is not observed directly, but only through the sampled stochastic quadratic $q_t$.
\begin{lemma}
    \label{lem:ons-identity}
    For $t < \tau$, we have the following identity:
    \begin{align}
        \frac12
        \left(\normm{\mu_{t+1}-x_\star}{\Sigma_{t+1}^{-1}}^2
        -\normm{\mu_t-x_\star}{\Sigma_t^{-1}}^2
        \right)
         & =
        \eta\bigl(q_t(x_\star)-q_t(\mu_t)\bigr)
        +\frac{\eta^2}{2}\normm{\bar{g}_t}{\Sigma_{t+1}}^2.
    \end{align}
\end{lemma}
\begin{proof}
    Since $t < \tau_{\mathrm{cov}}$ we have $\Sigma_{t+1}, \Sigma_t \succ 0$.
    Therefore, we obtain
    \begin{align}
        \frac12\normm{\mu_{t+1}-x_\star}{\Sigma_{t+1}^{-1}}^2
         & =\frac12\normm{\mu_t-x_\star-\eta\Sigma_{t+1}g_t}
        {\Sigma_{t+1}^{-1}}^2                                 \\
         & =\frac12\normm{\mu_t-x_\star}{\Sigma_{t+1}^{-1}}^2
        +\eta\ip{g_t}{x_\star-\mu_t}
        +\frac{\eta^2}{2}\normm{g_t}{\Sigma_{t+1}}^2          \\
         & =\frac12\normm{\mu_t-x_\star}{\Sigma_t^{-1}}^2
        +\frac{\lambda\eta}{2}(\mu_t-x_\star)^\top H_t (\mu_t-x_\star)
        +\eta\ip{g_t}{x_\star-\mu_t}
        +\frac{\eta^2}{2}\normm{g_t}{\Sigma_{t+1}}^2          \\
         & =\frac12\normm{\mu_t-x_\star}{\Sigma_t^{-1}}^2
        +\eta\bigl(q_t(x_\star)-q_t(\mu_t)\bigr)
        +\frac{\eta^2}{2}\normm{g_t}{\Sigma_{t+1}}^2.
        \qedhere
    \end{align}
\end{proof}
The ONS identity from \cref{lem:ons-identity} gives us
the following telescoping identity for the stochastic quadratic regret
\begin{align}
    \operatorname{qReg}_{1:\tau -1}(x_{\star})
    \triangleq
    \sum_{t=1}^{\tau-1}
    q_t(\mu_t)-q_t(x_\star)
    =
    \frac{R_1^2 - R_\tau^2}{2\eta }
    +\frac{\eta}{2}\sum_{t=1}^{\tau-1}\normm{g_t}{\Sigma_{t+1}}^2.
    \label{eq:ons-telescope}
\end{align}
To control \eqref{eq:regret-switch-focus} through this identity,
we need to bound the difference between the sampled quadratics $q_t$
and the population quadratics $\bar{q}_t$.
Writing $D_t=\bar{q}_t(x_{\star})-q_t(x_{\star})$, which satisfies $\E_t[D_t]=0$,
the resulting error is
\begin{align}
    \left|
    \overline{\operatorname{qReg}}_{1:\tau-1}(x_{\star})-\operatorname{qReg}_{1:\tau-1}(x_{\star})
    \right|
    =
    \left|
    \sum_{t=1}^{\tau-1}
    D_t
    \right|\,,
\end{align}
where we used $q_t(\mu_t)=\bar{q}_t(\mu_t) = 0$.
Putting this back into \eqref{eq:ons-telescope} gives
\begin{align}
    \sum_{t=1}^{\tau-1}
    \bar{q}_t(\mu_t)-\bar{q}_t(x_\star)
     & \leq
    \frac{R_1^2}{2\eta }
    +\frac{\eta}{2}\sum_{t=1}^{\tau-1}\normm{g_t}{\Sigma_{t+1}}^2
    +\left|
    \sum_{t=1}^{\tau-1}
    D_t
    \right|\,,
\end{align}

\proofstepinline{3.2}
First, we deal with the sum $\frac{\eta}{2}\sum_{t=1}^{\tau - 1}\normm{g_t}{\Sigma_{t+1}}^2$.
We will work with the slightly time-perturbed quantity $\normm{g_t}{\Sigma_{t}}^2$ in place of $\normm{g_t}{\Sigma_{t+1}}^2$ for now.
By the definition of $g_t$ in \eqref{eq:sampled-hermite-coefficients},
on the event $\mathcal E_{\mathrm{global}}$, we have
\begin{align}
    \normm{g_t}{\Sigma_t}^2
    =Y_t^2\norm{Z_t}^2
    \leq
    4d\ell Y_t^2
    ,
    \label{eq:energy-feedback-identity}
\end{align}
where we used the bound $\norm{Z_t}^2\leq 4d\ell$ on the event $\mathcal E_{\mathrm{global}}$.
Next, to control the sum $\sum_{t=1}^{\tau-1}Y_t^2$ we use the following version of Freedman's inequality.
\begin{restatable}{lemma}{freedmannonnegative}
    \label{lem:freedman-nonnegative}
    Let $(U_i)_{i=1}^T$ be an adapted sequence.
    Then, for every $b,S,z>0$,
    \begin{align}
        \Pp\left(
        \sum_{i=1}^T U_i^2
        \geq 2S+4bz,
        \quad
        \sum_{i=1}^T\E_i[U_i^2]\leq S,
        \quad
        \max_{1\leq i\leq T}U_i^2\leq b
        \right)
        \leq2e^{-z}.
        \label{eq:freedman-nonnegative}
    \end{align}
\end{restatable}
To apply this lemma, with $U_t=Y_t$, we first bound the absolute size of the feedback $|Y_t|^2$.
since $f(A_t)\in[0,1]$, we have
\begin{align}
    |Y_t|=\left|\pi_\eps(X_t)(f(A_t)+\xi_t)+\frac{2v(X_t)}{\eps}\right|
    \leq
    \pi_\eps(X_t)(1+|\xi_t|)+\frac{2v(X_t)}{\eps}.
    \label{eq:energy-feedback-decomposition}
\end{align}
Thus we need to control both the observation noise $|\xi_t|$ and the capped Minkowski functional $v(X_t)$.
To bound $v(X_t)$ for $t < \tau$, write $W_t=X_t-\mu_t$, then the subadditivity of $\pi(\cdot)$ gives
\begin{align}
    v(X_t)\leq v(\mu_t)+\frac{\pi(W_t)}{1-\eps}
    \leq v(\mu_t)+2\pi(W_t)
    \leq v(\mu_t)+2(3 + 2\sqrt\ell)
    \,,
    \label{eq:energy-sample-gauge}
\end{align}
where we used the bound $\pi(W_t)\leq 3 + 2\sqrt\ell$ on the event $\mathcal E_{\mathrm{global}}$.
We also need to bound $v(\mu_t)$ in \eqref{eq:energy-sample-gauge}.
Projecting the mean onto $\K_\eps$ would make this term zero; our unprojected mean can leave $\K_\eps$.
Instead, we use the stopping time $\tau_{\mathrm{shell}}$ defined in \eqref{eq:energy-shell-stopping-time}.
On $\mathcal E_{\mathrm{global}}$, for every $t < \tau$, by \eqref{eq:energy-feedback-decomposition} and \eqref{eq:energy-sample-gauge} we have
\begin{align}
    |Y_t|
    \leq \pi_\eps(X_t)(1+|\xi_t|)+2\frac{\left(v(\mu_t) + 2 \pi(W_t)\right)}{\eps}
    \leq \frac{4\ell^{3/2}}{\eps}
    \label{eq:energy-sample-gauge-bound}
\end{align}
where we used
the numerical inequality $\ell\geq640$.
All together,
for $t < \tau$, on the event $\mathcal E_{\mathrm{global}}$, we have
\begin{align}
    \normm{g_t}{\Sigma_t}^2
    =Y_t^2\norm{Z_t}^2
    \leq\frac{64d\ell^4}{\eps^2}\,.
    \label{eq:energy-sample-gauge-bound-2}
\end{align}
To apply \cref{lem:freedman-nonnegative}, we also need to bound the conditional second moment $\sum_{t=1}^{\tau-1}\E_t[Y_t^2]$.
To this end, introduce the stopping time
\begin{align}
    \tau_{\mathrm{moment}}(S)
    =\min\left\{
    t \in [n]: \sum_{i=1}^{t} \E_i[Y_i^2] > S
    \right\}\wedge (n+1)\,,
\end{align}
for some fixed threshold $S>0$ to be specified later.
Now, let
\begin{align}
    \mathcal E_Y
    \triangleq
    \left\{
    \sum_{t=1}^{\tau-1}Y_t^2
    \leq
    2S + \frac{4d\ell^5}{\eps^2}
    \right\}\,.
    \label{eq:interval-feedback-bound}
\end{align}
\begin{lemma}
    The event $\mathcal E_\mathrm{global} \cap \mathcal E_Y^c$ holds with probability at most $\cdots$.
\end{lemma}
\begin{proof}
    Apply \cref{lem:freedman-nonnegative}, with  $z=\ell/64$ and $b=16\ell^3/\eps^2$...
\end{proof}
Therefore, on the event $\mathcal E_\mathrm{global} \cap \mathcal E_Y$, we get
\begin{align}
    \frac{\eta}{2}\sum_{t=1}^{\tau-1}\normm{g_t}{\Sigma_t}^2
    =
    \frac{\eta}{2}\sum_{t=1}^{\tau-1}Y_t^2\|Z_t\|^2
    \leq
    2\eta d\ell
    \sum_{t=1}^{\tau-1}Y_t^2
    \leq
    4\eta d\ell
    \left(
    S + \frac{2d\ell^5}{\eps^2}
    \right)\,.
    \label{eq:interval-energy-bound}
\end{align}
Finally, we compare $\normm{g_t}{\Sigma_t}^2$ and $\normm{g_t}{\Sigma_{t+1}}^2$.
On $\mathcal E_{\mathrm{global}}$, for every $t<\tau$, we have
\begin{align}
    \lambda\eta\left\|
    \Sigma_t^{1/2}H_t\Sigma_t^{1/2}
    \right\|_{\op}
     & =\lambda\eta|Y_t|\|Z_tZ_t^\top-I_d\|_{\op}
    \leq\frac{16\lambda\eta d\ell^{5/2}}{\eps}
    \leq\frac12,
\end{align}
where the last inequality follows from our parameter choices.
The update $\Sigma_{t+1}^{-1}
    =\Sigma_t^{-1}+\lambda\eta H_t$ therefore implies
$\frac23\Sigma_t
    \preceq\Sigma_{t+1}
    \preceq2\Sigma_t$,
and hence
$
    \frac23\normm{g_t}{\Sigma_t}^2
    \leq\normm{g_t}{\Sigma_{t+1}}^2
    \leq2\normm{g_t}{\Sigma_t}^2
$.
Thus, \eqref{eq:interval-energy-bound} gives
\begin{align}
    \frac{\eta}{2}\sum_{t=1}^{\tau-1}\normm{g_t}{\Sigma_{t+1}}^2
    \leq
    8\eta d\ell
    \left(
    S + \frac{2d\ell^5}{\eps^2}
    \right)\,.
\end{align}

\proofstepinline{3.3}
Going back to the quadratic regret, we now need to control the martingale difference sum $\sum_{t=1}^{\tau}D_t$, which is a martingale difference sequence as $\E_t[D_t]=0$.
We use the following self-normalized martingale inequality \citet[Theorem~3.26]{bercu2015concentration} to control this sum.
\begin{restatable}{lemma}{selfnormalisedmartingale}
    \label{lem:self-normalised-martingale}
    Let $(U_i)_{i=1}^T$ be a martingale-difference sequence. Then, for every $z,v>0$,
    \begin{align}
        \Pp\left(
        \left|\sum_{i=1}^T U_i\right|\geq z\quad\text{and}\quad
        \sum_{i=1}^T \left(U_i^2+\E_i[U_i^2]\right)\leq v
        \right)
        \leq
        2\exp\left(-\frac{z^2}{2v}\right).
        \label{eq:self-normalised-martingale}
    \end{align}
\end{restatable}
To apply it, we bound the realized and conditional second moments together.
Since $\bar{q}_t(x_{\star})=\E_t[q_t(x_\star)]$, a straightforward calculation gives
\begin{align}
    \sum_{t=1}^{\tau-1}\left(D_t^2+\E_t[D_t^2]\right)
     & \leq
    2\sum_{t=1}^{\tau-1}
    q_t(x_{\star})^2
    +
    2\sum_{t=1}^{\tau-1}
    \E_t[q_t(x_{\star})^2]\,,
    \label{eq:martingale-variance}
\end{align}
where we used $-2\bar{q}_t(x_{\star})q_t(x_{\star}) \leq \bar{q}_t(x_{\star})^2+q_t(x_{\star})^2$ and Jensen's inequality
$\bar{q}_t(x_{\star})^2\leq\E_t[q_t(x_{\star})^2]$.
Now, let $G_t = \langle \Sigma_t^{-1/2}(x_\star-\mu_t)/R_t, Z_t\rangle$
when $R_t>0$, and $G_t = 0$ when $R_t=0$.
% Then we can bound the second moment of $q_t(x_{\circ})$ as
% \begin{align}
%     \E_t\left[q_t(x_{\circ})^2\right]
%      & \leq R_t(x_{\circ})^2\E_t\left[
%         Y_t^2\left(
%         |G_t(x_{\circ})|+\frac12|G_t(x_{\circ})^2-1|
%         \right)^2\right]
%     \\
%      & \quad
%     \leq 4R_t(x_{\circ})^2\E_t\left[
%         Y_t^2\left(G_t(x_{\circ})^2+1\right)^2\right]
%     \leq 8F\E_t\left[
%         Y_t^2\left(G_t(x_{\circ})^2+1\right)^2\right].
% \end{align}
Then just from the definition of $q_t(x_\star)$, we have
\begin{align}
    q_t(x_{\star})^2
     & =
    Y_t^2\left[
        \left\langle \Sigma_t^{-1/2}(x_\star - \mu_t), Z_t \right\rangle
        + \frac{\lambda}{2}\left(
        \left\langle \Sigma_t^{-1/2}(x_\star - \mu_t), Z_t \right\rangle^2
        - \normm{x_\star - \mu_t}{\Sigma_t^{-1}}^2
        \right)
    \right]^2 \\
     & =
    Y_t^2\left[
        R_t G_t
        + \frac{\lambda}{2} R_t^2 \left(G_t^2 - 1\right)
    \right]^2 \\
     & \leq
    2R_t^2Y_t^2\left(G_t^2+1\right)^2
    \leq
    4F Y_t^2\left(G_t^2+1\right)^2
    \,,
    \label{eq:quadratic-estimate-decomposition}
\end{align}
where we used $(a+b)^2\leq2a^2+2b^2$ and the focus bound $R_t^2\leq\frac{2}{128^2\lambda^2 \ell}$, which also gives $\lambda R_t<1$.
Now, observe that the direction defining $G_t$ is then an
$\F_{t-1}$-measurable unit vector, so as long as $R_t>0$, the conditional distribution of
$G_t$ given $\F_{t-1}$ is $\N(0,1)$.
% Therefore, $|G_t| \leq 2\sqrt{\ell}$ with high probability by standard tail bounds.
This justifies the definition of the event
\begin{align}
    \mathcal E_G
    =\left\{
    |G_t|\leq2\sqrt\ell
    \text{ for all }1\leq t< \tau
    \right\}
    \,,
    \label{eq:gaussian-score-event}
\end{align}
on which we have $(G_t^2+1)^2\leq25\ell^2$.
Now, on the event $\mathcal E_G$, we can bound the realized part of the RHS of \eqref{eq:martingale-variance} as
\begin{align}
    2\sum_{t=1}^{\tau-1}
    q_t(x_\star)^2
    \leq
    2\sum_{t=1}^{\tau-1}
    4F Y_t^2\left(G_t^2+1\right)^2
    \leq
    200 F \ell^2 \sum_{t=1}^{\tau-1}
    Y_t^2
    \,,
    \label{eq:quadratic-bound}
\end{align}
and similarly for the conditional part of the RHS of \eqref{eq:martingale-variance}:
\begin{align}
    2\sum_{t=1}^{\tau-1}
    \E_t[q_t(x_\star)^2]
     & \leq
    200 F \ell^2 \sum_{t=1}^{\tau-1}
    \E_t[Y_t^2]
    +
    8F \sum_{t=1}^{\tau-1}\E_t\left[Y_t^2(G_t^2+1)^2\1_{\{|G_t|>2\sqrt\ell\}}\right]
    \\
     & \leq
    200 F \ell^2 \sum_{t=1}^{\tau-1}
    \E_t[Y_t^2]
    +
    8Fe^{-\ell/16}
    \,,
\end{align}
where the last inequality uses \cref{lem:weighted-gaussian-tail} to bound the summed weighted Gaussian tails.
Now, we can reuse our bounds from the previous step.
Then, on \eqref{eq:interval-feedback-bound} we also have $\sum_{t=1}^{\tau-1}Y_t^2\leq2S+4d\ell^5/\eps^2$.
Thus, absorbing the tail term into the coefficient gives
\begin{align}
    \sum_{t=1}^{\tau-1}\left(D_t^2+\E_t[D_t^2]\right)
     & \leq300F\ell^2\left(3S+\frac{4d\ell^5}{\eps^2}\right).
\end{align}
Define the event
\begin{align}
    \mathcal E_{q}
    =\left\{
    \left|\sum_{t=1}^{\tau-1}D_t\right|
    \leq\sqrt{F\ell^3\left(S+\frac{d\ell^5}{\eps^2}\right)}
    \right\},
    \label{eq:self-normalized-quadratic-event}
\end{align}
which holds by high probability due to \cref{lem:self-normalised-martingale}.
\begin{lemma}
    The event $\mathcal E_{\mathrm{global}} \cap \mathcal E_G \cap \mathcal E_{q}^c$ holds with probability at most ...
\end{lemma}
\proofstepinline{4}
Combining Steps~2--3.3, on the event
$\mathcal E_{\mathrm{global}}\cap\mathcal E_Y\cap\mathcal E_G\cap\mathcal E_q$,
the extension regret before the stopping time satisfies
\begin{align}
    \operatorname{eReg}_{1:\tau-1}(x_\star)
     & \leq
    \frac{R_1^2}{2\eta}
    +8\eta d\ell\left(S+\frac{2d\ell^5}{\eps^2}\right)
    +\sqrt{F\ell^3\left(S+\frac{d\ell^5}{\eps^2}\right)}
    +2\sum_{t=1}^{\tau-1}\tr(\bar H_t\Sigma_t)+1.
    \label{eq:stopped-regret-before-trace}
\end{align}
Here we used $(\tau-1)n^{-7}/\eps\leq n^{-6}/\eps\leq1$.
The remaining quantities to control are the thresholds $F,S$ and the trace sum.
We can bound the trace sum now, without choosing $F$ or $S$.

\begin{lemma}[Stopped trace bound]
    \label{lem:stopped-trace}
    Under the parameter choices of \cref{thm:main}, with $\eps=\eps_0\leq1/2$,
    \begin{align}
        \eta\sum_{t=1}^{\tau-1}\tr(\bar H_t\Sigma_t)
        \leq
        2\eta\sum_{t=1}^{\tau-1}\tr(\bar H_t\bar\Sigma_t)
        \leq\frac{8d\log(1+n)}{\lambda}.
        \label{eq:stopped-trace-budget}
    \end{align}
\end{lemma}
This bound holds pathwise; its proof is in \cref{app:trace}.
Substituting \eqref{eq:stopped-trace-budget} into
\eqref{eq:stopped-regret-before-trace}, and using the initialization
$R_1^2=\sigma^{-2}\norm{x_\star}^2\leq8d^3M^2$, gives
\begin{align}
    \operatorname{eReg}_{1:\tau-1}(x_\star)
     & \leq
    \frac{4d^3M^2}{\eta}
    +8\eta d\ell\left(S+\frac{2d\ell^5}{\eps^2}\right)
    +\sqrt{F\ell^3\left(S+\frac{d\ell^5}{\eps^2}\right)}
    +\frac{16d\log(1+n)}{\lambda\eta}+1.
    \label{eq:stopped-regret-after-trace}
\end{align}
We now fix the two remaining thresholds as
\begin{align}
    S=16n,
    \qquad
    F=\frac{1}{1000\lambda^2\ell^3}
    =1000d^3M^2\ell^5.
    \label{eq:global-threshold-choices}
\end{align}
Substituting these choices and the parameters of \cref{thm:main} into
\eqref{eq:stopped-regret-after-trace}, using
$\log(1+n)\leq\ell/128$ and $\ell\geq640$, gives
\begin{align}
    \operatorname{eReg}_{1:\tau-1}(x_\star)
    \leq\frac{F}{4\eta}
    =\frac{125}{4}d^2M\ell^5\sqrt n.
    \label{eq:global-stopped-regret}
\end{align}
The remaining steps of the proof are devoted to showing that $\tau=n+1$
on a high-probability event contained in
$\mathcal E_{\mathrm{global}}\cap\mathcal E_Y\cap\mathcal E_G\cap\mathcal E_q$.
On that event, \eqref{eq:global-stopped-regret} bounds the full-horizon
extension regret. Intersecting also with $\mathcal E_{\mathrm{reg}}$ from
Step~0 then gives the desired realized regret bound.



\newpage
\proofstepinline{4}
The previous steps~2--4 proves that, with high probability, for any fixed $\F_{s-1}$-measurable comparator $x_{\circ}$
and any interval $s\leq u< \tau_\mathrm{cov} \wedge \tau_{\mathrm{shell}} \wedge \tau_{\mathrm{focus}}(x_{\circ}; F) \wedge \tau_{\mathrm{tr}}(T) \wedge \tau_{\mathrm{moment}}(S)$ the bound
\begin{align}
    \overline{\operatorname{qReg}}_{s:u}(x_{\circ})
     & \leq
    \frac{R_s(x_{\circ})^2 - R_{u+1}(x_{\circ})^2}{2\eta }
    +\frac{\eta}{2}\sum_{t=s}^{u}\normm{g_t}{\Sigma_{t+1}}^2
    +\left|
    \sum_{t=s}^{u}
    D_t(x_{\circ})
    \right| \\
     & \leq
    \frac{F}{\eta }
    +8\eta d\ell\left(S+\frac{2d\ell^5}{\eps^2}\right)
    +\sqrt{F\ell^3\left(S+\frac{d\ell^5}{\eps^2}\right)}.
\end{align}
Further, together with \eqref{eq:regret-switch-focus}, this gives a bound on the extended regret $\operatorname{eReg}_{s:u}(x_{\circ})$ as
\begin{align}
    \operatorname{eReg}_{s:u}(x_{\circ})
     & \leq
    \frac{F}{\eta }
    +8\eta d\ell\left(S+\frac{2d\ell^5}{\eps^2}\right)
    +\sqrt{F\ell^3\left(S+\frac{d\ell^5}{\eps^2}\right)}
    +
    2T
    +1\,,
    \label{eq:regret-switch-focus-2}
\end{align}
This is a valid regret bound for $(s,u) = (1,\tau-1)$, where $\tau=\tau_{\mathrm{cov}} \wedge \tau_{\mathrm{focus}}(x_{\circ}; F) \wedge \tau_{\mathrm{tr}}(T) \wedge \tau_{\mathrm{moment}}(S) \wedge \tau_{\mathrm{shell}}$.

\proofstepinline{6}
Let $\bar{F}$ be a parameter to be chosen later, and let
\begin{align}
    \tau = \tau_{\mathrm{cov}} \wedge \tau_{\mathrm{tr}}(T) \wedge \tau_{\mathrm{moment}}(S) \wedge \tau_{\mathrm{shell}} \wedge u,\quad\text{and}\quad \bar{\tau} = \tau_{\mathrm{focus}}(x_{\circ}; \bar{F}) \wedge u\,.
\end{align}
We want to prove that for the correct choice of $\bar{F}$, as a function of other parameters like $S$ and $T$, we have, with high probability,
\begin{align}
    \bar{\tau} \geq \tau\,.
\end{align}
To this end, we start with the ONS telescoping sum:
\begin{align}
    \frac{1}{2}R_{\bar{\tau}\wedge \tau}(x_{\circ})^2
    =
    \frac{1}{2}R_s(x_{\circ})^2
    +\frac{\eta^2}{2}
    \sum_{t=s}^{\bar{\tau}\wedge \tau-1} \normm{g_t}{\Sigma_{t+1}}^2
    -\eta\, \operatorname{qReg}_{s:\bar{\tau}\wedge \tau-1}(x_{\circ})\,.
    \label{eq:ons-telescoping-sum}
\end{align}
We get back to the first two terms on RHS of \eqref{eq:ons-telescoping-sum}.
To upper bound the third term, we need a lower bound on $\operatorname{qReg}_{s:\bar{\tau}\wedge \tau-1}(x_{\circ})$.
To this end, write
\begin{align}
    \operatorname{qReg}_{s:\bar{\tau}\wedge \tau-1}(x_{\circ})
    \geq
    \overline{\operatorname{qReg}}_{s:\bar{\tau}\wedge \tau-1}(x_{\circ})
    - \left|
    \sum_{t=s}^{\bar{\tau}\wedge \tau-1} D_t(x_{\circ})
    \right|\,.
\end{align}
Then, from the regret switching lemma, we have
\begin{align}
    \overline{\operatorname{qReg}}_{s:\bar{\tau}\wedge \tau-1}(x_{\circ})
    \geq
    \operatorname{eReg}_{s:\bar{\tau}\wedge \tau-1}(x_{\circ})
    - 2\sum_{t=s}^{\bar{\tau}\wedge \tau-1} \tr\left(\bar{H}_t\Sigma_t\right) - \frac{n^{-7}}{\eps}(\tau \wedge \bar{\tau} - s)\,.
\end{align}
Now, let
\begin{align}
    E \triangleq e(x_{\circ}) - e(x_\star) \geq 0\,,
\end{align}
then
\begin{align}
    \operatorname{eReg}_{s:\bar{\tau}\wedge \tau-1}(x_{\circ})
    \geq
    -E(\bar{\tau}\wedge \tau - s)\,.
\end{align}
By chaining all these inequalities, we obtain an upper bound on \eqref{eq:ons-telescoping-sum}:
\begin{align}
    \frac{1}{2}R_{\bar{\tau}\wedge \tau}(x_{\circ})^2
     & \leq
    \frac{1}{2}R_s(x_{\circ})^2
    +\frac{\eta^2}{2}
    \sum_{t=s}^{\bar{\tau}\wedge \tau-1} \normm{g_t}{\Sigma_{t+1}}^2                                                                                                                  \\
     & +\eta \left( E(\bar{\tau}\wedge \tau - s) + 2\sum_{t=s}^{\bar{\tau}\wedge \tau-1} \tr\left(\bar{H}_t\Sigma_t\right) + \frac{n^{-7}}{\eps}(\tau \wedge \bar{\tau} - s) + \left|
    \sum_{t=s}^{\bar{\tau}\wedge \tau-1} D_t(x_{\circ})
    \right| \right)\,.
\end{align}
Now, let $L = u - s$, and apply the bounds we have derived on the high probability events  to get
\begin{align}
    \frac{1}{2}R_{\bar{\tau}\wedge\tau}(x_{\circ})^2
     & \leq
    \frac{1}{2}R_s(x_{\circ})^2
    +8\eta^2d\ell\left(S+\frac{2d\ell^5}{\eps^2}\right)
    +\eta(EL+2T+1)
    +\eta\sqrt{\bar{F}\ell^3\left(S+\frac{d\ell^5}{\eps^2}\right)}
    \triangleq \Lambda(x_{\circ}, S, L, T, \bar{F})
    \,.
    \label{eq:Lambda-definition}
\end{align}
Here we use the energy bound from Step~4.2,
\eqref{eq:self-normalized-quadratic-event} with $F=\bar{F}$,
the trace stopping condition, and $Ln^{-7}/\eps\leq1$.
So for $\bar{F}$, such that
\begin{align}
    \Lambda(x_{\circ}, S, L, T, \bar{F})\leq \bar{F} \leq \frac{1}{128^2\lambda^2 \ell}\,,
\end{align}
we have $R_{\bar{\tau}\wedge \tau}(x_{\circ})^2 \leq 2\bar{F}$, and consequently
\begin{align}
    \tau_{\mathrm{focus}}(x_{\circ}; \bar{F}) \wedge u \geq \tau \wedge u\,.
\end{align}


\proofstepinline{7}
Let $s \leq u \leq \tau_{\mathrm{cov}} \wedge \tau_{\mathrm{shell}}$,
and let $L = u - s + 1$.
By \cref{lem:uniform-trace,lem:single-round-feedback}, almost surely,
\begin{align}
    \sum_{t=s}^{u}\tr(\bar{H}_t\Sigma_t)
     & \leq T_L\triangleq\frac{5L}{\eps},
     & \sum_{t=s}^{u}\E_t[Y_t^2]
     & \leq S_L\triangleq\frac{8L\ell^2}{\eps^2},
    \label{eq:confinement-local-budget}
\end{align}
then for $x_{\circ} = \mu_s$ in \eqref{eq:Lambda-definition} we get
\begin{align}
    \Lambda(\mu_s, S_L, L, T_L, \bar{F})
\end{align}







\proofstepinline{8}
We will show that on a high-probability event, it holds that
\begin{align}
    \tau_{\mathrm{shell}} \geq \tau_{\mathrm{cov}}\wedge\tau_{\mathrm{focus}}(x_\star; F_\star)\,,
    \quad
    \text{for }\quad
    F_\star=\frac{1}{1000\lambda^2\ell^3}\,.
\end{align}
Let $x_{\circ} \in \R^d$ be a fixed $\F_{s-1}$-measurable comparator, and let
\begin{align}
    F_L=16\frac{\eta}{\eps}L\ell,
\end{align}
and $L$ is small enough that $F_L\leq F_\star$,
Let $s \leq u < \tau_{\mathrm{cov}}\wedge\tau_{\mathrm{focus}}(x_\star; F_\star)\wedge\tau_{\mathrm{focus}(x_{\circ}; F_L)}$ be arbitrary such that $L=u-s+1$.
By \cref{lem:uniform-trace,lem:single-round-feedback}, almost surely,
\begin{align}
    \sum_{t=s}^{u}\tr(\bar{H}_t\Sigma_t)
     & \leq T_L\triangleq\frac{5L}{\eps},
     & \sum_{t=s}^{u}\E_t[Y_t^2]
     & \leq S_L\triangleq\frac{8L\ell^2}{\eps^2},
    \label{eq:confinement-local-budgets}
\end{align}
Then, on the event $\mathcal E_Z\cap\mathcal E_{\mathrm{noise}}\cap\mathcal E_{\mathrm{gauge}} \cap \mathcal{E}_{G}^{(s:u)}(x_{\circ}) \cap \mathcal E_{q}^{(s:u)}(x_{\circ}) \cap \mathcal{E}_Y^{(s:u)}$, we have
\begin{align}
    \operatorname{eReg}_{s:u}(x_{\circ})
     & \leq
    \frac{F_L}{\eta }
    +8\eta d\ell\left(S_L+\frac{2d\ell^5}{\eps^2}\right)
    +\sqrt{F_L\ell^3\left(S_L+\frac{d\ell^5}{\eps^2}\right)}
    +
    2T_L
    +1\,.
\end{align}






\vspace{1em}
\hrule
\vspace{1em}
To control $v(\mu_t)$, first observe that nonnegativity of $f$ and the
definition of $e$ give
\begin{align}
    2v(x)\leq\eps e(x)
    \qquad\text{for all }x\in\R^d.
\end{align}
Moreover, convexity of $e$ and $\E_t[X_t]=\mu_t$ imply
$e(\mu_t)\leq\E_t[e(X_t)]$ by Jensen's inequality. Consequently,
\begin{align}
    \eps\sum_{t=s}^{u}\bigl(e(\mu_t)-e(x_{\circ})\bigr)
     & \leq\eps\sum_{t=s}^{u}\bigl(\E_t[e(X_t)]-e(x_{\circ})\bigr)
    =\eps\operatorname{eReg}_{s:u}(x_{\circ})\,.
\end{align}

Next, fix a deterministic interval for $s\leq u<\tau_{\mathrm{cov}}\wedge\tau_{\mathrm{shell}}$ of length $L=u-s+1$.
By \cref{lem:uniform-trace,lem:single-round-feedback}, almost surely,
\begin{align}
    \sum_{t=s}^{u}\tr(\bar{H}_t\Sigma_t)
     & \leq T_L:=\frac{5L}{\eps},
     & \sum_{t=s}^{u}\E_t[Y_t^2]
     & \leq S_L:=\frac{8L\ell^2}{\eps^2}.
    \label{eq:confinement-local-budgets}
\end{align}
Thus neither $\tau_{\mathrm{tr}}(T_L)$ nor
$\tau_{\mathrm{moment}}(S_L)$ restricts this interval further.
For comparison with its starting mean $\mu_s$, set
\begin{align}
    k=\left\lceil\frac{\eps F_\star}{\eta}\right\rceil,
    \qquad
    F_L=16\frac{\eta}{\eps}L\ell,
    \qquad 1\leq L\leq k.
\end{align}
Our parameter choices give $\eta/\eps=2/(25d\ell^8)\leq F_\star$, so
\begin{align}
    F_L\leq32F_\star\ell
    =\frac{32}{1000\lambda^2\ell^2}
    \leq\frac{1}{128^2\lambda^2\ell}.
\end{align}
Hence $F_L$ is an admissible focus threshold in Steps~2--5.

Work on $\mathcal E_Z\cap\mathcal E_{\mathrm{noise}}
    \cap\mathcal E_{\mathrm{gauge}}$ and the existing events
$\mathcal E_{q}^{(s:j)}(\mu_s)$ from Step~4.3, for the prefixes
used below with budgets $F=F_L$ and $S=S_L$.
The pointwise bound \eqref{eq:energy-sample-gauge-bound-2} and
$\Sigma_{t+1}\preceq2\Sigma_t$ give
\begin{align}
    \frac{\eps\eta}{2}\sum_{t=s}^{u}
    \normm{g_t}{\Sigma_{t+1}}^2
     & \leq64\frac{\eta}{\eps}Ld\ell^4
    =\frac{128L}{25\ell^4}\leq L.
    \label{eq:confinement-local-energy}
\end{align}
For every prefix ending at $j<\tau_{\mathrm{focus}}(\mu_s;F_L)$,
the martingale bound in Step~4.3 becomes
\begin{align}
    \eps\left|\sum_{t=s}^{j}D_t(\mu_s)\right|
     & \leq\eps\sqrt{F_L\ell^3
        \left(S_L+\frac{d\ell^5}{\eps^2}\right)}
    =\sqrt{\frac{256L^2}{25d\ell^2}+\frac{32L\ell}{25}}
    \leq L+2\sqrt{L\ell}.
    \label{eq:confinement-local-martingale}
\end{align}

We now show that this local focus condition holds throughout the
interval, including its ending mean $\mu_{u+1}$.
Since $v(\mu_s)\leq\ell$, the extension formula gives
$\eps e(\mu_s)\leq\eps+(2+\eps)\ell\leq3\ell$.
Before a first local focus violation, nonnegativity of $e$,
\eqref{eq:regret-switch-focus} on a single round, and the trace bound imply
\begin{align}
    \eps \bar{q}_t(\mu_s)
     & \leq\eps e(\mu_s)+10+n^{-7}\leq4\ell.
\end{align}
If $j\leq u+1$ were the first state with
$\tfrac12R_j(\mu_s)^2>F_L$, apply the ONS identity through round $j-1$.
All its summands precede the focus violation, and $R_s(\mu_s)=0$,
so the preceding bounds yield
\begin{align}
    \frac12R_j(\mu_s)^2
     & \leq\frac{\eta}{\eps}
    \left(4L\ell+2L+2\sqrt{L\ell}\right)
    <F_L,
\end{align}
a contradiction. Consequently,
\begin{align}
    \max_{s\leq t\leq u+1}\frac12R_t(\mu_s)^2\leq F_L.
    \label{eq:confinement-local-focus}
\end{align}
We can therefore apply Step~5 on the whole interval with comparator
$\mu_s$, retaining $R_s(\mu_s)=0$ and dropping the ending potential.
Together with Jensen's inequality and the bounds above, this gives
\begin{align}
    \eps\sum_{t=s}^{u}\bigl(e(\mu_t)-e(\mu_s)\bigr)
     & \leq\eps\operatorname{eReg}_{s:u}(\mu_s)
    \leq13L+2\sqrt{L\ell}.
    \label{eq:confinement-local-regret}
\end{align}

We next bound the average extension value over the last $k$ rounds.
For the global comparator $x_\star$, take $s=1$ in the definition of
$\tau_{\mathrm{focus}}(x_\star;F_\star)$.
Fix $t<\tau_{\mathrm{cov}}\wedge\tau_{\mathrm{shell}}
    \wedge\tau_{\mathrm{focus}}(x_\star;F_\star)$, and put
$m=\min\{k,t\}$ and $s=t-m+1$.
If $t<k$, then $s=1$, so \eqref{eq:confinement-local-regret} and
$\eps e(\mu_1)=\eps f(0)\leq1/2$ give
\begin{align}
    \frac{\eps}{m}\sum_{j=s}^{t}e(\mu_j)
     & \leq\frac12+13+2\sqrt{\frac{\ell}{m}}.
\end{align}
If $t\geq k$, use $x_\star$ as comparator on $[s,t]$, with budgets
$F=F_\star$, $T=T_k$, and $S=S_k$.
We also work on the existing events
$\mathcal E_{q}^{(s:t)}(x_\star)$ with these budgets for all such windows.
Since $F_\star/k\leq\eta/\eps$, their martingale bound gives
\begin{align}
    \frac{\eps}{k}\left|\sum_{j=s}^{t}D_j(x_\star)\right|
     & \leq\frac{\sqrt{F_\star\ell^3(8k\ell^2+d\ell^5)}}{k}
    \leq\sqrt{\frac{16}{25d\ell^3}+\frac{2}{25k}}
    \leq1.
\end{align}
Combining the ONS bound from Step~5 with
\cref{eq:confinement-local-budgets,eq:confinement-local-energy}, and using
$\eps e(x_\star)\leq1/2$, yields
\begin{align}
    \frac{\eps}{k}\sum_{j=s}^{t}e(\mu_j)
     & \leq\eps e(x_\star)+\frac{\eps F_\star}{\eta k}
    +12+\frac{\eps}{k}
    \leq14.
\end{align}
Thus, in both cases,
\begin{align}
    \frac{\eps}{m}\sum_{j=s}^{t}e(\mu_j)
     & \leq14+2\sqrt{\frac{\ell}{m}}.
    \label{eq:confinement-window-average}
\end{align}

To pass from this average to the ending mean, telescope the averages
over successive suffixes:
\begin{align}
    \eps e(\mu_t)
     & =\frac{\eps}{m}\sum_{j=s}^{t}e(\mu_j)
    +\eps\sum_{r=1}^{m-1}\frac{1}{r(r+1)}
    \sum_{j=t-r+1}^{t}\bigl(e(\mu_j)-e(\mu_{t-r})\bigr).
\end{align}
Each inner sum is the regret against the starting mean on the interval
$[t-r,t]$ of length $r+1\leq k$, with its zero first summand omitted.
Hence \cref{eq:confinement-local-regret,eq:confinement-window-average} imply
\begin{align}
    \eps e(\mu_t)
     & \leq14+2\sqrt{\frac{\ell}{m}}
    +13\sum_{r=1}^{m-1}\frac1r
    +2\sqrt\ell\sum_{r=1}^{m-1}\frac{1}{r\sqrt{r+1}}
    \leq27+13\log n+8\sqrt\ell
    \leq\frac\ell2.
    \label{eq:confinement-pointwise-extension}
\end{align}
Here we used $\sum_{r=1}^{m-1}1/r\leq1+\log n$,
$\sum_{r=1}^{\infty}1/(r\sqrt{r+1})
    \leq1+\int_1^\infty x^{-3/2}\,dx=3$,
and $\log n\leq\ell/128$ with $\ell\geq640$.
Since $2v\leq\eps e$, this proves $v(\mu_t)\leq\ell/4$ before all three stops.


Finally, suppose that
$\tau_{\mathrm{shell}}<\tau_{\mathrm{cov}}
    \wedge\tau_{\mathrm{focus}}(x_\star;F_\star)$.
Since $v(\mu_1)=0$, we have $\tau_{\mathrm{shell}}\geq2$.
Set $t=\tau_{\mathrm{shell}}-1$.
The covariance bounds $\Sigma_t\preceq2\sigma^2I_d\preceq I_d$
and $\Sigma_{t+1}\preceq2\Sigma_t$ give
$\|\Sigma_{t+1}\Sigma_t^{-1/2}\|_{\op}\leq2$.
The update and the pointwise feedback and Gaussian bounds therefore imply
\begin{align}
    \norm{\mu_{t+1}-\mu_t}
     & \leq2\eta|Y_t|\norm{Z_t}
    \leq16\frac{\eta}{\eps}\sqrt d\,\ell^2.
\end{align}
As $v$ is $1/(1-\eps)\leq2$-Lipschitz,
\begin{align}
    v(\mu_{\tau_{\mathrm{shell}}})
     & \leq\frac\ell4+32\frac{\eta}{\eps}\sqrt d\,\ell^2
    =\frac\ell4+\frac{64}{25\sqrt d\,\ell^6}
    <\ell,
\end{align}
contradicting the definition of $\tau_{\mathrm{shell}}$.
This proves $\tau_{\mathrm{shell}}\geq\tau_{\mathrm{cov}}
    \wedge\tau_{\mathrm{focus}}(x_\star;F_\star)$ on the intersection of the
existing events used above.

\newpage
\proofstepinline{6}
The remaining input to Steps~3--4 is a
deterministic budget for the sum of conditional feedback moments.
We enforce this budget globally by introducing more stopping times:
\begin{align}
    \tau_{\mathrm{moment}}
    =\min\left\{t\in[n]:\sum_{j=1}^t\E_j[Y_j^2]>16n\right\}\wedge n,
    \,\,\text{and}\,\,
    \tau_{\mathrm{stable}}
    =\tau_{\mathrm{cov}}\wedge\tau_{\mathrm{focus}}
    \wedge\tau_{\mathrm{shell}}\wedge\tau_{\mathrm{moment}}.
    \label{eq:global-stopping-time}
\end{align}
The sum can cross $16n$ on the stopping round. To control this overshoot,
\cref{lem:single-round-feedback} in the appendix gives the uniform bound
$\E_t[Y_t^2]\leq8\ell^2/\eps^2$ for every
$t\leq\tau_{\mathrm{cov}}\wedge\tau_{\mathrm{shell}}$.
Consequently, writing $\tau=\tau_{\mathrm{stable}}$,
\begin{align}
    S_{\mathrm{global}}\triangleq\sum_{t=1}^{\tau}\E_t[Y_t^2]
    \leq16n+\frac{8\ell^2}{\eps^2}
    \leq17n\,.
    \label{eq:global-moment-budget}
\end{align}
Thus we can take $S=17n$ for $(s,u)=(1,\tau)$ in \cref{lem:energy-feedback-event,lem:interval-quadratic-concentration}.
First, \cref{lem:energy-feedback-event}, with $S=17n$, gives
\begin{align}
    \sum_{t=1}^{\tau}Y_t^2
    \leq17n+\frac{\sqrt{34n}\,\ell^2}{\eps}
    +\frac{2\ell^4}{\eps^2}
    \leq20n
    \label{eq:global-feedback-bound}
\end{align}
on $\mathcal E_Z\cap\mathcal E_{\mathrm{noise}}\cap\mathcal E_{\mathrm{gauge}}$,
apart from an event of probability at most $2e^{-\ell/16}$.

For the quadratic error, the initialization gives
$\frac12R_1(x_\star)^2\leq4d^3M^2<B$.
The focus stopping rule therefore gives
$R_t(x_\star)\leq\sqrt{2B}$ and
$\lambda R_t(x_\star)\leq\lambda\sqrt{2B}<1$ for $t\leq\tau$,
so Step~4 applies with $F=B$.
Using the deterministic bound $V_{1:n}(B,17n)$ for the stopped
self-normalizer, its argument gives
\begin{align}
    \left|\sum_{t=1}^{\tau}D_t(x_\star)\right|
    \leq\sqrt{\frac{\ell V_{1:n}(B,17n)}8}
    \leq53\sqrt{Bn}\,\ell^{3/2}.
    \label{eq:global-quadratic-error}
\end{align}
Here we used $\ell^4/\eps^2\leq n$ and
$V_{1:n}(B,17n)\leq21624Bn\ell^2$.
This costs another $2e^{-\ell/16}$ on the feedback event from Step~3
and the score event $\{|G_t(x_\star)|\leq2\sqrt\ell\text{ for all }t\leq\tau\}$.
The Gaussian and noise tail bounds give
\begin{align}
    \Pp\bigl(\overline{\mathcal E}_Z\cup
    \overline{\mathcal E}_{\mathrm{noise}}\cup
    \overline{\mathcal E}_{\mathrm{gauge}}\bigr)
    \leq4ne^{-\ell},
    \qquad
    \Pp\left(\max_{t\leq\tau}|G_t(x_\star)|>2\sqrt\ell\right)
    \leq2ne^{-2\ell}.
\end{align}
Therefore all these bounds hold together with probability at least
$1-4ne^{-\ell}-2ne^{-2\ell}-4e^{-\ell/16}\geq1-\delta/4$.

It remains to replace $\Sigma_t$ by $\Sigma_{t+1}$
in the gradient energy, including at the final update.
On the same event, the feedback and Gaussian norm bounds imply
\begin{align}
    \lambda\eta\left\|
    \Sigma_t^{1/2}H_t\Sigma_t^{1/2}
    \right\|_{\op}
     & =\lambda\eta|Y_t|\|Z_tZ_t^\top-I_d\|_{\op}
    \leq\frac{16\lambda\eta}{\eps}(d+\ell)\ell^{3/2}
    <\frac12
    \qquad(t\leq\tau),
\end{align}
where we used the parameter choices in \cref{thm:main} and $M\sqrt d\geq1$.
The precision update thus satisfies
\begin{align}
    \frac12\Sigma_t^{-1}
    \preceq\Sigma_{t+1}^{-1}
    \preceq\frac32\Sigma_t^{-1},
    \qquad
    \Sigma_{t+1}\preceq2\Sigma_t
    \qquad(t\leq\tau).
    \label{eq:global-one-step-comparison}
\end{align}
In particular, the final update is positive definite, so the ONS
identity is valid through round $\tau$.
Combining \eqref{eq:energy-feedback-identity} and
\eqref{eq:global-feedback-bound} now gives
\begin{align}
    \frac{\eta}{2}\sum_{t=1}^{\tau}
    \normm{g_t}{\Sigma_{t+1}}^2
    \leq\eta\sum_{t=1}^{\tau}
    \normm{g_t}{\Sigma_t}^2
    \leq60\eta n(d+\ell)
    \leq80\eta dn\ell.
    \label{eq:global-gradient-energy}
\end{align}
Substituting this and \eqref{eq:global-quadratic-error} into the bound
from Step~2 yields the global quadratic-regret bound
\begin{align}
    \overline{\operatorname{qReg}}_{1:\tau}(x_\star)
    \leq
    \frac{R_1(x_\star)^2-R_{\tau+1}(x_\star)^2}{2\eta}
    +80\eta dn\ell+53\sqrt{Bn}\,\ell^{3/2}.
    \label{eq:global-quadratic-regret}
\end{align}
This bound holds up to $\tau_{\mathrm{stable}}$; removing the stopping
conditions and controlling the trace term are the remaining tasks.

BREAK
The trace sum in \eqref{eq:regret-switch-focus} is controlled by the following lemma, which is proved in \cref{app:trace}.
\begin{lemma}
    \label{lem:trace}
    We have
    \begin{align}
        \eta\sum_{t=1}^{\tau_{\mathrm{cov}}}
        \tr(\bar{H}_t\Sigma_t)
        \leq
        2\eta\sum_{t=1}^{\tau_{\mathrm{cov}}}
        \tr(\bar{H}_t\bar{\Sigma}_t)
        \leq\frac{8d\log(1 + n)}{\lambda}.
        \label{eq:trace-budget}
    \end{align}
\end{lemma}
For short intervals we also need a trace bound that does not grow with
the dimension. The next estimate is uniform over the Gaussian mean.
\begin{lemma}[Uniform trace bound]
    \label{lem:uniform-trace}
    Suppose that $\sigma=(M\sqrt{2d})^{-1}$, $0<\eps\leq1/2$, and
    $0\prec\Sigma\preceq2\sigma^2I_d$. For every $\mu\in\R^d$,
    \begin{align}
        0\leq\eps\tr\left(
        \langle D^2e,\gamma_{\mu,\Sigma}\rangle\Sigma
        \right)\leq5.
        \label{eq:uniform-trace}
    \end{align}
\end{lemma}
The proof is also in \cref{app:trace}.

Let $x_\star\in\argmin_{x\in\K_\eps}f(x)$.
This point exists because $\K_\eps$ is a compact subset of
$\operatorname{int}\K$, where the finite convex function $f$ is continuous.
Moreover, $x/\pi_\eps(x)\in\K_\eps$ for every $x\in\R^d$, so agreement on
$\K_\eps$ and \cref{lem:extension} show that $x_\star$ is a global
minimizer of $e$. Consequently,
\begin{align}
    \E_t[e(X_t)]-e(x_\star)\geq0.
    \label{eq:stationarity-sign}
\end{align}
For the analysis, define
\begin{align}
    \alpha & =\frac{\eta}{\eps},
           & k                                & =\left\lceil\frac B\alpha\right\rceil,
    \label{eq:analysis-budgets}
    \\
    R_t(x) & =\normm{\mu_t-x}{\Sigma_t^{-1}},
           & \operatorname{eReg}_m            & =\sum_{t=1}^m
    \left(\E_t[e(X_t)]-e(x_\star)\right).
    \label{eq:analysis-potentials}
\end{align}
The focus condition alone does not give a dimension-free bound on
$v(\mu_t)$. We therefore stop separately at a violation of this bound or
of a predictable feedback budget, and subsequently rule out all four
stopping conditions. These stopping times are used only in the proof;
the algorithm does not compute $x_\star$ or any conditional expectation.

In addition to $\tau_{\mathrm{cov}}$, define
\begin{align}
    \tau_{\mathrm{shell}}
     & =\min\{t\in[n]:v(\mu_{t+1})>\ell\}\wedge n,
    \\
    \tau_{\mathrm{moment}}
     & =\min\left\{t\in[n]:\sum_{j=1}^t\E_j[Y_j^2]>16n\right\}\wedge n,
    \\
    \tau_{\mathrm{stable}}
     & =\tau_{\mathrm{cov}}\wedge\tau_{\mathrm{focus}}
    \wedge\tau_{\mathrm{shell}}\wedge\tau_{\mathrm{moment}}.
    \label{eq:safe-stopping-time}
\end{align}
All four are $(\F_t)$-stopping times. The initial focus potential satisfies
\begin{align}
    \frac12R_1(x_\star)^2
    \leq4d^3M^2<B,
    \qquad v(\mu_1)=0,
    \label{eq:initial-focus}
\end{align}
where we used $\mu_1=0$, $\Sigma_1^{-1}=\sigma^{-2}I_d$, and
$\norm{x_\star}\leq2d$. Hence, writing $\tau=\tau_{\mathrm{stable}}$,
minimality gives
\begin{align}
    \frac12R_t(x_\star)^2 & \leq B,
                          & v(\mu_t)                                              & \leq \ell,
                          & \frac12\bar{\Sigma}_t^{-1}
                          & \preceq\Sigma_t^{-1}\preceq\frac32\bar{\Sigma}_t^{-1}
    \qquad(t\leq\tau).
    \label{eq:stability-stopped-invariants}
\end{align}
In particular, $\Sigma_t\preceq2\sigma^2I_d\preceq I_d$.
The roundwise bound in \cref{lem:feedback} allows a single overshoot of
the moment budget:
\begin{align}
    \sum_{t=1}^\tau\E_t[Y_t^2]
    \leq16n+\frac{8\ell^2}{\eps^2}
    \leq17n.
    \label{eq:stopped-moment-budget}
\end{align}
The last inequality follows from
$n\eps^2=10^4d^4M^2\ell^{16}\geq8\ell^2$.
This bound is available before any stopping condition is removed.

The next proposition contains the stochastic estimates used in both the
stability and regret arguments.
\begin{proposition}
    \label{prop:concentration}
    Suppose that $\eps=\eps_0\leq1/2$. With probability at least
    $1-\delta/4$, simultaneously,
    \begin{align}
        \begin{aligned}
            \sum_{t=1}^{\tau_{\mathrm{stable}}}Y_t^2       & \leq20n,
            \\
            \max_{t\leq\tau_{\mathrm{stable}}}|\eps Y_t|   & \leq4\ell^{3/2},
            \\
            \max_{t\leq\tau_{\mathrm{stable}}}\norm{Z_t}^2 & \leq3(d+\ell),
        \end{aligned}
        \label{eq:stopped-increment-bounds}
    \end{align}
    \begin{align}
        \left|
        \sum_{t=1}^{\tau_{\mathrm{stable}}}
        \bigl(q_t(x_\star)-\bar{q}_t(x_\star)\bigr)
        \right|
         & \leq8\sqrt{Bn}\,\ell,
        \label{eq:q-concentration}
        \\
        \sum_{t=1}^{\tau_{\mathrm{stable}}}
        \normm{g_t}{\Sigma_t}^2
         & \leq80dn\ell,
        \label{eq:gradient-energy}
    \end{align}
    and
    \begin{align}
        -100\ell^{3/2}\sqrt{dn}\,
        \bar{\Sigma}_{\tau_{\mathrm{stable}}+1}^{-1}
        \preceq
        \sum_{t=1}^{\tau_{\mathrm{stable}}}(H_t-\bar{H}_t)
        \preceq
        100\ell^{3/2}\sqrt{dn}\,
        \bar{\Sigma}_{\tau_{\mathrm{stable}}+1}^{-1}.
        \label{eq:H-concentration}
    \end{align}
\end{proposition}
The proof is in \cref{app:concentration}. The additional estimate needed
without projection is a bound on the extension value at each mean.
It follows from ONS on intervals of length at most $k$, using the
uniform trace bound in \cref{lem:uniform-trace}.
\begin{lemma}[Confinement]
    \label{lem:confinement}
    There is an event of probability at least $1-\delta/4$ on whose
    intersection with the event of \cref{prop:concentration} the following
    holds. If the covariance and focus conditions hold at round
    $\tau=\tau_{\mathrm{stable}}$, then
    \begin{align}
        \max_{t\leq\tau+1}v(\mu_t)\leq\frac\ell2.
        \label{eq:confinement}
    \end{align}
\end{lemma}
The proof is in \cref{app:confinement}. It uses the stopped moment and
shell bounds, not the conclusion of \cref{thm:main}.

\begin{proof}[Proof of \cref{thm:main}]
    If $\eps_0>1/2$, then
    \begin{align}
        \operatorname{Reg}_n\leq n
        <200d^2M\ell^8\sqrt n,
    \end{align}
    which proves \eqref{eq:main-regret}. Suppose henceforth that
    $\eps=\eps_0\leq1/2$. In this case,
    \begin{align}
        \alpha=\frac{2}{25d\ell^8},
        \qquad
        n\geq4\cdot10^4d^4M^2\ell^{16},
        \qquad
        32B\ell\leq n,
        \qquad k<n.
        \label{eq:active-horizon}
    \end{align}
    The last two inequalities follow by substitution, using
    $dM^2\geq1$ and $\ell\geq640$.

    \proofstep{1}{Stability}
    Write $\tau=\tau_{\mathrm{stable}}$ and work on the intersection of the
    events in \cref{prop:concentration,lem:confinement}. Its probability
    is at least $1-\delta/2$.

    \proofstepinline{1.1}
    We first show that the covariance condition holds at round $\tau$.
    The coefficient in \eqref{eq:H-concentration} satisfies
    \begin{align}
        100\lambda\eta\ell^{3/2}\sqrt{dn}
        =\frac{4}{5\ell^{5/2}}<\frac14.
        \label{eq:main-stability-H-small}
    \end{align}
    Since every earlier covariance comparison holds, the proposed
    precision at round $\tau$ is
    \begin{align}
        \Sigma_\tau^{-1}+\lambda\eta H_\tau
        =\bar{\Sigma}_{\tau+1}^{-1}
        +\lambda\eta\sum_{t=1}^\tau(H_t-\bar{H}_t).
    \end{align}
    Therefore \eqref{eq:H-concentration} gives
    \begin{align}
        \frac34\bar{\Sigma}_{\tau+1}^{-1}
        \preceq
        \Sigma_\tau^{-1}+\lambda\eta H_\tau
        \preceq\frac54\bar{\Sigma}_{\tau+1}^{-1}.
        \label{eq:main-strict-comparison}
    \end{align}
    The update is accepted and the covariance condition does not fail.
    The comparisons before $\tau$ and \eqref{eq:main-strict-comparison}
    at $\tau$ give, for every $t\leq\tau$,
    \begin{align}
        \Sigma_{t+1}^{-1}
        \succeq\frac12\bar{\Sigma}_{t+1}^{-1}
        \succeq\frac12\bar{\Sigma}_t^{-1}
        \succeq\frac13\Sigma_t^{-1}.
    \end{align}
    Inverting and applying \eqref{eq:gradient-energy} yields
    \begin{align}
        \sum_{t=1}^\tau\normm{g_t}{\Sigma_{t+1}}^2
        \leq3\sum_{t=1}^\tau\normm{g_t}{\Sigma_t}^2
        \leq240dn\ell.
        \label{eq:main-next-gradient-energy}
    \end{align}
    Similarly, $\Sigma_t\preceq2\bar{\Sigma}_t$ and
    \cref{lem:trace} give
    \begin{align}
        \eta\sum_{t=1}^\tau\tr(\bar{H}_t\Sigma_t)
        \leq2\eta\sum_{t=1}^\tau\tr(\bar{H}_t\bar{\Sigma}_t)
        \leq\frac{8d\log(1+n)}{\lambda}.
        \label{eq:main-stability-trace}
    \end{align}

    \proofstepinline{1.2}
    We next control the focus potential and the extension regret together.
    We apply \cref{cor:bracket-ell} on a slightly larger region, which
    will also be used in the confinement proof. The covariance comparison
    gives $0\prec\Sigma_t\preceq2\sigma^2I_d$, and
    \eqref{eq:extension-lipschitz} gives $\Lip(e)\leq8/\eps$.
    For every $x$ with $\frac12R_t(x)^2\leq32B\ell$, the focus condition holds because
    \begin{align}
        2\ell\lambda^2\normm{\mu_t-x}{\Sigma_t^{-1}}^2
         & =2\lambda^2R_t(x)^2\ell
        \leq4\lambda^2(32B\ell)\ell
        =\frac{16}{125\ell}
        \leq\frac1{5000}<\frac1{64^2}.
        \label{eq:main-stability-bracket-focus}
    \end{align}
    Thus \cref{cor:regret-switch,cor:bracket-ell} give
    \begin{align}
        \bar{q}_t(x)
         & \leq e(x)-\E_t[e(X_t)]
        +\frac{17}{16}\tr(\bar{H}_t\Sigma_t)+\frac{n^{-8}}{\eps}
        \qquad\bigl(\tfrac12R_t(x)^2\leq32B\ell\bigr).
        \label{eq:main-q-upper}
    \end{align}

    Summing \eqref{eq:ons-step} with $x=x_\star$, applying
    \eqref{eq:main-q-upper}, and then using the concentration and trace
    estimates gives
    \begin{align}
        \frac12R_{\tau+1}(x_\star)^2+\eta\operatorname{eReg}_\tau
         & \leq \frac12R_1(x_\star)^2
        +\eta\sum_{t=1}^\tau
        \bigl(q_t(x_\star)-\bar{q}_t(x_\star)\bigr)
        \nonumber                                                       \\
         & \quad+\frac{17\eta}{16}\sum_{t=1}^\tau\tr(\bar{H}_t\Sigma_t)
        +\frac{\eta^2}{2}\sum_{t=1}^\tau
        \normm{g_t}{\Sigma_{t+1}}^2
        +\alpha\tau n^{-8}
        \nonumber                                                       \\
         & \leq4d^3M^2
        +\frac{16d\log(1+n)}{\lambda}
        +8\eta\sqrt{Bn}\,\ell
        +120\eta^2dn\ell+1
        \leq\frac B4.
        \label{eq:main-focus-decomposition}
    \end{align}
    To check the final inequality, divide its preceding expression by
    $B$ and use $\log(1+n)\leq\ell/128$:
    \begin{align}
        \frac{4}{1000\ell^5}
        +\frac18
        +\frac{64}{\sqrt{1000d}\,\ell^{3/2}}
        +\frac{192}{25\ell^4}
        +\frac1B
        <\frac14.
        \label{eq:main-budget-check}
    \end{align}
    Here $M\sqrt d\geq1$ bounds the trace term and
    $\ell\geq640$ bounds the displayed numerical expression.
    Nonnegativity of $\operatorname{eReg}_\tau$ now implies
    $\frac12R_{\tau+1}(x_\star)^2\leq B/4$, so the focus condition does not fail.
    The same inequality also gives $\operatorname{eReg}_\tau\leq B/(4\eta)$.

    \proofstepinline{1.3}
    We now improve the stopped feedback budget. The self-bound in
    \cref{lem:feedback}, followed by \eqref{eq:main-focus-decomposition},
    gives
    \begin{align}
        \sum_{t=1}^\tau\E_t[Y_t^2]
         & \leq4n+\frac{32\ell}{\eps}\operatorname{eReg}_\tau
        +\frac{32ne^{-\ell}}{\eps^2}
        \\
         & \leq4n+\frac{8\ell B}{\eta\eps}
        +\frac{32ne^{-\ell}}{\eps^2}
        =4n+\frac{10n}{\ell^2}
        +\frac{32ne^{-\ell}}{\eps^2}
        \leq6n<16n.
        \label{eq:predictable-Y2}
    \end{align}
    The tail term is at most $n$, since $\eps^{-2}\leq n$ and
    $32ne^{-\ell}\leq1$. Thus the moment condition does not fail.
    Notice the order of the argument: \eqref{eq:stopped-moment-budget}
    was used to prove concentration, and the resulting extension-regret
    bound is used only now to improve that budget.

    \proofstepinline{1.4}
    The covariance and focus conditions have been verified at $\tau$.
    Therefore \cref{lem:confinement} applies and gives
    $v(\mu_{\tau+1})\leq\ell/2<\ell$. The shell condition does not fail either.
    If $\tau<n$, all four conditions would hold at the first round on
    which one of them fails, a contradiction. Hence
    \begin{align}
        \tau=\tau_{\mathrm{stable}}=n.
        \label{eq:main-stability}
    \end{align}
    The preceding argument also verifies the conditions at round $n$,
    so every proposed precision is positive definite, including the last.

    \proofstep{2}{Regret bound}
    Azuma's inequality and $f(A_t)\in[0,1]$ give, with probability at least
    $1-\delta/4$,
    \begin{align}
        \sum_{t=1}^n\left(f(A_t)-\E_t[f(A_t)]\right)
        \leq\sqrt{2n\log\frac4\delta}
        \leq\sqrt{2n\ell}.
        \label{eq:action-martingale}
    \end{align}
    Intersecting with the events used in \texttt{Step 1} costs at most
    $3\delta/4\leq\delta$ in total.

    \proofstepinline{2.1}
    With $\tau=n$, the joint bound \eqref{eq:main-focus-decomposition}
    immediately yields
    \begin{align}
        \sum_{t=1}^n\left(\E_t[e(X_t)]-e(x_\star)\right)
        \leq\frac B{4\eta}
        =\frac{125}{4}d^2M\ell^5\sqrt n.
        \label{eq:main-extension-regret}
    \end{align}

    \proofstepinline{2.2}
    We finally pass from the extension regret to the realized regret.
    For every $x\in\K$, the definition of $x_\star$, convexity, and
    $f\in[0,1]$ give
    \begin{align}
        f(x_\star)
        \leq f((1-\eps)x)
        \leq(1-\eps)f(x)+\eps f(0)
        \leq f(x)+\eps.
    \end{align}
    Taking the infimum over $x\in\K$ yields
    \begin{align}
        f(x_\star)-\inf_{x\in\K}f(x)\leq\eps.
        \label{eq:shrink-comparator}
    \end{align}
    Moreover, $A_t=X_t/\pi_\eps(X_t)\in\K_\eps$ regardless of $\mu_t$,
    so \cref{lem:extension} gives $f(A_t)=e(A_t)\leq e(X_t)$, while
    $f(x_\star)=e(x_\star)$. Therefore,
    \begin{align}
        \operatorname{Reg}_n
         & =\sum_{t=1}^n\left(f(A_t)-\E_t[f(A_t)]\right)
        +\sum_{t=1}^n\left(\E_t[f(A_t)]-f(x_\star)\right)
        +n\left(f(x_\star)-\inf_{x\in\K}f(x)\right)
        \nonumber                                        \\
         & \leq\sqrt{2n\ell}
        +\frac{125}{4}d^2M\ell^5\sqrt n
        +100d^2M\ell^8\sqrt n
        \nonumber                                        \\
         & \leq\left(\sqrt2+\frac{125}{4}+100\right)
        d^2M\ell^8\sqrt n
        <200d^2M\ell^8\sqrt n.
    \end{align}
    The penultimate inequality uses $M\geq d^{-1/2}$ and $n,d,\ell\geq1$.
\end{proof}

\newpage
\begin{thebibliography}{9}
    \raggedright

    \bibitem[Bercu et al.(2015)Bercu, Delyon, and Rio]{bercu2015concentration}
    Bernard Bercu, Bernard Delyon, and Emmanuel Rio.
    \newblock \emph{Concentration Inequalities for Sums and Martingales}.
    \newblock SpringerBriefs in Mathematics. Springer, Cham, 2015.
    \newblock \href{https://doi.org/10.1007/978-3-319-22099-4}
    {doi:10.1007/978-3-319-22099-4}.

    \bibitem[Evans and Gariepy(2015)]{evansgariepy2015measure}
    Lawrence C. Evans and Ronald F. Gariepy.
    \newblock \emph{Measure Theory and Fine Properties of Functions}.
    \newblock Revised edition, CRC Press, 2015.

    \bibitem[Fokkema et al.(2024)Fokkema, van der Hoeven, Lattimore, and Mayo]{fokkema2024online}
    Hidde Fokkema, Dirk van der Hoeven, Tor Lattimore, and Jack J. Mayo.
    \newblock Online Newton method for bandit convex optimisation: Extended abstract.
    \newblock In \emph{Proceedings of the 37th Conference on Learning Theory},
    volume 247 of \emph{Proceedings of Machine Learning Research}, pages
    1713--1714, 2024.
    \newblock Full version: \href{https://arxiv.org/abs/2406.06506}{arXiv:2406.06506}.

    \bibitem[Lattimore(2026)]{lattimore2026bandit}
    Tor Lattimore.
    \newblock \emph{Bandit Convex Optimisation}.
    \newblock Cambridge University Press, 2026.
    \newblock \href{https://doi.org/10.1017/9781009607551}{doi:10.1017/9781009607551}.

    \bibitem[Lattimore and Gy\"orgy(2023)]{lattimore2023second}
    Tor Lattimore and Andr\'as Gy\"orgy.
    \newblock A second-order method for stochastic bandit convex optimisation.
    \newblock In \emph{Proceedings of the 36th Conference on Learning Theory},
    volume 195 of \emph{Proceedings of Machine Learning Research}, pages
    2067--2094, 2023.

    \bibitem[Laurent and Massart(2000)]{laurent2000adaptive}
    B\'eatrice Laurent and Pascal Massart.
    \newblock Adaptive estimation of a quadratic functional by model selection.
    \newblock \emph{The Annals of Statistics}, 28(5):1302--1338, 2000.
    \newblock \href{https://doi.org/10.1214/aos/1015957395}
    {doi:10.1214/aos/1015957395}.

    \bibitem[Vershynin(2018)]{vershynin2018high}
    Roman Vershynin.
    \newblock \emph{High-Dimensional Probability: An Introduction with
        Applications in Data Science}.
    \newblock Cambridge University Press, 2018.


\end{thebibliography}

\newpage
\appendix
\section{Properties of the convex extension}
\label{app:extension}
\begin{lemma}[{\citet[Proposition~3.23]{lattimore2026bandit}}]
    \label{lem:extension}
    The function $\bar{e}:
        x
        \longmapsto
        \pi_\eps(x)
        f\!\left(\frac{x}{\pi_\eps(x)}\right)
        +
        \frac{1-\eps}{\eps}v(x)
    $ is convex, agrees with $f$ on $\K_\eps$, and
    \begin{align}
        \Lip(\bar e)\leq\frac{2}{\eps(1-\eps)}\,.
    \end{align}
    Moreover,
    $e\!\left(\frac{x}{\pi_\eps(x)}\right)\leq e(x)$
    for all $x\in\R^d$.
\end{lemma}
Now, observe that $e(x)=\bar e(x)+\frac{1+\eps}{\eps}v(x)$, so $e$ is also convex, agrees with $f$ on $\K_\eps$ since $v(x)=0$ for $x\in\K_\eps$.
Moreover, $v$ is $1/(1-\eps)$-Lipschitz, so $e$ is also Lipschitz with
\begin{align}
    \Lip(e)\leq\frac{2}{\eps(1-\eps)}+\frac{1+\eps}{\eps}\cdot\frac{1}{1-\eps}\leq\frac{8}{\eps}\,,
\end{align}
when $\eps\in(0,1/2]$.

\section{Distributional derivatives}
\label{app:distributional-derivatives}
A general convex function $h:\R^d\to\R$ may not be differentiable everywhere and
its Hessian may not exist in the classical sense.
Traditionally, this issue is handled through a mollification process: convolve $h$ with a smooth kernel to obtain a smooth approximation, and then take limits.
However, we will take a different approach that avoids mollification and works directly with the original function $h$.
We employ the theory of distributions, which allows us to define derivatives in a weak sense.
Practically, first we recall the definition of the Schwartz space $\mathcal S(\R^d)$, which consists of smooth functions that decay faster than any polynomial at infinity, along with all their derivatives.
We write $Dh$ and $D^2h$ for the first and second distributional derivatives of $h$, which
are defined as linear functionals on the Schwartz space $\mathcal S(\R^d)$ of smooth rapidly decaying functions.
\begin{align}
    \langle Dh,\varphi\rangle
     & =-\int_{\R^d}h(x)\nabla\varphi(x)\,dx,
     &
    \langle D^2h,\varphi\rangle
     & =\int_{\R^d}h(x)\nabla^2\varphi(x)\,dx,
    \qquad
    \varphi\in\mathcal S(\R^d),
\end{align}
where $\varphi$ is the Gaussian density $\gamma_{\mu,\Sigma}$ for our purposes,
which is a Schwartz function.
The crucial fact here is that if $h$ is convex, then $D^2h$ is a positive-semidefinite matrix-valued \textit{tempered} Radon measure and the Gaussian density $\gamma_{\mu,\Sigma}$ is a Schwartz function, so the pairing $\langle D^2h,\gamma_{\mu,\Sigma}\rangle$ is well-defined and yields a positive-semidefinite matrix.
\alireza{This subsection needs more explanations and citations.}

\section{Hessian Measures}
First we record the standard fact that the distributional Hessian of a convex function is a positive-semidefinite matrix-valued Radon measure.
\begin{lemma}[Hessian measure of a convex function]
    \label{lem:convex-hessian-measure}
    Let $r:\R^d\to\R$ be convex and Lipschitz.
    Then the distributional Hessian $D^2r$ is a positive-semidefinite
    matrix-valued Radon measure.
\end{lemma}
\begin{proof}
    This is standard; see \citet[Section~6.3]{evansgariepy2015measure}.
\end{proof}
Next lemma gives the Gaussian averaging identity for the Hessian measure, which is the key to the direct Hermite coefficients.
\begin{lemma}[Gaussian averaging of the Hessian measure]
    \label{lem:gaussian-hessian-average}
    Let $r:\R^d\to\R$ be convex and Lipschitz, let
    $X=\mu+\Sigma^{1/2}Z$, where $Z\sim\N(0,I_d)$ and $\Sigma\succ0$, and
    let $\gamma_{\mu,\Sigma}$ be the density of $X$.  Then
    \begin{align}
        \langle D^2r,\gamma_{\mu,\Sigma}\rangle
        =
        \int_{\R^d}
        \gamma_{\mu,\Sigma}(x)\,D^2r(dx)
        =
        \E\left[
            r(X)\Sigma^{-1/2}(ZZ^\top-I_d)\Sigma^{-1/2}
            \right].
        \label{eq:gaussian-hessian-average}
    \end{align}
\end{lemma}
\begin{proof}
    By \cref{lem:convex-hessian-measure}
    $D^2r$ is a positive-semidefinite matrix-valued Radon measure.
    Since $r$ is Lipschitz, $D^2r \in \mathcal S'(\R^d)$ is a tempered distribution,
    and the Gaussian density $\gamma_{\mu,\Sigma}\in\mathcal S(\R^d)$ is a Schwartz function,
    which means that $\langle D^2r,\gamma_{\mu,\Sigma}\rangle < \infty$ and
    \begin{align}
        \langle D^2r,\gamma_{\mu,\Sigma}\rangle
        =
        \langle r,D^2\gamma_{\mu,\Sigma}\rangle
        =
        \E\left[
            r(X)\Sigma^{-1/2}(ZZ^\top-I_d)\Sigma^{-1/2}
            \right],
    \end{align}
    where the first equality is integration by parts applied twice.
\end{proof}

\section{Regret switching}
\label{app:regret-switch}
\begin{lemma}
    \label{lem:regret-switch}
    Let $X\sim\N(\mu,\Sigma)$, then for every $x\in\R^d$,
    \begin{align}
        \E[e(X)]-e(x)
        =\tr(\bar{H}\Sigma)+\ip{\bar{g}}{\mu-x}-\E[D_e(x\,\|\,X)]\,,
        \label{eq:regret-switch}
    \end{align}
    where $\bar{g}=\langle De,\gamma_{\mu,\Sigma}\rangle$ and
    $\bar{H}=\langle D^2e,\gamma_{\mu,\Sigma}\rangle$.
\end{lemma}
\begin{proof}
    Splitting $X-x=(X-\mu)+(\mu-x)$ in a rearrangement of \eqref{eq:bregman} and taking expectations gives
    \begin{align}
        \E[e(X)]-e(x)
         & =
        \E\left[\ip{\nabla e(X)}{X-\mu}\right]
        +
        \ip{\E[\nabla e(X)]}{\mu-x}
        -
        \E[D_e(x\,\|\,X)]\,,
    \end{align}
    and the proof is completed by showing that $\bar{g}=\E[\nabla e(X)]$ and
    $\tr(\bar{H}\Sigma)=\E[\ip{\nabla e(X)}{X-\mu}]$.

    Since $e$ is Lipschitz, it is differentiable Lebesgue-almost everywhere,
    and its almost-everywhere partial derivatives are bounded by $\Lip(e)$
    and represent its weak derivatives.
    Applying this identity in every coordinate yields $\bar{g}=\E[\nabla e(X)]$,
    which proves the first identity.

    We next identify the trace term. Since
    $\bar{H}=\langle D^2e,\gamma_{\mu,\Sigma}\rangle$,
    distributional integration by parts gives
    \begin{align}
        \tr(\bar{H}\Sigma)
        =
        \sum_{i,j=1}^d
        \Sigma_{ij}
        \left\langle
        \partial_{ij}e,
        \gamma_{\mu,\Sigma}
        \right\rangle
        =
        -\sum_{i,j=1}^d
        \Sigma_{ij}
        \left\langle
        \partial_i e,
        \partial_j\gamma_{\mu,\Sigma}
        \right\rangle
        =
        -\int_{\R^d}
        \left\langle
        \nabla e(x),
        \Sigma\nabla\gamma_{\mu,\Sigma}(x)
        \right\rangle
        \,dx.
    \end{align}
    The Gaussian score identity
    $\Sigma\nabla\gamma_{\mu,\Sigma}(x)
        =
        -(x-\mu)\gamma_{\mu,\Sigma}(x)$
    therefore yields
    \begin{align}
        \tr(\bar{H}\Sigma)
         & =
        \int_{\R^d}
        \left\langle
        \nabla e(x),
        x-\mu
        \right\rangle
        \gamma_{\mu,\Sigma}(x)\,dx
        =
        \E\left[
            \left\langle
            \nabla e(X),
            X-\mu
            \right\rangle
            \right].
        \qedhere
    \end{align}
\end{proof}

\section{Curvature and Bregman divergence}
\label{app:bracket}
This appendix gives the general comparison underlying \cref{cor:bracket-ell}.
We state the result with a free tail parameter $\beta$, derive the corollary
for the algorithm parameters, and then prove the general result.

\begin{proposition}[General curvature comparison]
    \label{prop:bracket}
    Let $e:\R^d\to\R$ be convex and $L$-Lipschitz, and let
    $X\sim\N(\mu,\Sigma)$ with $\Sigma\succ0$.
    Write $\bar{H}=\langle D^2e,\gamma_{\mu,\Sigma}\rangle$.
    Then for any $\beta\in(0,1/2]$ and $\lambda\in(0, 1/8]$
    such that $\lambda\left(d+\log\frac2\beta\right) \leq \frac{1}{64}$,
    for all $x\in\R^d$ satisfying
    \begin{align}
        \lambda\normm{\mu-x}{\Sigma^{-1}}
        \sqrt{\log\frac2\beta}                \leq \frac{1}{64},
        \label{eq:bracket-focus}
    \end{align}
    it holds that
    \begin{align}
        \frac\lambda2\normm{\mu-x}{\bar{H}}^2
        -\E[D_e(x\,\|\,X)]
        \leq
        \frac1{16}\tr(\bar{H}\Sigma)
        +\frac{\lambda L\beta}{\sqrt{\pi}}
        \norm{\mu-x}\normm{\mu-x}{\Sigma^{-1}}.
        \label{eq:bracket}
    \end{align}
\end{proposition}

\subsection{Derivation of the corollary}
The choice $\beta=e^{-\ell}$ makes the tail remainder negligible throughout
the focus region once the covariance is bounded by $2\sigma^2I_d$.

\begin{proof}[Proof of \cref{cor:bracket-ell}]
    Apply \cref{prop:bracket} with $\beta=e^{-\ell}$.
    Since $M\geq d^{-1/2}$, $d\geq1$ and $\ell\geq640$, we have
    $0<\lambda\leq1/8$, $0<e^{-\ell}\leq1/2$, and
    \begin{align}
        \lambda\bigl(d+\log(2e^\ell)\bigr)
        \leq\frac{1+2\ell}{1000\ell^4}
        \leq\frac1{64}.
    \end{align}
    The condition on $x$ implies \eqref{eq:bracket-focus} because
    $\log(2e^\ell)\leq2\ell$.
    Moreover, the covariance bound gives
    $\norm{\mu-x}\leq\sqrt2\,\sigma\normm{\mu-x}{\Sigma^{-1}}$.
    Together with \eqref{eq:bracket-ell-focus}, this bounds the remainder by
    \begin{align}
        \frac{8\lambda e^{-\ell}}{\eps\sqrt\pi}
        \norm{\mu-x}\normm{\mu-x}{\Sigma^{-1}}
         & \leq\frac{\sigma e^{-\ell}}{512\sqrt{2\pi}\lambda\eps\ell}
        \leq\frac{d\ell^3e^{-\ell}}{\eps}
        \leq\frac{n^{-8}}{\eps}.
    \end{align}
    Here we used the parameter choices, $\ell^3\leq e^{\ell/2}$, and
    $de^{-\ell/2}=d(128nd/\delta)^{-64}\leq n^{-8}$.
\end{proof}

\subsection{Proof of the general comparison}
We first record the measure notation and the curvature representation of
expected Bregman divergence used in the proof.
Let $M$ be a positive-semidefinite matrix-valued Radon measure and let
$\nu=\tr M$ be its trace measure.
Positivity implies that every entry of $M$ is absolutely continuous with
respect to $\nu$.\alireza{Needs proof or citation.}
Hence, by the Radon--Nikodym theorem, there is a
measurable matrix-valued function $Q$ such that\alireza{Needs proof.}
\begin{align}
    M(dz)=Q(z)\nu(dz),
    \qquad
    Q(z)\succeq0,
    \qquad
    \tr Q(z)=1
    \quad\text{for $\nu$-almost every }z.
    \label{eq:matrix-measure-polar}
\end{align}
Therefore, for a nonnegative measurable function $\rho$ and a measurable vector field
$b$, we use the notation
\begin{align}
    \int_B \rho(z)b(z)^\top M(dz)b(z)
    \triangleq
    \int_B
    \rho(z)b(z)^\top Q(z)b(z)\,\nu(dz).
    \label{eq:matrix-measure-quadratic-form}
\end{align}
In particular, this integral is nonnegative, possibly with value $+\infty$.

\begin{lemma}[Averaged Bregman curvature]
    \label{lem:averaged-bregman-curvature}
    Let $h:\R^d\to\R$ be convex and Lipschitz, let $\bar{x}\in\R^d$, and let
    $Z\sim\N(0,I_d)$.  For $s\in(0,1]$, let $p_s$ denote the density of $Z_s=(1-s)\bar{x}+sZ$.
    Then
    \begin{align}
        \E[D_h(\bar{x}\,\|\,Z)]
        =
        \int_0^1\frac1s
        \int_{\R^d}
        p_s(z)(z-\bar{x})^\top D^2h(dz)(z-\bar{x})\,ds.
        \label{eq:averaged-bregman-curvature}
    \end{align}
\end{lemma}
\begin{proof}
    TODO. It is just FTC.
\end{proof}


\begin{proof}[Proof of \cref{prop:bracket}]
    Write $\gamma=\gamma_{0,I_d}$ and $\gamma_V=\gamma_{0,V}$ for $V\succ0$.
    First, we switch to the standardized case by defining
    \begin{align}
        \bar{x}=\Sigma^{-\frac12}(x-\mu),
        \qquad
        h(z)=e(\mu+\Sigma^{1/2}z),
        \quad\text{and}
        \quad
        Z=\Sigma^{-\frac12}(X - \mu) \sim \N(0,I_d)\,.
    \end{align}
    Since $\Sigma^{\frac12}\bar{x}=x-\mu$, the chain rule at
    differentiability points of $e$ gives
    $D_h(\bar{x}\,\|\,Z)=D_e(x\,\|\,X)$
    almost surely.

    The plan is to define a set $\mathcal A\subset \R^d$ and control the integrals over $\mathcal A$ and $\mathcal A^c$ separately.
    To this end, write
    \begin{align}
        \begin{aligned}
            \normm{\mu-x}{\bar{H}}^2
             & =
            \bar{x}^\top \Sigma^{\frac12}\bar{H}\Sigma^{\frac12}\bar{x}
            \\
             & =
            \int_{\R^d}\gamma(z)\bar{x}^\top D^2h(dz)\bar{x}
            \\
             & =
            \int_{\mathcal A}\gamma(z)\bar{x}^\top D^2h(dz)\bar{x}
            +
            \int_{\mathcal A^c}\gamma(z)\bar{x}^\top D^2h(dz)\bar{x}
        \end{aligned}
        \label{eq:standardized-comparator-curvature}
    \end{align}
    In step 1, we upper bound the integral over $\mathcal A$ by the sum of expected Bregman divergence $\E[D_e(x\,\|\,X)]$ and a trace term $\tr(\bar{H}\Sigma)$.
    In step 2, we upper bound the integral over $\mathcal A^c$ by a tail term that is small under the assumptions of the proposition.

    \proofstepinline{1}
    By \cref{lem:averaged-bregman-curvature}, we have
    \begin{align}
        \E[D_e(x\,\|\,X)]
        =
        \E[D_h(\bar{x}\,\|\,Z)]
        =
        \int_0^1\frac1s
        \int_{\R^d}
        p_s(z)(z-\bar{x})^\top D^2h(dz)(z-\bar{x})\,ds\,,
        \label{eq:bregman-density}
    \end{align}
    where $p_s$ denotes the density of $(1-s)\bar{x}+sZ$.
    Define the set\footnote{
        If $\bar{x}=0$, then the left-hand side of \eqref{eq:bracket} is nonpositive, so the inequality holds trivially.
    }
    \begin{align}
        \mathcal A=\left\{z:
        \norm{z}^2\leq2d\log2+4\log\frac2\beta
        \quad
        \text{and}
        \quad
        \left|\ip{\frac{\bar{x}}{\norm{\bar{x}}}}{z}\right|^2
        \leq4\log\frac{2\sqrt2}{\beta}
        \right\}.
        \label{eq:A}
    \end{align}

    \proofstepinline{1.1}
    Next we show that $\frac12\gamma(z)\leq p_s(z)$ for $s\in[1-2\lambda,1]$ and $z\in\mathcal A$.
    Since $sZ+(1-s)\bar{x}\sim\N((1-s)\bar{x},s^2I_d)$,
    we have
    \begin{align}
        \log\frac{p_s(z)}{\gamma(z)}
         & =-d\log s
        -\frac{1-s^2}{2s^2}\norm{z}^2
        +\frac{1-s}{s^2}\ip{\bar{x}}{z}
        -\frac{(1-s)^2}{2s^2}\norm{\bar{x}}^2\,.
        \label{eq:density-ratio}
    \end{align}
    For $s\in[1-2\lambda,1]$, the assumption $\lambda\leq1/8$ implies $s\geq\frac34$, and $0\leq1-s\leq2\lambda$.
    Consequently,
    \begin{align}
        \frac{1-s^2}{2s^2}
        \leq
        \frac{1-s}{s^2}
        \leq
        \frac{2\lambda}{(3/4)^2}
        =
        \frac{32}{9}\lambda
        \leq4\lambda,
        \quad\text{and}\quad
        \frac{(1-s)^2}{2s^2}
        \leq
        \frac{(2\lambda)^2}{2(3/4)^2}
        =
        \frac{32}{9}\lambda^2
        \leq4\lambda^2.
    \end{align}
    Moreover, $-d\log s\geq0$. Therefore,
    \eqref{eq:density-ratio} yields
    \begin{align}
        \log\frac{p_s(z)}{\gamma(z)}
        \geq
        -4\lambda\norm{z}^2
        -4\lambda\left|\ip{\bar{x}}{z}\right|
        -4\lambda^2\norm{\bar{x}}^2.
        \label{eq:density-ratio-lower}
    \end{align}
    Now suppose for $z\in\mathcal A$, then
    together with the assumption $\lambda\left(d+\log\frac2\beta\right)\leq \frac{1}{64}$ we have
    \begin{align}
        \lambda\norm{z}^2
         & \leq
        \lambda\left(
        2d\log2+4\log\left(\frac2\beta\right)
        \right)
        \leq
        4\lambda\left(d+\log\left(\frac2\beta\right)\right)
        \leq
        \frac{1}{16}\,,
        \label{ineq:lambda-z-squared-bound}
    \end{align}
    and, again, from $z\in\mathcal A$ we get
    \begin{align}
        \lambda\left|\ip{\bar{x}}{z}\right|
        \leq
        2\lambda\norm{\bar{x}}
        \sqrt{\log\frac{2\sqrt2}{\beta}}
        \leq
        2\lambda\norm{\bar{x}}
        \sqrt{
            \log\left(
            \frac{2}{\beta}
            \right)
            +\frac12\log2
        }
        \leq
        \sqrt5\,
        \lambda\norm{\bar{x}}\sqrt{\log\left(\frac{2}{\beta}\right)}
        \leq
        \frac{\sqrt5}{64}\,,
    \end{align}
    where in the penultimate inequality we used $
        \frac12\log2
        \leq \frac14\log\left(\frac{2}{\beta}\right)$ since $\beta<1/2$,
    and in the last inequality we used \eqref{eq:bracket-focus}.
    The last inequality also gives
    $\lambda^2\norm{\bar{x}}^2
        \leq
        \frac{1}{2^{12} \log\left(\frac{2}{\beta}\right)}
        \leq
        \frac{1}{2^{12} \log4}$.
    Substituting these three bounds into
    \eqref{eq:density-ratio-lower} gives
    \begin{align}
        \log\frac{p_s(z)}{\gamma(z)}
        \geq
        -4\left(\frac{1}{16}+\frac{\sqrt5}{64}+\frac{1}{2^{12} \log4}\right)
        \geq
        \log\frac12\,.
    \end{align}
    Therefore, we have established
    $p_s(z)
        \geq
        \frac12\gamma(z)$
    for all
    $s\in[1-2\lambda,1]$
    and
    $z\in\mathcal A$.

    \proofstepinline{1.2}
    Since the measure integrand in \eqref{eq:bregman-density} is
    nonnegative, we may restrict both integrals to
    $s\in[1-2\lambda,1]$ and $z\in\mathcal A$,
    which together with $-\log(1-2\lambda)\geq2\lambda$ gives
    \begin{align}
        \E[D_h(\bar{x}\,\|\,Z)]
         & \geq
        \int_{1-2\lambda}^1\frac1s
        \int_{\mathcal A}
        p_s(z)(z-\bar{x})^\top D^2h(dz)(z-\bar{x})\,ds
        \\
         & \geq
        \frac12
        \int_{1-2\lambda}^1\frac{ds}{s}
        \int_{\mathcal A}
        \gamma(z)(z-\bar{x})^\top D^2h(dz)(z-\bar{x})
        \\
         & =
        \frac{-\log(1-2\lambda)}{2}
        \int_{\mathcal A}
        \gamma(z)(z-\bar{x})^\top D^2h(dz)(z-\bar{x})
        \\
         & \geq
        \lambda
        \int_{\mathcal A}
        \gamma(z)(z-\bar{x})^\top D^2h(dz)(z-\bar{x}).
        \label{eq:bregman-lower}
    \end{align}
    Apply \eqref{eq:matrix-measure-polar} to $D^2h$ and write
    $\nu=\tr(D^2h)$, and $D^2h(dz)=Q(z)\nu(dz)$,
    where $Q(z)\succeq0$ and $\tr Q(z)=1$ for $\nu$-almost every $z$.
    For every positive-semidefinite matrix $G$, by the triangle inequality we have
    \begin{align}
        \frac12\bar{x}^\top G\bar{x}
        \leq
        (z-\bar{x})^\top G(z-\bar{x})+z^\top Gz.
    \end{align}
    Applying this inequality to $G=Q(z)$, multiplying by
    $\lambda\gamma(z)$, and integrating over $\mathcal A$ gives
    \begin{align}
        \frac\lambda2
        \int_{\mathcal A}
        \gamma(z)\bar{x}^\top D^2h(dz)\bar{x}
        \leq{} &
        \lambda
        \int_{\mathcal A}
        \gamma(z)(z-\bar{x})^\top D^2h(dz)(z-\bar{x})
        +
        \lambda
        \int_{\mathcal A}
        \gamma(z)z^\top D^2h(dz)z.
    \end{align}
    Moreover, for $\nu$-almost every $z$,
    $z^\top Q(z)z
        \leq
        \norm{z}^2\tr Q(z)
        =
        \norm{z}^2.$
    On $\mathcal A$, \eqref{ineq:lambda-z-squared-bound} gives $\lambda\norm{z}^2\leq1/16$.
    Hence,
    \begin{align}
        \lambda
        \int_{\mathcal A}
        \gamma(z)z^\top D^2h(dz)z
        \leq
        \frac1{16}
        \int_{\mathcal A}\gamma(z)\,\nu(dz)
        \leq
        \frac1{16}
        \int_{\R^d}\gamma(z)\tr(D^2h)(dz)
        =
        \frac1{16}\tr(\bar{H}\Sigma),
        \label{eq:trace-upper-bound}
    \end{align}
    where the last equality is followed from
    \begin{align}
        \int_{\R^d}\gamma(z)\,\tr(D^2h)(dz)
        %  & =
        % \tr(\langle D^2h, \gamma\rangle)
        =
        \tr(\langle h, D^2\gamma\rangle)
        =
        \tr\left(\E\left[h(Z)(ZZ^\top-I_d)\right]\right)
        =
        \tr(\Sigma^{\frac12}\bar{H}\Sigma^{\frac12})
        =
        \tr(\bar{H}\Sigma).
        \label{eq:standardized-trace-curvature}
    \end{align}
    Finally, combining \cref{eq:bregman-lower,eq:bregman-lower,eq:trace-upper-bound} we obtain
    \begin{align}
        \frac\lambda2
        \int_{\mathcal A}
        \gamma(z)\bar{x}^\top D^2h(dz)\bar{x}
        \leq
        \E[D_h(\bar{x}\,\|\,Z)]
        +
        \frac1{16}\tr(\bar{H}\Sigma)\,,
        \label{eq:bracket-step1}
    \end{align}
    which concludes step 1.

    \proofstepinline{2}
    It remains to bound the curvature measure on $\mathcal A^c$.  Set
    \begin{align}
        u=\frac{\bar{x}}{\norm{\bar{x}}},
        \qquad
        D^2_{\bar{x}}\triangleq\bar{x}^\top D^2h\,\bar{x}.
    \end{align}
    Since $D^2h$ is positive semidefinite, $D^2_{\bar{x}}$ is a nonnegative Radon
    measure.\alireza{we could've just used $D_{\bar{x}}^2$ from the beginning.}
    Direct comparison of Gaussian densities gives
    \begin{align}
        \begin{aligned}
            \gamma(z)\1\left\{\norm z^2>
            2d\log2+4\log\frac2\beta\right\}
             & \leq
            \frac\beta2\gamma_{2I_d}(z)\,\,\,\text{and}\,\,\,
            \\
            \gamma(z)\1\left\{|\ip{u}{z}|^2>
            4\log\frac{2\sqrt2}{\beta}\right\}
             & \leq
            \frac\beta2\gamma_{I_d+uu^\top}(z).
        \end{aligned}
        \label{eq:directional-tail}
    \end{align}
    Since $\mathcal A^c$ is contained in the union of the two events above
    and $D^2_{\bar{x}}$ is nonnegative, integrating these inequalities gives
    \begin{align}
        \int_{\mathcal A^c}\gamma(z)\,D^2_{\bar{x}}(dz)
        \leq
        \frac\beta2
        \left(
        \int_{\R^d}\gamma_{2I_d}(z)\,D^2_{\bar{x}}(dz)
        +
        \int_{\R^d}\gamma_{I_d+uu^\top}(z)\,D^2_{\bar{x}}(dz)
        \right).
        \label{eq:tail-curvature-domination}
    \end{align}
    We next bound each integral on the right-hand side of
    \eqref{eq:tail-curvature-domination}.  Fix $V\in\{2I_d,I_d+uu^\top\}$.
    Since $D^2_{\bar{x}}=\bar{x}^\top D^2h\,\bar{x}$, distributional integration by parts against the Schwartz function
    $\gamma_V$ gives
    \begin{align}
        \int_{\R^d}\gamma_V(z)\,D^2_{\bar{x}}(dz)
         & =
        -\int_{\R^d}
        \ip{\nabla h(z)}{\bar{x}}
        \ip{\nabla\gamma_V(z)}{\bar{x}}\,dz
        =
        \int_{\R^d}
        \ip{\nabla h(z)}{\bar{x}}
        \ip{V^{-1}\bar{x}}{z}
        \gamma_V(z)\,dz,
        \label{eq:wider-gaussian-ibp}
    \end{align}
    where we used
    $\nabla\gamma_V(z)=-V^{-1}z\,\gamma_V(z)$.
    At every differentiability point of $e$, the chain rule gives
    \begin{align}
        \ip{\nabla h(z)}{\bar{x}}
         & =
        \ip{\nabla e(\mu+\Sigma^{1/2}z)}
        {\Sigma^{1/2}\bar{x}}
        =
        \ip{\nabla e(\mu+\Sigma^{1/2}z)}
        {x-\mu}.
    \end{align}
    Since $e$ is $L$-Lipschitz, it follows that
    \begin{align}
        \left|\ip{\nabla h(z)}{\bar{x}}\right|
        \leq
        L\norm{x-\mu}
        \qquad\text{for almost every }z.
    \end{align}
    Therefore,
    \begin{align}
        \int_{\R^d}\gamma_V(z)\,D^2_{\bar{x}}(dz)
         & \leq
        L\norm{x-\mu}
        \int_{\R^d}
        \left|\ip{V^{-1}\bar{x}}{z}\right|
        \gamma_V(z)\,dz
        =
        L\norm{x-\mu}
        \sqrt{\frac{2}{\pi}}
        \sqrt{\bar{x}^\top V^{-1}\bar{x}}.
        \label{eq:wider-gaussian-bound}
    \end{align}
    For $V=2I_d$, we have
    $\bar{x}^\top V^{-1}\bar{x}=\frac12\norm{\bar{x}}^2$.
    For $V=I_d+uu^\top$, since $\norm u=1$,
    we have
    $V^{-1}
        =
        I_d-\frac12uu^\top$.
    Moreover, $\bar{x}=\norm{\bar{x}}\,u$, and therefore
    $\bar{x}^\top V^{-1}\bar{x}=
        \bar{x}^\top
        \left(I_d-\frac12uu^\top\right)\bar{x}
        =\frac12\norm{\bar{x}}^2$.
    Thus, for either
    $V\in\{2I_d,I_d+uu^\top\}$,
    we have
    $\sqrt{\bar{x}^\top V^{-1}\bar{x}} = \frac{\norm{\bar{x}}}{\sqrt2}$.
    Substituting this into \eqref{eq:wider-gaussian-bound} gives
    \begin{align}
        \int_{\R^d}\gamma_V(z)\,D^2_{\bar{x}}(dz)
        \leq
        \frac{L}{\sqrt\pi}
        \norm{x-\mu}\norm{\bar{x}}.
        \label{eq:wider-gaussian}
    \end{align}
    Combining \eqref{eq:tail-curvature-domination} and
    \eqref{eq:wider-gaussian} now yields
    \begin{align}
        \frac{\lambda}{2}\int_{\mathcal A^c}\gamma(z)\,D^2_{\bar{x}}(dz)
        \leq
        \frac{\beta L \lambda}{2\sqrt\pi}
        \norm{x-\mu}\norm{\bar{x}}.
        \label{eq:bracket-step2}
    \end{align}
    Finally, the definition of $\bar{x}$ gives
    $\norm{\bar{x}}^2=\normm{\mu-x}{\Sigma^{-1}}^2$,
    completing the second step.
    Plugging \eqref{eq:bracket-step1} and \eqref{eq:bracket-step2}
    into \eqref{eq:standardized-comparator-curvature}
    completes the proof.
\end{proof}

\section{Online Newton Step}
\label{app:ons}
\begin{lemma}[ONS]
    \label{lem:ons-step}
    On any round in \cref{alg:direct-hermite}, every $x\in\R^d$ satisfies
    \begin{align}
        \frac12
        \left(\normm{\mu_{t+1}-x}{\Sigma_{t+1}^{-1}}^2
        -\normm{\mu_t-x}{\Sigma_t^{-1}}^2\right)
        =
        \eta\left(
        q_t(x)-q_t(\mu_t)
        \right)
        +\frac{\eta^2}{2}\normm{g_t}{\Sigma_{t+1}}^2.
        \label{eq:ons-step}
    \end{align}
\end{lemma}
\begin{proof}
    The squared-distance expansion in \cref{sec:proof-sketch} uses only
    the update equations. The same calculation applies with $x_\star$
    replaced by any $x\in\R^d$ and with
    $\bar{q}_t,\bar{g}_t,\bar{H}_t,\bar{\Sigma}_t$ replaced by
    $q_t,g_t,H_t,\Sigma_t$,
    respectively, giving the stated identity.
\end{proof}

\section{Proof of \cref{lem:trace}}
\label{app:trace}

\begin{lemma}[Predictable spectrum]
    \label{lem:predictable-spectrum}
    Suppose that $\eps=\eps_0\leq1/2$. Then
    \begin{align}
        \|\bar{\Sigma}_{t+1}^{-1}\|_{\op}^{\frac12}
        \leq
        \sigma^{-1}+t\,,
        % \lambda\eta\Lip(e)
        \qquad(1\leq t\leq \tau_{\mathrm{cov}}),
        \label{eq:predictable-spectrum}
    \end{align}
\end{lemma}

\begin{proof}[Proof of \cref{lem:predictable-spectrum}]
    Let $L=\Lip(e)$.
    To prove \eqref{eq:predictable-spectrum}, write
    \begin{align}
        \norm{\bar{\Sigma}_{t+1}^{-1}}_{\op}
        =
        \norm{\bar{\Sigma}_{t}^{-1} + \lambda\eta \bar{H}_t}_{\op}
        \leq
        \norm{\bar{\Sigma}_{t}^{-1}}_{\op} + \lambda\eta\norm{\bar{H}_t}_{\op}\,.
    \end{align}
    Next, distributional integration by parts for every $v\in\mathbb{S}_1^d$ gives
    \begin{align}
        \norm{\bar{H}_t}_{\op}
         & =
        \sup_{\|v\|=1} v^\top \bar{H}_t v
        =
        \sup_{\|v\|=1}
        \int_{\R^d}e(x) D^2_v\gamma_{\mu_t,\Sigma_t}(x)\,dx
        =
        \sup_{\|v\|=1}
        -\int_{\R^d}\partial_ve(x)\,
        \partial_v\gamma_{\mu_t,\Sigma_t}(x)\,dx
        \\
         & =
        \sup_{\|v\|=1}
        \E_t\left[
            \ip{\nabla e(X_t)}{v}
            \ip{\Sigma_t^{-1}(X_t-\mu_t)}{v}
            \right]
        \leq
        \sup_{\|v\|=1}
        L \|v\|_{\Sigma_t^{-1}}
        =
        L \norm{\Sigma_t^{-1}}_{\op}^{1/2},
        \label{eq:H-quadratic-bound}
    \end{align}
    where the inequality is by Cauchy--Schwarz and the Lipschitz property of $e$.
    Plugging this bound into the previous inequality gives
    \begin{align}
        \begin{aligned}
            \norm{\bar{\Sigma}_{t+1}^{-1}}_{\op}
             & \leq
            \norm{\bar{\Sigma}_{t}^{-1}}_{\op} + \lambda\eta L \norm{\Sigma_t^{-1}}_{\op}^{1/2}
            \\
             & \leq
            \norm{\bar{\Sigma}_{t}^{-1}}_{\op} + \sqrt{\frac{3}{2}}\lambda\eta L\norm{\bar{\Sigma}_{t}^{-1}}_{\op}^{1/2}
            \\
             & \leq
            \left(\norm{\bar{\Sigma}_{t}^{-1}}_{\op}^{1/2} + \lambda\eta L\right)^2\,,
        \end{aligned}
        \label{eq:predictable-spectrum-bound}
    \end{align}
    where the second inequality uses \eqref{eq:covariance-stopping-time}.
    Now, by taking square roots and iterating, we get
    \begin{align}
        \norm{\bar{\Sigma}_{t+1}^{-1}}_{\op}^{1/2}
        \leq
        \norm{\bar{\Sigma}_1^{-1}}_{\op}^{1/2} + t\lambda\eta L
        = \sigma^{-1}+t\lambda\eta L
        \leq
        \sigma^{-1}
        +\frac{64t}{10^5d^{5/2}M\ell^{12}}
        \leq
        \sigma^{-1} + t\,,
    \end{align}
    which proves \eqref{eq:predictable-spectrum}.
\end{proof}

\begin{proof}[Proof of \cref{lem:trace}]
    For $t\leq\tau_{\mathrm{cov}}$, the definition of
    $\tau_{\mathrm{cov}}$ and the initialization give
    $\Sigma_t^{-1}\succeq\frac12\bar{\Sigma}_t^{-1}$, and hence
    $\Sigma_t\preceq2\bar{\Sigma}_t$.
    Since $\bar{H}_t\succeq0$, we have
    $\tr(\bar{H}_t\Sigma_t)\leq2\tr(\bar{H}_t\bar{\Sigma}_t)$,
    which proves the first inequality in \eqref{eq:trace-budget}.
    To prove the second, let $L=\Lip(e)$.
    By cyclicity of the trace,
    \begin{align}
        \eta\sum_{t=1}^{\tau_{\mathrm{cov}}}
        \tr(\bar{H}_t\bar{\Sigma}_t)
        =
        \frac1\lambda\sum_{t=1}^{\tau_{\mathrm{cov}}}
        \tr\left(
        \lambda\eta\bar{\Sigma}_t^{1/2}\bar{H}_t
        \bar{\Sigma}_t^{1/2}
        \right).
        \label{eq:normalised-trace-sum}
    \end{align}
    For every $t\leq\tau_{\mathrm{cov}}$ and every unit vector $u$, the integration-by-parts
    argument used for \eqref{eq:H-quadratic-bound}, followed by
    \eqref{eq:covariance-stopping-time}, gives
    \begin{align}
        u^\top\bar{\Sigma}_t^{1/2}\bar{H}_t
        \bar{\Sigma}_t^{1/2}u
         & \leq
        L\norm{\bar{\Sigma}_t^{1/2}u}
        \sqrt{u^\top\bar{\Sigma}_t^{1/2}
        \Sigma_t^{-1}\bar{\Sigma}_t^{1/2}u}
        \leq
        \sqrt{\frac32}\,L
        \norm{\bar{\Sigma}_t^{1/2}u}
        \leq
        \sqrt{\frac32}\,L\sigma.
        \label{eq:normalised-H}
    \end{align}
    The last inequality uses
    $\bar{\Sigma}_t\preceq\bar{\Sigma}_1=\sigma^2I_d$.
    Therefore, by the parameter choices,
    \begin{align}
        \norm{\eta\lambda\bar{\Sigma}_t^{1/2}\bar{H}_t\bar{\Sigma}_t^{1/2}}_{\op}
        \leq
        \sqrt{\frac32}\lambda\eta L\sigma
        \leq
        \frac1{1000}
        \leq
        1\,.
    \end{align}
    Returning to
    \eqref{eq:normalised-trace-sum} and applying
    $s\leq2\log(1+s)$ to eigenvalues of $\lambda\eta\bar{\Sigma}_t^{1/2}\bar{H}_t\bar{\Sigma}_t^{1/2}$
    which are in $[0,1]$, we obtain
    \begin{align}
        \eta\sum_{t=1}^{\tau_{\mathrm{cov}}}
        \tr(\bar{H}_t\bar{\Sigma}_t)
         & \leq
        \frac2\lambda\sum_{t=1}^{\tau_{\mathrm{cov}}}
        \log\det\left(
        I_d+\lambda\eta\bar{\Sigma}_t^{1/2}\bar{H}_t
        \bar{\Sigma}_t^{1/2}
        \right)                     \\
         & =
        \frac2\lambda\sum_{t=1}^{\tau_{\mathrm{cov}}}
        \log\frac{\det(\bar{\Sigma}_{t+1}^{-1})}
        {\det(\bar{\Sigma}_t^{-1})} \\
         & =
        \frac2\lambda
        \log\frac{\det(\bar{\Sigma}_{\tau_{\mathrm{cov}}+1}^{-1})}
        {\det(\bar{\Sigma}_1^{-1})} \\
         & \leq
        \frac{2d}{\lambda}
        \log\left(
        \sigma^2
        \norm{\bar{\Sigma}_{\tau_{\mathrm{cov}}+1}^{-1}}_{\op}
        \right)                     \\
         & =
        \frac{4d}{\lambda}
        \log\left(
        \sigma
        \norm{\bar{\Sigma}_{\tau_{\mathrm{cov}}+1}^{-1}}_{\op}^{1/2}
        \right)                     \\
         & \leq
        \frac{4d}{\lambda}
        \log\left(1+n\right)\,,
    \end{align}
    where the last inequality
    uses \cref{lem:predictable-spectrum},
    $\sigma\leq1$, and $\tau_{\mathrm{cov}}\leq n$.
    % Finally, by the lemma assumption and $\tau_{\mathrm{cov}}\leq n$ give
    % \begin{align}
    %     \log\left(1+\tau_{\mathrm{cov}}\lambda\eta L\sigma\right)
    %     \leq\log(1+n)\,,
    % \end{align}
    % which proves \eqref{eq:trace-budget}.
\end{proof}


\begin{proof}[Proof of \cref{lem:stopped-trace}]
    Apply the argument of the proof of \cref{lem:trace} with every upper
    summation limit replaced by $\tau-1$. Since
    $\tau-1\leq n\wedge\tau_{\mathrm{cov}}$, all covariance comparisons
    used there hold, and the final spectral bound applies with an endpoint
    at most $n$. This gives \eqref{eq:stopped-trace-budget}.
\end{proof}

\begin{proof}[Proof of \cref{lem:uniform-trace}]
    Write $h=\eps e$ and
    $\K^\circ=\{a:\sup_{z\in\K}\ip{a}{z}\leq1\}$.
    We first show that $\nabla h(x)\in5\K^\circ$ at almost every $x$.
    At a differentiability point in $\operatorname{int}\K_\eps$, we have
    $x+\eps z\in\K$ for every $z\in\K$. Therefore
    \begin{align}
        \ip{\nabla h(x)}{z}
        =\eps\ip{\nabla f(x)}{z}
        \leq f(x+\eps z)-f(x)\leq1.
    \end{align}
    At a differentiability point outside $\K_\eps$, radial homogeneity
    gives $h(tx)+2=t(h(x)+2)$ for $t$ near one, and hence
    $\ip{\nabla h(x)}{x}=h(x)+2$. Applying the supporting-hyperplane
    inequality at $(1-\eps)z$, for $z\in\K$, gives
    \begin{align}
        (1-\eps)\ip{\nabla h(x)}{z}
        \leq h((1-\eps)z)+2\leq\eps+2.
    \end{align}
    Thus $\ip{\nabla h(x)}{z}\leq5$. The boundary of $\K_\eps$ has
    Lebesgue measure zero, so the two cases establish the claimed
    inclusion almost everywhere. In particular,
    $\ip{\nabla h(x)}{w}\leq5\pi(w)$ for every $w\in\R^d$.

    Let $X\sim\N(\mu,\Sigma)$ and $W=X-\mu$. Integration by parts as
    in \cref{lem:regret-switch} now gives
    \begin{align}
        \eps\tr\left(\langle D^2e,\gamma_{\mu,\Sigma}\rangle\Sigma\right)
         & =\E\left[\ip{\nabla h(X)}{W}\right]
        \leq5\E[\pi(W)]
        \leq5\sqrt{2d}\,\sigma\E[\pi(U)]
        \leq5.
    \end{align}
    For the penultimate inequality, add an independent centered Gaussian
    with covariance $2\sigma^2I_d-\Sigma$, apply Jensen's inequality to
    the convex gauge, and use the independent radius and direction of
    the resulting isotropic Gaussian. Nonnegativity follows from
    convexity of $e$.
\end{proof}

\section{Feedback control}
\label{app:feedback}
\begin{lemma}[Single-round feedback moment]
    \label{lem:single-round-feedback}
    For every $t\leq\tau_{\mathrm{cov}}\wedge\tau_{\mathrm{shell}}$,
    \begin{align}
        \E_t[Y_t^2]\leq\frac{8\ell^2}{\eps^2}.
        \label{eq:single-round-feedback}
    \end{align}
\end{lemma}
\begin{proof}
    Write $W_t=X_t-\mu_t$ and $r_t=\pi_\eps(X_t)=1+v(X_t)$.
    The covariance and shell conditions give
    $\Sigma_t\preceq2\sigma^2I_d$ and $v(\mu_t)\leq\ell$.
    By subadditivity of the gauge and $\eps\leq1/2$,
    \begin{align}
        r_t\leq1+\ell+2\pi(W_t).
    \end{align}
    Conditionally on $\F_{t-1}$, $W_t$ is centered Gaussian with
    covariance $\Sigma_t$. Thus
    \cref{lem:gaussian-gauge-fourth-moment} and Minkowski's inequality give
    \begin{align}
        \E_t[r_t^2]^{1/2}
        \leq1+\ell+2\E_t[\pi(W_t)^2]^{1/2}
        \leq\ell+7.
    \end{align}
    The noise assumption \eqref{eq:noise} gives conditional mean zero
    and conditional second moment at most one. Since
    $Y_t=e(X_t)+r_t\xi_t$ and
    $0\leq\eps e(X_t)\leq(2+\eps)r_t\leq(5/2)r_t$,
    \begin{align}
        \eps^2\E_t[Y_t^2]
        =\eps^2\E_t\left[e(X_t)^2+r_t^2\E_t[\xi_t^2\mid X_t]\right]
        \leq7\E_t[r_t^2]
        \leq7(\ell+7)^2
        \leq8\ell^2,
    \end{align}
    where the last inequality uses $\ell\geq640$.
\end{proof}

\begin{lemma}[Gaussian gauge fourth moment]
    \label{lem:gaussian-gauge-fourth-moment}
    Let $\sigma=(M\sqrt{2d})^{-1}$ and $W\sim\N(0,\Sigma)$ with
    $0\preceq\Sigma\preceq2\sigma^2I_d$. Then
    \begin{align}
        \E[\pi(W)^4]^{1/4}\leq3.
    \end{align}
\end{lemma}
\begin{proof}
    Adding independent centered Gaussian noise to make the covariance
    $2\sigma^2I_d$ and applying Jensen's inequality to the convex gauge
    can only increase its expectation. Decomposing the resulting
    isotropic Gaussian into its independent radius and direction gives
    \begin{align}
        \E[\pi(W)]
        \leq\sqrt{2d}\,\sigma\E[\pi(U)]\leq1.
    \end{align}
    Since $\mathbb B_1\subseteq\K$, the gauge $\pi$ is $1$-Lipschitz.
    Also $\Sigma\preceq2\sigma^2I_d\preceq I_d$, so Gaussian
    concentration gives
    \begin{align}
        \Pp\left(|\pi(W)-\E[\pi(W)]|>r\right)
        \leq2e^{-r^2/2}\qquad(r>0).
    \end{align}
    Integrating this tail bound yields
    \begin{align}
        \E\left[|\pi(W)-\E[\pi(W)]|^4\right]
        \leq8\int_0^\infty r^3e^{-r^2/2}\,dr=16.
    \end{align}
    Minkowski's inequality therefore gives
    \begin{align}
        \E[\pi(W)^4]^{1/4}
        \leq\E[\pi(W)]
        +\E\left[|\pi(W)-\E[\pi(W)]|^4\right]^{1/4}
        \leq1+2=3.
    \end{align}
\end{proof}

\begin{lemma}[Weighted Gaussian tail]
    \label{lem:weighted-gaussian-tail}
    Assume the parameter choices of \cref{thm:main} with $\eps=\eps_0\leq1/2$.
    For every fixed comparator $x_{\circ}$ and every
    $t\leq\tau_{\mathrm{cov}}$ with $t<\tau_{\mathrm{shell}}$,
    \begin{align}
        \E_t\left[Y_t^2(G_t(x_{\circ})^2+1)^2
        \1_{\overline{\mathcal E}_G^{(t:t)}(x_{\circ})}\right]
        \leq\frac{e^{-\ell/4}}{\eps^2}
        \leq e^{-\ell/8}.
    \end{align}
    Consequently, for every interval
    $1\leq s\leq u\leq\tau_{\mathrm{cov}}$ with $u<\tau_{\mathrm{shell}}$,
    \begin{align}
        \sum_{t=s}^{u}\E_t\left[Y_t^2(G_t(x_{\circ})^2+1)^2
        \1_{\overline{\mathcal E}_G^{(t:t)}(x_{\circ})}\right]
        \leq e^{-\ell/16}.
    \end{align}
\end{lemma}
\begin{proof}
    If $R_t(x_{\circ})=0$, then $G_t(x_{\circ})=0$ and the tail term vanishes.
    Otherwise, $G_t(x_{\circ})$ is conditionally standard Gaussian.
    The noise assumption \eqref{eq:noise} and
    $0\leq\eps e(X_t)\leq(5/2)\pi_\eps(X_t)$ give
    \begin{align}
        \E_t[Y_t^2\mid X_t]
        =e(X_t)^2+\pi_\eps(X_t)^2\E_t[\xi_t^2\mid X_t]
        \leq\frac7{\eps^2}\pi_\eps(X_t)^2.
    \end{align}
    The covariance and shell conditions, subadditivity of the gauge,
    and \cref{lem:gaussian-gauge-fourth-moment} imply
    \begin{align}
        \E_t[\pi_\eps(X_t)^4]^{1/4}
        \leq1+v(\mu_t)+2\E_t[\pi(X_t-\mu_t)^4]^{1/4}
        \leq\ell+7,
    \end{align}
    so $\E_t[\pi_\eps(X_t)^4]^{1/2}\leq4\ell^2$.
    Gaussian moments and the Gaussian tail bound give
    \begin{align}
        \E_t[(G_t(x_{\circ})^2+1)^8]
         & \leq2^7(15!!+1),
         &
        \Pp_t\left(\overline{\mathcal E}_G^{(t:t)}(x_{\circ})\right)
         & \leq2e^{-2\ell}.
    \end{align}
    Since $G_t(x_{\circ})$ is determined by $X_t$, conditioning first on
    $X_t$ and applying Cauchy--Schwarz twice yields
    \begin{align}
         & \E_t\left[Y_t^2(G_t(x_{\circ})^2+1)^2
        \1_{\overline{\mathcal E}_G^{(t:t)}(x_{\circ})}\right]
         & \leq\frac7{\eps^2}
        \E_t[\pi_\eps(X_t)^4]^{1/2}
        \E_t[(G_t(x_{\circ})^2+1)^8]^{1/4}
        \Pp_t\left(\overline{\mathcal E}_G^{(t:t)}(x_{\circ})\right)^{1/4}
        \nonumber                                         \\
         & \qquad\leq\frac{5600\ell^2}{\eps^2}e^{-\ell/2}
        \leq\frac{e^{-\ell/4}}{\eps^2},
    \end{align}
    where the last inequality uses $\ell\geq640$.
    The parameter choices and $M\geq d^{-1/2}$ give
    $\eps^{-2}=n/(10^4d^4M^2\ell^{16})\leq n\leq e^{\ell/128}$,
    and hence $\eps^{-2}e^{-\ell/4}\leq e^{-\ell/8}$.
    Summing over the interval proves the second claim, since
    \begin{align}
        (u-s+1)e^{-\ell/8}
        \leq ne^{-\ell/8}
        \leq e^{-15\ell/128}
        \leq e^{-\ell/16}.
    \end{align}
\end{proof}

\begin{lemma}[Interval feedback concentration]
    \label{lem:interval-feedback-concentration}
    For every deterministic interval $[s,u]\subseteq[1,n]$ and deterministic $S>0$, the event $\mathcal E_Y^{(s:u)}(S)$ defined by \eqref{eq:interval-feedback-bound} satisfies
    \begin{align}
        \Pp\left(\overline{\mathcal E}_Y^{(s:u)}(S)\right)
        \leq2e^{-\ell/16}.
    \end{align}
\end{lemma}
\begin{proof}
    Set $b=16\ell^3/\eps^2$ and, for $s\leq t\leq u$, define
    \begin{align}
        I_t=\1\{t\leq\tau_{\mathrm{cov}}\wedge\tau_{\mathrm{shell}}\},
        \qquad U_t=Y_t^2\1\{|\eps Y_t|\leq4\ell^{3/2}\}.
    \end{align}
    The indicator $I_t$ is predictable and $0\leq U_t\leq b$, so the centered increments $I_t(U_t-\E_t[U_t])$ have conditional mean zero and absolute value at most $b$.
    Their accumulated conditional variance satisfies
    \begin{align}
        \sum_{t=s}^{u}\E_t\left[\left(I_t(U_t-\E_t[U_t])\right)^2\right]
        \leq b\sum_{t=s}^{u}I_t\E_t[U_t]
        \leq b\sum_{t=s}^{u}I_t\E_t[Y_t^2].
    \end{align}
    Apply \cref{lem:freedman} with variance budget $bS$ and $z=\sqrt{\ell/8}$, enlarging its linear term.
    Except on an event of probability at most $2e^{-\ell/16}$, we then have
    \begin{align}
        \sum_{t=s}^{u}I_t(U_t-\E_t[U_t])
        \leq\frac{\sqrt{2S}\,\ell^2}{\eps}+\frac{2\ell^4}{\eps^2}
    \end{align}
    whenever $\sum_{t=s}^{u}I_t\E_t[Y_t^2]\leq S$.
    Under \eqref{eq:energy-feedback-budget}, every $I_t$ on the interval equals one.
    If $\mathcal E_{\mathrm{noise}}\cap\mathcal E_{\mathrm{gauge}}$ also holds, \eqref{eq:energy-sample-gauge-bound} gives $U_t=Y_t^2$ throughout the interval.
    Finally, $\sum_{t=s}^{u}\E_t[U_t]\leq\sum_{t=s}^{u}\E_t[Y_t^2]\leq S$.
    Adding this bound to the preceding martingale bound proves \eqref{eq:interval-feedback-bound}.
\end{proof}

Throughout this section, write $\tau=\tau_{\mathrm{stable}}$. Define
\begin{align}
    \mathcal E_Y
    =\left\{
    \sum_{t=1}^\tau Y_t^2\leq20n,
    \quad
    \max_{t\leq\tau}|Y_t|\leq\frac{4\ell^{3/2}}{\eps}
    \right\}.
    \label{eq:feedback-good-event}
\end{align}
\begin{lemma}
    \label{lem:feedback}
    Suppose that $\eps=\eps_0\leq1/2$. On every round $t\leq\tau$,
    \begin{align}
        \E_t[Y_t^2]\leq\frac{8\ell^2}{\eps^2}.
        \label{eq:feedback-roundwise}
    \end{align}
    Moreover,
    \begin{align}
        \sum_{t=1}^\tau\E_t[Y_t^2]
        \leq4n+\frac{32\ell}{\eps}\operatorname{eReg}_\tau
        +\frac{32ne^{-\ell}}{\eps^2},
        \label{eq:feedback-self-bound}
    \end{align}
    and $\Pp(\mathcal E_Y)\geq1-\delta/32$.
\end{lemma}
\begin{proof}
    Write
    \begin{align}
        r_t=\pi_\eps(X_t)=1+v(X_t),
        \qquad W_t=X_t-\mu_t.
    \end{align}
    We first record the Gaussian bounds used in both parts of the proof.
    Subadditivity of the gauge gives, without any feasibility assumption
    on $\mu_t$,
    \begin{align}
        v(X_t)
         & \leq v(\mu_t)+\frac{\pi(W_t)}{1-\eps}
        \leq \ell+2\pi(W_t).
        \label{eq:feedback-v-minkowski}
    \end{align}
    The stopped covariance comparison gives
    $\Sigma_t\preceq2\sigma^2I_d\preceq I_d$. As in the proof of
    \cref{lem:uniform-trace}, Gaussian covariance comparison yields
    \begin{align}
        \E_t[\pi(W_t)]
        \leq\sqrt{2d}\,\sigma\E[\pi(U)]\leq1.
        \label{eq:feedback-minkowski-mean}
    \end{align}
    Since $\pi$ is $1$-Lipschitz, Gaussian concentration gives
    \begin{align}
        \Pp_t\left(\pi(W_t)>\E_t[\pi(W_t)]+s\right)
        \leq\exp\left(-\frac{s^2}{2\|\Sigma_t\|_{\op}}\right)
        \leq e^{-s^2/2}.
        \label{eq:feedback-minkowski-tail}
    \end{align}
    In particular,
    \begin{align}
        \Pp_t(2\pi(W_t)>s)\leq e^{-s^2/32}
        \qquad(s\geq4).
        \label{eq:feedback-v-tail}
    \end{align}
    The two-sided version of \eqref{eq:feedback-minkowski-tail} also gives
    \begin{align}
        \E_t\left[\left|\pi(W_t)-\E_t[\pi(W_t)]\right|^4\right]
                               & \leq8\int_0^\infty s^3e^{-s^2/2}\,ds=16,
        \\
        \E_t[\pi(W_t)^4]^{1/4} & \leq3,
                               & \E_t[r_t^4]^{1/4}                        & \leq \ell+7.
        \label{eq:feedback-gauge-moments}
    \end{align}
    The last two bounds use Minkowski's inequality and
    $r_t\leq1+\ell+2\pi(W_t)$.

    \proofstepinline{1}
    Conditioning on $X_t$ first, the noise assumption gives
    \begin{align}
        \E_t[Y_t^2\mid X_t]
         & =e(X_t)^2+r_t^2\E_t[\xi_t^2\mid X_t]
        \leq e(X_t)^2+r_t^2.
        \label{eq:Y2-conditional-simple}
    \end{align}
    Since $\eps e(X_t)\leq(2+\eps)r_t\leq(5/2)r_t$,
    \begin{align}
        \eps^2\E_t[Y_t^2]
        \leq7\E_t[r_t^2]
        \leq7(\ell+7)^2
        \leq8\ell^2.
    \end{align}
    This proves \eqref{eq:feedback-roundwise}, using $\ell\geq640$.
    In particular, it justifies the overshoot calculation in
    \eqref{eq:stopped-moment-budget} independently of any regret bound.

    To obtain the self-bound, start instead from
    $e(X_t)\leq1+3v(X_t)/\eps$, which gives
    \begin{align}
        \sum_{t=1}^\tau\E_t[Y_t^2]
        \leq4n+\frac{32}{\eps^2}
        \sum_{t=1}^\tau\E_t[v(X_t)^2].
        \label{eq:feedback-start}
    \end{align}
    Set $a=16\sqrt\ell$. Splitting at $\ell+a$, and then integrating
    \eqref{eq:feedback-v-tail}, yields
    \begin{align}
        \E_t[v(X_t)^2]
         & \leq(\ell+a)\E_t[v(X_t)]
        +\E_t\left[v(X_t)^2\1\{v(X_t)>\ell+a\}\right]
        \\
         & \leq2\ell\E_t[v(X_t)]
        +\left(2\ell^2+2a^2+64\right)e^{-a^2/32}
        \\
         & \leq2\ell\E_t[v(X_t)]+e^{-\ell}.
        \label{eq:feedback-v2-split}
    \end{align}
    Indeed, the tail event implies $2\pi(W_t)>a$, and on that event
    $v(X_t)^2\leq2\ell^2+2(2\pi(W_t))^2$. Also $a\leq \ell$ and
    $(2\ell^2+512\ell+64)e^{-8\ell}\leq e^{-\ell}$.

    Finally, $A_t\in\K_\eps$ implies $f(A_t)\geq f(x_\star)\geq0$.
    Therefore the definition of the extension gives the pointwise bound
    \begin{align}
        e(X_t)-e(x_\star)
         & =r_tf(A_t)-f(x_\star)+\frac2\eps v(X_t)
        \geq\frac2\eps v(X_t).
        \label{eq:extension-violation}
    \end{align}
    Combining \cref{eq:feedback-start,eq:feedback-v2-split,eq:extension-violation}
    proves \eqref{eq:feedback-self-bound}.

    \proofstepinline{2}
    It remains to prove the realized feedback bound. For this calculation
    only, put
    \begin{align}
        y_t=\eps Y_t,
        \qquad N=n\eps^2,
        \qquad b=16\ell^3.
    \end{align}
    By \eqref{eq:feedback-minkowski-tail} and the noise assumption, the
    event
    $\{\pi(W_t)\leq1+2\sqrt\ell,\ |\xi_t|\leq\sqrt\ell\}$
    has conditional probability at least $1-3e^{-\ell}$.
    On this event, $v(X_t)\leq2\ell$, $r_t\leq3\ell$, and
    \begin{align}
        |y_t|
        \leq\eps r_t+2v(X_t)+\eps r_t|\xi_t|
        \leq\frac{11}{2}\ell+\frac32\ell^{3/2}
        \leq4\ell^{3/2}.
        \label{eq:feedback-cap}
    \end{align}
    A union bound over the stopped rounds therefore costs at most
    $3ne^{-\ell}$.

    Define the predictable activity indicator $I_t=\1\{t\leq\tau\}$ and
    the truncated variable
    $U_t=y_t^2\1\{|y_t|\leq4\ell^{3/2}\}$.
    The martingale differences $I_t(U_t-\E_t[U_t])$ are bounded in
    absolute value by $b$, and \eqref{eq:stopped-moment-budget} gives
    \begin{align}
        \sum_{t=1}^n I_t\E_t[(U_t-\E_t[U_t])^2]
        \leq b\sum_{t=1}^\tau\E_t[y_t^2]
        \leq17bN.
    \end{align}
    Apply \cref{lem:freedman} with tail exponent $\ell/16$, using the
    weaker bound $\sqrt{2wz}+2bz$ when the variance is at most $w$.
    Outside an event of probability $2e^{-\ell/16}$,
    \begin{align}
        \sum_{t=1}^\tau U_t
         & \leq17N+\sqrt{34}\,\sqrt N\,\ell^2+2\ell^4
        \leq18N\leq20N.
        \label{eq:feedback-truncated-sum}
    \end{align}
    For the second inequality, use
    \begin{align}
        \frac{\ell^2}{\sqrt N}
        =\frac{1}{100d^2M\ell^6}
        \leq\frac{1}{100\ell^6}.
        \label{eq:feedback-cap-to-budget}
    \end{align}
    On the event where truncation changes no stopped sample,
    \eqref{eq:feedback-truncated-sum} is the desired bound on
    $\sum Y_t^2$. Thus
    \begin{align}
        \Pp(\overline{\mathcal E}_Y)
        \leq3ne^{-\ell}+2e^{-\ell/16}
        \leq\frac{\delta}{32},
    \end{align}
    where the final inequality follows by substituting
    $e^{-\ell/16}=(128nd/\delta)^{-8}$.
\end{proof}

\section{Proof of \cref{lem:quadratic-second-moment}}
\label{app:quadratic-second-moment}
% RESERVED BETTER PROOF WRITING
% Now, for the moment suppose that
% \begin{align}
%     \lambda R_t(x_{\circ}) \leq 1
%     \qquad\text{and}
%     \qquad
%     R_t(x_{\circ}) \leq \sqrt{2F},\,
% \end{align}
% for some $F>0$ to be specified later, then we can bound the second moment of $q_t(x_{\circ})$ as
% \begin{align}
%     \E_t\left[q_t(x_{\circ})^2\right]
%      & \leq R_t(x_{\circ})^2\E_t\left[
%         Y_t^2\left(
%         |G_t(x_{\circ})|+\frac12|G_t(x_{\circ})^2-1|
%         \right)^2\right]
%     \\
%      & \quad
%     \leq 4R_t(x_{\circ})^2\E_t\left[
%         Y_t^2\left(G_t(x_{\circ})^2+1\right)^2\right]
%     \leq 8F\E_t\left[
%         Y_t^2\left(G_t(x_{\circ})^2+1\right)^2\right].
% \end{align}
% Now, observe that the direction defining $G_t(x_{\circ})$ is then an
% $\F_{t-1}$-measurable unit vector, so the conditional distribution of
% $G_t(x_{\circ})$ given $\F_{t-1}$ is $\N(0,1)$.
% Therefore, define the event that the Gaussian random variable is bounded on the interval $[s,u]$ as
% \begin{align}
%     \mathcal E_G^{(s:u)}(x_{\circ})
%     =\left\{
%     |G_t(x_{\circ})|\leq2\sqrt\ell
%     \text{ for all }s\leq t\leq u
%     \right\}\,.
%     \label{eq:gaussian-score-event}
% \end{align}
% This event holds with high probability, as $\Pp_s\left(\overline{\mathcal E}_G^{(s:u)}(x_{\circ})\right) \leq 2(u-s+1)e^{-2\ell}$.
% Next, we split the expectation according to the
% single-round event $\mathcal E_G^{(t:t)}(x_{\circ})$ which gives
% \begin{align}
%     \E_t\left[Y_t^2(G_t(x_{\circ})^2+1)^2\right]
%      & =\E_t\left[Y_t^2(G_t(x_{\circ})^2+1)^2
%     \1_{\mathcal E_G^{(t:t)}(x_{\circ})}\right]
%     +\E_t\left[Y_t^2(G_t(x_{\circ})^2+1)^2
%     \1_{\overline{\mathcal E}_G^{(t:t)}(x_{\circ})}\right]
%     \nonumber                               \\
%      & \leq25\ell^2\E_t[Y_t^2]
%     +\E_t\left[Y_t^2(G_t(x_{\circ})^2+1)^2
%     \1_{\overline{\mathcal E}_G^{(t:t)}(x_{\circ})}\right].
%     \label{eq:gaussian-score-moment-split}
% \end{align}
% To control the remaining term, we first condition on $X_t$.
% By \eqref{eq:noise} and $\eps e(X_t)\leq(5/2)\pi_\eps(X_t)$,
% \begin{align}
%     \E_t[Y_t^2\mid X_t]
%      & =e(X_t)^2+\pi_\eps(X_t)^2\E_t[\xi_t^2\mid X_t]
%     \leq e(X_t)^2+\pi_\eps(X_t)^2
%     \leq\frac7{\eps^2}\pi_\eps(X_t)^2.
% \end{align}
% Since $G_t(x_{\circ})$ is determined by $X_t$, conditioning
% first and then applying Cauchy--Schwarz gives
% \begin{align}
%     \E_t\left[Y_t^2(G_t(x_{\circ})^2+1)^2
%         \1_{\{|G_t(x_{\circ})|>2\sqrt\ell\}}\right]
%     % \nonumber                              \\
%     %  & \quad\leq\frac7{\eps^2}\E_t\left[
%     % \pi_\eps(X_t)^2(G_t(x_{\circ})^2+1)^2
%     % \1_{\{|G_t(x_{\circ})|>2\sqrt\ell\}}\right]
%     % \nonumber                              \\
%     \leq\frac7{\eps^2}
%     \E_t[\pi_\eps(X_t)^4]^{1/2}
%     \E_t\left[(G_t(x_{\circ})^2+1)^4
%         \1_{\{|G_t(x_{\circ})|>2\sqrt\ell\}}\right]^{1/2}.
%     \label{eq:gaussian-score-tail-separation}
% \end{align}
% We apply Cauchy--Schwarz once more to extract the small probability of the Gaussian tail:
% \begin{align}
%     \E_t\left[(G_t(x_{\circ})^2+1)^4
%         \1_{\{|G_t(x_{\circ})|>2\sqrt\ell\}}\right]^{1/2}
%     \leq\E_t\left[(G_t(x_{\circ})^2+1)^8\right]^{1/4}
%     \Pp_t\left(|G_t(x_{\circ})|>2\sqrt\ell\right)^{1/4}
%     % \leq\bigl(2^8(15!!+1)\bigr)^{1/4}e^{-\ell/2}
%     \leq200e^{-\ell/2}\,,
%     \label{eq:gaussian-score-tail-factor}
% \end{align}
% where we used the fact that the sixteenth moment of a conditionally standard Gaussian is $15!!$, so $\E_t[(G_t(x_{\circ})^2+1)^8]\leq2^7(15!!+1)$.
% It remains to control $\E_t[\pi_\eps(X_t)^4]^{1/2}$ in
% \eqref{eq:gaussian-score-tail-separation}. To do so, we work with the
% conditional $L^4$ norm. Subadditivity of the gauge gives
% \begin{align}
%     \pi_\eps(X_t)
%     \leq1+v(\mu_t)+\frac{\pi(X_t-\mu_t)}{1-\eps}
%     \leq1+v(\mu_t)+2\pi(X_t-\mu_t).
% \end{align}
% Taking the conditional $L^4$ norm and applying Minkowski's inequality,
% \begin{align}
%     \E_t[\pi_\eps(X_t)^4]^{1/4}
%     \leq1+v(\mu_t)+2\E_t[\pi(X_t-\mu_t)^4]^{1/4}
%     \leq v(\mu_t) + 7.
% \end{align}
% where the last inequality follows since \cref{lem:gaussian-gauge-fourth-moment} bounds
% the norm $\E_t[\pi(X_t-\mu_t)^4]^{1/4}$ by $3$.
% We are therefore left with controlling $v(\mu_t)$.
% Projecting the mean onto $\K_\eps$ would make this term zero, but our unprojected mean can
% leave $\K_\eps$.
% To track this quantity, we define our third stopping time
% \begin{align}
%     \tau_{\mathrm{shell}}
%     =\min\{t\in[n]:v(\mu_{t+1})>\ell\}\wedge n.
%     \label{eq:shell-stopping-time}
% \end{align}
% Since $v(\mu_1)=0$, we have $v(\mu_t)\leq\ell$ for every
% $t\leq\tau_{\mathrm{shell}}$. Thus the factor needed in
% \eqref{eq:gaussian-score-tail-separation} satisfies
% \begin{align}
%     \E_t[\pi_\eps(X_t)^4]^{1/2}\leq(\ell+7)^2\leq4\ell^2
%     \qquad(t\leq\tau_{\mathrm{cov}}\wedge \tau_{\mathrm{focus}} \wedge\tau_{\mathrm{shell}}).
%     \label{eq:stopped-extension-gauge-moment}
% \end{align}
% For these rounds, substituting
% \cref{eq:gaussian-score-tail-separation,eq:gaussian-score-tail-factor,eq:stopped-extension-gauge-moment}
% into \eqref{eq:gaussian-score-moment-split} gives
% \begin{align}
%     \E_t\left[Y_t^2(G_t(x_{\circ})^2+1)^2\right]
%     \leq25\ell^2\E_t[Y_t^2]+\frac{5600\ell^2}{\eps^2}e^{-\ell/2}
%     \leq25\ell^2\E_t[Y_t^2]+\frac{e^{-\ell/4}}{\eps^2},
%     \label{eq:gaussian-score-weighted-second-moment}
% \end{align}
% where the last inequality uses $\ell\geq640$.
% Using $\E_t[D_t(x_{\circ})^2]\leq\E_t[q_t(x_{\circ})^2]$ and
% summing over an interval $[s,u]$ on which the preceding conditions hold gives
% \begin{align}
%     \sum_{t=s}^{u}\E_t[D_t(x_{\circ})^2]
%     \leq\sum_{t=s}^{u}\E_t[q_t(x_{\circ})^2]
%     \leq200F\ell^2\sum_{t=s}^{u}\E_t[Y_t^2]
%     +\frac{8F(u-s+1)}{\eps^2}e^{-\ell/4}.
%     \label{eq:quadratic-conditional-variance}
% \end{align}
% For now, suppose that we have a bound on $\sum_{t=s}^{u}\E_t[Y_t^2]$.
% We will establish the required bound separately in the global and local
% arguments.
\begin{proof}
    Fix a round $t\leq\tau_{\mathrm{cov}}\wedge\tau_{\mathrm{shell}}$
    satisfying the radius conditions. The scalar representation of
    $q_t(x_{\circ})$, together with $\lambda R_t(x_{\circ})\leq1$, gives
    \begin{align}
        |q_t(x_{\circ})|
        \leq|Y_t|R_t(x_{\circ})
        \left(|G_t(x_{\circ})|+\frac12|G_t(x_{\circ})^2-1|\right)
        \leq2|Y_t|R_t(x_{\circ})(G_t(x_{\circ})^2+1).
    \end{align}
    Since $R_t(x_{\circ})^2\leq2F$, it follows that
    \begin{align}
        \E_t[q_t(x_{\circ})^2]
        \leq8F\E_t[Y_t^2(G_t(x_{\circ})^2+1)^2].
    \end{align}
    To bound this weighted second moment, split according to the
    single-round Gaussian-score event $\mathcal E_G^{(t:t)}(x_{\circ})$.
    On this event, $(G_t(x_{\circ})^2+1)^2\leq25\ell^2$, so
    \begin{align}
        \E_t[Y_t^2(G_t(x_{\circ})^2+1)^2]
        \leq25\ell^2\E_t[Y_t^2]
        +\E_t\left[Y_t^2(G_t(x_{\circ})^2+1)^2
        \1_{\overline{\mathcal E}_G^{(t:t)}(x_{\circ})}\right].
    \end{align}
    For the remaining term, we first condition on $X_t$ to remove the
    observation noise. By \eqref{eq:noise} and
    $\eps e(X_t)\leq(5/2)\pi_\eps(X_t)$,
    \begin{align}
        \E_t[Y_t^2\mid X_t]
        =e(X_t)^2+\pi_\eps(X_t)^2\E_t[\xi_t^2\mid X_t]
        \leq\frac7{\eps^2}\pi_\eps(X_t)^2.
    \end{align}
    Because $G_t(x_{\circ})$ is determined by $X_t$, Cauchy--Schwarz
    therefore gives
    \begin{align}
        \E_t\left[Y_t^2(G_t(x_{\circ})^2+1)^2
        \1_{\overline{\mathcal E}_G^{(t:t)}(x_{\circ})}\right]
        \leq\frac7{\eps^2}
        \E_t[\pi_\eps(X_t)^4]^{1/2}
        \E_t\left[(G_t(x_{\circ})^2+1)^4
        \1_{\overline{\mathcal E}_G^{(t:t)}(x_{\circ})}\right]^{1/2}.
    \end{align}
    The second factor is small because $G_t(x_{\circ})$ is conditionally
    standard Gaussian. Applying Cauchy--Schwarz again,
    \begin{align}
        \E_t\left[(G_t(x_{\circ})^2+1)^4
        \1_{\overline{\mathcal E}_G^{(t:t)}(x_{\circ})}\right]^{1/2}
        \leq\E_t[(G_t(x_{\circ})^2+1)^8]^{1/4}
        \Pp_t\left(\overline{\mathcal E}_G^{(t:t)}(x_{\circ})\right)^{1/4}
        \leq200e^{-\ell/2}.
    \end{align}
    Here we used
    $\E_t[(G_t(x_{\circ})^2+1)^8]\leq2^7(15!!+1)$ and the Gaussian tail
    bound $\Pp_t(|G_t(x_{\circ})|>2\sqrt\ell)\leq2e^{-2\ell}$.

    For the first factor, subadditivity of the gauge and Minkowski's
    inequality yield
    \begin{align}
        \E_t[\pi_\eps(X_t)^4]^{1/4}
        \leq1+v(\mu_t)+2\E_t[\pi(W_t)^4]^{1/4}
        \leq\ell+7.
    \end{align}
    The last inequality uses the shell condition $v(\mu_t)\leq\ell$
    and \cref{lem:gaussian-gauge-fourth-moment}, which applies because
    $\Sigma_t\preceq2\sigma^2I_d$.
    Thus $\E_t[\pi_\eps(X_t)^4]^{1/2}\leq4\ell^2$.
    Combining the two factors gives
    \begin{align}
        \E_t\left[Y_t^2(G_t(x_{\circ})^2+1)^2
        \1_{\overline{\mathcal E}_G^{(t:t)}(x_{\circ})}\right]
        \leq\frac{5600\ell^2}{\eps^2}e^{-\ell/2}
        \leq\frac{e^{-\ell/4}}{\eps^2},
    \end{align}
    where the last inequality uses $\ell\geq640$.
    By the parameter choices in \cref{thm:main} and $M\geq d^{-1/2}$,
    $\eps^{-2}=n/(10^4d^4M^2\ell^{16})\leq n\leq e^{\ell/128}$.
    Hence $\eps^{-2}e^{-\ell/4}\leq e^{-\ell/8}$.
    Substituting into the weighted second-moment bound proves
    \begin{align}
        \E_t[q_t(x_{\circ})^2]
        \leq200F\ell^2\E_t[Y_t^2]
        +8Fe^{-\ell/8}.
    \end{align}
    Since $\bar{q}_t(x_{\circ})=\E_t[q_t(x_{\circ})]$,
    \begin{align}
        \E_t[D_t(x_{\circ})^2]
        =\E_t[q_t(x_{\circ})^2]-\bar{q}_t(x_{\circ})^2
        \leq\E_t[q_t(x_{\circ})^2].
    \end{align}
    This proves \eqref{eq:quadratic-second-moment}.

    For the realized square, on $\mathcal E_G^{(t:t)}(x_{\circ})$,
    \eqref{eq:quadratic-estimate-decomposition} and the radius conditions give
    \begin{align}
        q_t(x_{\circ})^2
        \leq8FY_t^2(G_t(x_{\circ})^2+1)^2
        \leq200F\ell^2Y_t^2.
        \label{eq:realized-quadratic-squares}
    \end{align}
    Jensen's inequality also gives
    \begin{align}
        D_t(x_{\circ})^2
        \leq2q_t(x_{\circ})^2+2\bar{q}_t(x_{\circ})^2
        \leq2q_t(x_{\circ})^2+2\E_t[q_t(x_{\circ})^2].
    \end{align}
    Using the conditional second-moment bound just proved yields
    \eqref{eq:realized-martingale-squares}.
\end{proof}

\section{Concentration}
\label{app:concentration}
Throughout this section, write $\tau=\tau_{\mathrm{stable}}$ and, for the
proof only, set
\begin{align}
    y_t=\eps Y_t,
    \qquad N=n\eps^2=10^4d^4M^2\ell^{16}.
    \label{eq:normalised-feedback}
\end{align}
Conditionally on $\F_{t-1}$, the standardized sample
$Z_t=\Sigma_t^{-1/2}(X_t-\mu_t)$ has law $\N(0,I_d)$.
Every stopped sum below is extended by zeros after its stopping time.
In particular, each activity indicator multiplying a centered increment
is predictable.

The following two-sided form of Freedman's inequality follows from
\citet[Theorem~3.10]{bercu2015concentration} applied to a martingale and its
negative.
\begin{lemma}[Freedman's inequality]
    \label{lem:freedman}
    Let $(D_t)_{t=1}^n$ be a square-integrable
    martingale-difference sequence such that, for some $b>0$,
    $|D_t|\leq b$ almost surely for every $t$. Then, for every $s,z>0$,
    \begin{align}
        \Pp\left(
        \left|\sum_{t=1}^n D_t\right|
        \geq sz+\frac{bz^2}{3},
        \quad
        \sum_{t=1}^n\E_t[D_t^2]\leq s^2
        \right)
        \leq
        2\exp\left(-\frac{z^2}{2}\right).
    \end{align}
\end{lemma}

% Preserve \label across autonum's handling of the repeated align environment.
\begingroup
\LetLtxMacro{\restatementlabel}{\label}
\freedmannonnegative*
\GlobalLetLtxMacro{\label}{\restatementlabel}
\endgroup

The next lemma is the two-sided form of
\citet[Theorem~3.26]{bercu2015concentration}.
% Preserve \label across autonum's handling of the repeated align environment.
\begingroup
\LetLtxMacro{\restatementlabel}{\label}
\selfnormalisedmartingale*
\GlobalLetLtxMacro{\label}{\restatementlabel}
\endgroup

\begin{lemma}[Weighted Gaussian score moments]
    \label{lem:gaussian-score-localisation}
    Suppose that $\eps=\eps_0\leq1/2$. On a round $t\leq\tau$, let
    $h_t$ be any $\F_{t-1}$-measurable unit vector and put
    $G_t=\ip{h_t}{Z_t}$. Then
    \begin{align}
        \E_t[y_t^2(G_t^2+1)^2]
        \leq25\ell^2\E_t[y_t^2]+e^{-\ell/4}.
        \label{eq:H-score-second-moment}
    \end{align}
\end{lemma}
\begin{proof}
    On $\mathcal A_t=\{|G_t|\leq2\sqrt\ell\}$, we have
    $(G_t^2+1)^2\leq25\ell^2$. Thus it suffices to bound the integral on
    $\overline{\mathcal A_t}$. By \eqref{eq:Y2-conditional-simple} and the bound
    $\eps e(X_t)\leq(5/2)r_t$ used in the feedback proof,
    \begin{align}
        \E_t[y_t^2(G_t^2+1)^2\1_{\overline{\mathcal A_t}}]
         & \leq7\E_t[r_t^2(G_t^2+1)^2\1_{\overline{\mathcal A_t}}]
        \\
         & \leq7\E_t[r_t^4]^{1/2}
        \E_t[(G_t^2+1)^4\1_{\overline{\mathcal A_t}}]^{1/2}.
    \end{align}
    The first factor is at most $(\ell+7)^2$ by
    \eqref{eq:feedback-gauge-moments}. For the second, Cauchy--Schwarz and
    the standard Gaussian sixteenth moment give
    \begin{align}
        \E_t[(G_t^2+1)^4\1_{\overline{\mathcal A_t}}]^{1/2}
         & \leq\E_t[(G_t^2+1)^8]^{1/4}\Pp_t(\overline{\mathcal A_t})^{1/4}
        \\
         & \leq\bigl(2^8(15!!+1)\bigr)^{1/4}e^{-\ell/2}
        \leq200e^{-\ell/2}.
    \end{align}
    Consequently the omitted integral is at most
    \begin{align}
        1400(\ell+7)^2e^{-\ell/2}\leq e^{-\ell/4},
    \end{align}
    using $\ell\geq640$. No independence between the noise and the
    sampled action is used.
\end{proof}

We will truncate only in the proof, never in the algorithm. Define
\begin{align}
    C_t=\{|y_t|\leq4\ell^{3/2},\ \norm{Z_t}^2\leq3(d+\ell)\}.
    \label{eq:common-truncation}
\end{align}
The feedback cap and the chi-square bound below give
$\Pp_t(\overline{C_t})\leq4e^{-\ell}$ on a stopped round.
The next estimate keeps the linear Gaussian score at its natural
second-moment scale, rather than bounding it by the squared Hessian score.
\begin{lemma}[Quadratic score bounds]
    \label{lem:quadratic-score-bounds}
    Let $x$ be predictable at round $t\leq\tau$, and suppose that
    $\frac12R_t(x)^2\leq D\leq32B\ell$. Put
    \begin{align}
        Q_t & =\eps q_t(x),
            & R_t           & =\normm{x-\mu_t}{\Sigma_t^{-1}},
            & G_t           & =\ip{\Sigma_t^{-1/2}(x-\mu_t)/R_t}{Z_t},
    \end{align}
    taking any predictable unit direction when $R_t=0$. Then
    \begin{align}
        \E_t[Q_t^2]
         & \leq10D\ell\,\E_t[y_t^2]+De^{-\ell/8},
        \label{eq:q-predictable-second-moment}
        \\
        |Q_t|\1_{C_t\cap\{|G_t|\leq2\sqrt\ell\}}
         & \leq12\sqrt D\,\ell^2,
        \label{eq:q-score-bound}
        \\
        \left|\E_t\left[Q_t\1_{\overline{C_t\cap\{|G_t|\leq2\sqrt\ell\}}}\right]\right|
         & \leq32\sqrt D\,\ell^2 e^{-\ell/2}.
        \label{eq:q-truncation-bias}
    \end{align}
\end{lemma}
\begin{proof}
    The estimators give the exact representation
    \begin{align}
        Q_t=y_tL_t,
        \qquad
        L_t=R_tG_t+\frac{\lambda R_t^2}{2}(G_t^2-1).
        \label{eq:q-score-representation}
    \end{align}
    The parameter calculation in
    \eqref{eq:main-stability-bracket-focus} gives
    $\lambda R_t\sqrt\ell\leq1/64$. Hence, on $|G_t|\leq2\sqrt\ell$,
    \begin{align}
        |L_t|
         & \leq R_t\sqrt\ell\left(2+\frac{4+1/\ell}{128}\right)
        \leq\frac{21}{10}R_t\sqrt\ell,
         & L_t^2                                                & \leq10D\ell.
        \label{eq:q-local-score-square}
    \end{align}
    On the complement we use the unconditional bound
    $L_t^2\leq6D(G_t^2+1)^2$. The tail calculation in the proof of
    \cref{lem:gaussian-score-localisation} therefore gives
    \begin{align}
        \E_t[Q_t^2]
        \leq10D\ell\,\E_t[y_t^2]+6De^{-\ell/4}
        \leq10D\ell\,\E_t[y_t^2]+De^{-\ell/8}.
    \end{align}
    Combining \eqref{eq:q-local-score-square} with the cap in $C_t$
    proves \eqref{eq:q-score-bound}.

    Finally, the scalar truncation fails with conditional probability at
    most $4e^{-\ell}+2e^{-2\ell}\leq6e^{-\ell}$.
    By \eqref{eq:feedback-roundwise}, the second moment just proved is
    at most $81D\ell^3$. Cauchy--Schwarz gives
    \begin{align}
        \left|\E_t\left[Q_t\1_{\overline{C_t\cap\{|G_t|\leq2\sqrt\ell\}}}\right]\right|
        \leq\sqrt{486D}\,\ell^{3/2}\,e^{-\ell/2}
        \leq32\sqrt D\,\ell^2 e^{-\ell/2}.
    \end{align}
\end{proof}

The following chi-square upper-tail inequality follows from
\citet[Lemma~1 and Eq.~(4.3)]{laurent2000adaptive}.
\begin{lemma}[Chi-square upper-tail bound]
    \label{lem:chi-square-upper-tail}
    Let $V\sim\chi_d^2$. Then, for every $x>0$,
    \begin{align}
        \Pp\left(
        V\geq d+2\sqrt{dx}+2x
        \right)
        \leq
        \exp(-x).
        \label{eq:chi-square-upper-tail}
    \end{align}
\end{lemma}

\begin{proof}[Proof of \cref{prop:concentration}]
    For reference, the following weaker form of \cref{lem:freedman}
    will suffice: centered increments bounded by $b$ with total
    conditional variance at most $w$ satisfy
    \begin{align}
        \Pp\left(\left|\sum_tD_t\right|>\sqrt{2wz}+2bz\right)
        \leq2e^{-z}.
        \label{eq:freedman-convenient}
    \end{align}
    All applications below have deterministic bounds on $b$ and $w$;
    random directional scales are handled by predictable truncation at
    deterministic dyadic levels.

    \proofstep{1}{Feedback and quadratic concentration}
    By \cref{lem:feedback}, $\mathcal E_Y$ fails with probability at most
    $\delta/32$. The chi-square bound and a union bound give
    \begin{align}
        \Pp(\overline{\mathcal E}_Z)\leq ne^{-\ell},
        \qquad
        \mathcal E_Z=\left\{
        \max_{t\leq\tau}\norm{Z_t}^2\leq3(d+\ell)\right\}.
        \label{eq:Z-norm-union}
    \end{align}
    This proves \eqref{eq:stopped-increment-bounds} on
    $\mathcal E_Y\cap\mathcal E_Z$.

    For the quadratic sum take $Q_t=\eps q_t(x_\star)$ and
    $D=B$ in \cref{lem:quadratic-score-bounds}. Truncate $Q_t$ on
    $C_t\cap\{|G_t|\leq2\sqrt\ell\}$ and subtract its conditional mean.
    Including the predictable indicator $\1\{t\leq\tau\}$, the resulting
    increments have absolute value at most $24\sqrt B\,\ell^2$ and total
    conditional variance at most
    \begin{align}
        10B\ell\sum_{t=1}^\tau\E_t[y_t^2]+Bn e^{-\ell/8}
        \leq171BN\ell.
        \label{eq:global-q-variance}
    \end{align}
    Here we used \eqref{eq:stopped-moment-budget} and
    $ne^{-\ell/8}\leq1\leq N\ell$.
    The sum of the truncation biases is at most $n^{-8}$: indeed
    $B,\ell\leq n$ by \eqref{eq:active-horizon}, and
    \begin{align}
        32n\sqrt B\,\ell^2 e^{-\ell/2}
        \leq32n^{7/2}(128nd/\delta)^{-64}\leq n^{-8}.
        \label{eq:scalar-total-bias}
    \end{align}
    Applying \eqref{eq:freedman-convenient} with $z=\ell/16$, and
    restoring the conditional means, gives
    \begin{align}
        \left|\eps\sum_{t=1}^\tau
        (q_t(x_\star)-\bar{q}_t(x_\star))\right|
         & \leq\sqrt{\frac{171}{8}}\sqrt{BN}\,\ell
        +3\sqrt B\,\ell^3+n^{-8}
        \leq8\sqrt{BN}\,\ell.
    \end{align}
    The last inequality uses \eqref{eq:feedback-cap-to-budget}.
    Besides $\overline{\mathcal E}_Y\cup\overline{\mathcal E}_Z$, the exceptional probability
    is at most $2ne^{-2\ell}+2e^{-\ell/16}$, accounting respectively for
    the scalar score maximum and Freedman's inequality.
    Dividing by $\eps$ proves \eqref{eq:q-concentration}.

    \proofstep{2}{Gradient energy}
    The exact identity
    $\normm{g_t}{\Sigma_t}^2=Y_t^2\norm{Z_t}^2$
    gives, on $\mathcal E_Y\cap\mathcal E_Z$,
    \begin{align}
        \sum_{t=1}^\tau\normm{g_t}{\Sigma_t}^2
        \leq3(d+\ell)\sum_{t=1}^\tau Y_t^2
        \leq80dn\ell.
    \end{align}
    We used $3(d+\ell)\leq4d\ell$, which holds for $\ell\geq3$.

    \proofstep{3}{Hessian concentration}
    Put $T=\eps\sum_{t=1}^\tau(H_t-\bar{H}_t)$ and
    $A=\bar{\Sigma}_{\tau+1}^{-1}$. Let $\mathcal U$ be a deterministic
    $(512dn^2)^{-1}$-net of the Euclidean unit sphere. The volumetric
    covering bound gives
    \begin{align}
        \log|\mathcal U|
        \leq d\log(1+1024dn^2)\leq\frac{d\ell}{64}.
        \label{eq:H-net-size}
    \end{align}
    For $u\in\mathcal U$, define
    \begin{align}
        a_{t,u}=u^\top\Sigma_t^{-1}u,
        \qquad
        G_{t,u}=
        \ip{\Sigma_t^{-1/2}u/\sqrt{a_{t,u}}}{Z_t}.
        \label{eq:H-standardisation}
    \end{align}
    Then
    \begin{align}
        \eps u^\top H_tu
        =a_{t,u}y_t(G_{t,u}^2-1).
        \label{eq:H-direction}
    \end{align}
    By covariance comparison and \cref{lem:predictable-spectrum},
    \begin{align}
        1\leq a_{t,u}\leq12dn^2,
        \qquad
        \kappa(A)\leq3dn^2.
        \label{eq:concentration-spectrum}
    \end{align}
    To verify these bounds, use
    $\Sigma_t^{-1}\succeq\sigma^{-2}I_d/2\succeq I_d$,
    $\Sigma_t^{-1}\preceq3\bar{\Sigma}_t^{-1}/2$, and
    \begin{align}
        \|\bar{\Sigma}_{t+1}^{-1}\|_{\op}
        \leq(\sigma^{-1}+n)^2
        \leq2(2d+n^2)\leq6dn^2,
        \qquad A\succeq2I_d.
    \end{align}

    Set $J=\lceil\log_2(12dn^2)\rceil$ and $b_j=2^j$ for
    $0\leq j\leq J$. For a fixed $(u,j)$, truncate the score in
    \eqref{eq:H-direction} on $C_t$, center it, and multiply by
    the predictable indicator
    $\1\{t\leq\tau,\ a_{t,u}\leq b_j\}$.
    Since $\norm{Z_t}^2+1\leq4d\ell$ on $C_t$, these centered increments
    have absolute value at most $32b_jd\ell^{5/2}$.
    By \cref{lem:gaussian-score-localisation} their total conditional
    variance is at most
    \begin{align}
        b_j^2\sum_{t=1}^\tau
        \left(25\ell^2\E_t[y_t^2]+e^{-\ell/4}\right)
        \leq426b_j^2N\ell^2.
        \label{eq:H-directional-predictable}
    \end{align}
    Apply \eqref{eq:freedman-convenient} with $z=d\ell/16$. The resulting
    bound is
    \begin{align}
        b_j\left(
        \sqrt{\frac{426}{8}}\sqrt{dN}\,\ell^{3/2}
        +4d^2\ell^{7/2}\right)
        \leq8b_j\sqrt{dN}\,\ell^{3/2}.
        \label{eq:H-peeled-bound}
    \end{align}
    For the last inequality, divide the second term by
    $b_j\sqrt{dN}\ell^{3/2}$ and use
    \begin{align}
        \frac{4d^{3/2}\ell^2}{\sqrt N}
        =\frac{1}{25\sqrt d\,M\ell^6}
        \leq\frac{1}{25\ell^6}.
    \end{align}

    We must also restore the untruncated conditional means.
    Cauchy--Schwarz, \cref{lem:gaussian-score-localisation}, and
    $\Pp_t(\overline{C_t})\leq4e^{-\ell}$ give the per-round bias bound
    \begin{align}
        \left|\E_t[a_{t,u}y_t(G_{t,u}^2-1)\1_{\overline{C_t}}]\right|
        \leq40a_{t,u}\ell^2 e^{-\ell/2}.
    \end{align}
    Since $a_{t,u}\leq(3/2)u^\top Au$ before stopping, the sum of these
    biases is at most $n^{-8}u^\top Au$.
    On the dyadic level matching $a_u^\star=\max_{t\leq\tau}a_{t,u}$,
    no directional summand is removed and
    $b_j\leq2a_u^\star\leq3u^\top Au$. Hence, on
    $\mathcal E_Y\cap\mathcal E_Z$, \eqref{eq:H-peeled-bound} gives
    \begin{align}
        |u^\top Tu|
        \leq30\sqrt{dN}\,\ell^{3/2}u^\top Au
        \qquad(u\in\mathcal U).
        \label{eq:H-net}
    \end{align}
    The dyadic level is selected only after applying the inequalities
    simultaneously to all deterministic levels.

    To pass from the net to all directions, write
    $r=(512dn^2)^{-1}$ and $c=30\sqrt{dN}\ell^{3/2}$.
    For any unit vector $w$, choose $u\in\mathcal U$ with
    $\norm{w-u}\leq r$. First \eqref{eq:H-net} implies
    $\|T\|_{\op}\leq2c\|A\|_{\op}$, since $2r\leq1/2$.
    Applying the approximation once more gives
    \begin{align}
        |w^\top Tw|
         & \leq c\,w^\top Aw+6cr\|A\|_{\op}
        \leq c(1+6r\kappa(A))w^\top Aw
        \leq2c\,w^\top Aw.
    \end{align}
    Thus $-100\sqrt{dN}\ell^{3/2}A\preceq T
        \preceq100\sqrt{dN}\ell^{3/2}A$; dividing by $\eps$ proves
    \eqref{eq:H-concentration}.

    Finally, \eqref{eq:H-net-size} and $\log(J+1)\leq\ell/128$ bound
    the Hessian failure probability by $2e^{-5d\ell/128}$.
    Summing the failure probabilities from all three steps gives
    \begin{align}
        \frac\delta{32}+ne^{-\ell}+2ne^{-2\ell}
        +2e^{-\ell/16}+2e^{-5d\ell/128}
        <\frac\delta4.
    \end{align}
    This proves the simultaneous probability statement.
\end{proof}

\begin{lemma}[Short-interval concentration]
    \label{lem:local-concentration}
    For each deterministic interval $[s,s+L-1]\subseteq[n]$ with
    $1\leq L\leq k$, consider the two comparator--threshold pairs
    \begin{align}
        x=\mu_s,\quad D=D_L:=16\alpha L\ell,
        \qquad\text{or}\qquad
        x=x_\star,\quad D=B.
        \label{eq:local-threshold}
    \end{align}
    Retain summands only through $\tau$ and through the first update
    whose ending potential $\frac12R_{t+1}(x)^2$ exceeds $D$. If the starting
    potential exceeds $D$, retain no summands. There is an event
    $\mathcal E_{\mathrm{loc}}$ of probability at least $1-\delta/4$
    such that, on $\mathcal E_{\mathrm{loc}}\cap\mathcal E_Y\cap\mathcal E_Z$,
    simultaneously for all these stopped sums,
    \begin{align}
        \left|\eps\sum
        \left(q_t(x)-\bar{q}_t(x)\right)\right|
        \leq8\sqrt{DL}\,\ell^3+1.
        \label{eq:local-q-concentration}
    \end{align}
\end{lemma}
\begin{proof}
    The comparator $\mu_s$ is $\F_{s-1}$-measurable, so each retained
    score has a predictable direction. Also
    $D_L\leq16(B+\alpha)\ell\leq32B\ell$ by $L\leq k$.
    Before the auxiliary distance stop, \cref{lem:quadratic-score-bounds}
    therefore applies with the deterministic threshold $D$.
    Truncate and center the normalized scalar score as in
    \texttt{Step 1} of the preceding proof. By
    \eqref{eq:feedback-roundwise}, its total conditional variance is
    at most $81DL\ell^3$, and its centered increments are bounded by
    $24\sqrt D\,\ell^2$. The total truncation bias is at most $n^{-8}$,
    as in \eqref{eq:scalar-total-bias}, because $D\leq32B\ell\leq n$.
    Applying \eqref{eq:freedman-convenient} with $z=\ell/16$ gives
    \begin{align}
        \left|\eps\sum(q_t(x)-\bar{q}_t(x))\right|
         & \leq\sqrt{\frac{81}{8}}\sqrt{DL}\,\ell^2
        +3\sqrt D\,\ell^3+n^{-8}
        \leq8\sqrt{DL}\,\ell^3+1.
    \end{align}
    The activity indicator includes the candidate first crossing,
    since its starting potential is still at most $D$.
    Thus no bound on a future iterate has been assumed in this application.

    There are at most $2n^2$ interval--comparator pairs and at most
    $2n^3$ scalar direction--round pairs. A union bound over the
    Freedman events and the scalar score maxima gives
    \begin{align}
        \Pp(\overline{\mathcal E}_{\mathrm{loc}})
        \leq4n^2e^{-\ell/16}+4n^3e^{-2\ell}
        <\frac\delta4.
    \end{align}
    The common truncation agrees with all stopped samples on
    $\mathcal E_Y\cap\mathcal E_Z$.
\end{proof}

\section{Proof of \cref{lem:confinement}}
\label{app:confinement}

Write $\tau=\tau_{\mathrm{stable}}$ and work on the intersection of
$\mathcal E_{\mathrm{loc}}$ from \cref{lem:local-concentration} with the
event of \cref{prop:concentration}. Suppose that the covariance and focus
conditions also hold at round $\tau$. Thus every update through $\tau$ is
accepted, $\frac12R_t(x_\star)^2\leq B$ for $t\leq\tau+1$, and
\begin{align}
    \Sigma_{t+1}\preceq3\Sigma_t,
    \qquad\Sigma_t\preceq I_d
    \qquad(t\leq\tau).
    \label{eq:local-covariance}
\end{align}
For this proof write $h=\eps e$. This is only a rescaling in the analysis;
the feedback and the estimators are unchanged. We first control the
regret on short intervals against their starting means, and then convert
these interval bounds into a bound on $h(\mu_t)$.

\proofstep{1}{Distance on a short interval}
Fix a deterministic interval $[s,s+L-1]\subseteq[1,\tau]$ with $L\leq k$.
Recall $D_L=16\alpha L\ell$. We claim that
\begin{align}
    \max_{s\leq j\leq s+L}\frac12R_j(\mu_s)^2\leq D_L.
    \label{eq:local-distance}
\end{align}
Stop this interval at the first update whose ending potential exceeds
$D_L$. Before that update, $\frac12R_j(\mu_s)^2\leq D_L\leq32B\ell$.
Since $h(\mu_s)\leq\eps+(2+\eps)\ell\leq3\ell$ and $h\geq0$,
\eqref{eq:main-q-upper} and \cref{lem:uniform-trace} give
\begin{align}
    \eps \bar{q}_j(\mu_s)
     & \leq h(\mu_s)-\E_j[h(X_j)]
    +\frac{17}{16}\eps\tr(\bar{H}_j\Sigma_j)+n^{-8}
    \leq4\ell.
    \label{eq:local-q-upper}
\end{align}
The increment bounds in \eqref{eq:stopped-increment-bounds} and
\eqref{eq:local-covariance} also give the roundwise energy bound
\begin{align}
    \frac{\eta^2}{2}\normm{g_j}{\Sigma_{j+1}}^2
     & \leq\frac{3\alpha^2}{2}(\eps Y_j)^2\norm{Z_j}^2
    \leq100\alpha^2d\ell^4.
    \label{eq:local-energy}
\end{align}
Indeed, the product of the numerical bounds is
$(3/2)\cdot16\cdot4=96\leq100$.

The ONS identity, $R_s(\mu_s)=0$, and the already-stopped concentration
bound \eqref{eq:local-q-concentration} now imply that any candidate first
crossing satisfies
\begin{align}
    \frac12R_{j+1}(\mu_s)^2
     & \leq4\alpha L\ell
    +\alpha\left(8\sqrt{D_LL}\,\ell^3+1\right)
    +100\alpha^2Ld\ell^4
    <\frac{D_L}{2}.
    \label{eq:local-distance-bootstrap}
\end{align}
To check the strict inequality, divide by $D_L=16\alpha L\ell$:
\begin{align}
    \frac14
    +2\sqrt{\alpha \ell}\,\ell^2
    +\frac{1}{16L\ell}
    +\frac{25}{4}\alpha d\ell^3
     & =\frac14
    +\frac{2\sqrt{2/25}}{\sqrt d\,\ell^{3/2}}
    +\frac{1}{16L\ell}
    +\frac{1}{2\ell^5}
    <\frac12.
\end{align}
This contradicts a first crossing, proving \eqref{eq:local-distance}.
The auxiliary distance stop was present in the concentration statement
before this argument, so the claim is not being used to justify its own
martingale bound.

\proofstep{2}{Regret to the starting mean}
For a comparator with $\frac12R_j(x)^2\leq32B\ell$, Jensen's inequality,
\eqref{eq:main-q-upper}, and \cref{lem:uniform-trace} give
\begin{align}
    h(\mu_j)-h(x)
    \leq-\eps \bar{q}_j(x)+6+n^{-8}.
    \label{eq:local-regret-switch}
\end{align}
By \texttt{Step 1}, the auxiliary distance stop can now be removed on
$[s,s+L-1]$ with comparator $\mu_s$. Summing
\eqref{eq:local-regret-switch} and the ONS identity, and dropping the
nonnegative ending potential, gives
\begin{align}
    \sum_{j=s}^{s+L-1}\bigl(h(\mu_j)-h(\mu_s)\bigr)
     & \leq6L+Ln^{-8}
    +100\alpha Ld\ell^4
    +8\sqrt{D_LL}\,\ell^3+1
    \leq10L.
    \label{eq:interval-regret}
\end{align}
For the last inequality, use
\begin{align}
    100\alpha d\ell^4 & =\frac8{\ell^4}<1,
                      & 8\sqrt{D_L/L}\,\ell^3
                      & =\frac{32\sqrt{2/25}}{\sqrt{d\ell}}<1,
    \label{eq:local-parameter-checks}
\end{align}
together with $n^{-8}\leq1$ and $1\leq L$.
The crucial point is that the right-hand side of
\eqref{eq:interval-regret} is a numerical constant times $L$, not a
quantity proportional to $dL$.

\proofstep{3}{Average over the last window}
Fix $T\leq\tau$, and put $m=\min\{k,T\}$ and $s=T-m+1$. We claim that
\begin{align}
    \frac1m\sum_{j=s}^T h(\mu_j)\leq11.
    \label{eq:window-average}
\end{align}
If $T<k$, use \eqref{eq:interval-regret} on $[1,T]$ and
$h(\mu_1)=\eps f(0)\leq1/2$.
If $T\geq k$, instead use the comparator $x_\star$ on $[s,T]$.
Its potential is at most $B$ throughout the interval, including its
ending time, by the stopped focus bound and the assumption at round
$\tau$. Thus \cref{lem:local-concentration} applies without an auxiliary
exit. The ONS identity, \eqref{eq:local-energy}, and
\eqref{eq:local-regret-switch} give
\begin{align}
    \frac1k\sum_{j=s}^T\bigl(h(\mu_j)-h(x_\star)\bigr)
     & \leq\frac{B}{\alpha k}+6+n^{-8}
    +100\alpha d\ell^4
    +8\sqrt{B/k}\,\ell^3+\frac1k
    \leq10.
\end{align}
Here $B/(\alpha k)\leq1$, $\sqrt{B/k}\leq\sqrt\alpha$, and
\begin{align}
    8\sqrt\alpha\,\ell^3
    =\frac{8\sqrt{2/25}}{\sqrt d\,\ell}<1.
\end{align}
Also \eqref{eq:local-parameter-checks} bounds the energy term, and
$n^{-8}+1/k\leq1$ because $n,k\geq2$.
Finally $h(x_\star)\leq1/2$, proving \eqref{eq:window-average} in both
cases.

\proofstep{4}{From the average to the ending mean}
For any real sequence $a_s,\ldots,a_T$ of length $m$, telescoping the
averages of its last $r$ terms gives
\begin{align}
    a_T
    =\frac1m\sum_{j=s}^T a_j
    +\sum_{r=1}^{m-1}\frac1{r(r+1)}
    \sum_{j=T-r+1}^T(a_j-a_{T-r}).
    \label{eq:suffix-average-identity}
\end{align}
Apply this identity with $a_j=h(\mu_j)$. Each inner sum is the regret on
an interval of length $r+1\leq k$ against its starting mean, with the
zero first summand omitted. Hence \eqref{eq:interval-regret} and
\eqref{eq:window-average} imply
\begin{align}
    h(\mu_T)
     & \leq11+10\sum_{r=1}^{m-1}\frac1r
    \leq21+10\log n
    \leq\frac\ell2
    \qquad(T\leq\tau).
    \label{eq:last-iterate-extension}
\end{align}
The last inequality uses $\log n\leq\ell/128$ and $\ell\geq640$.
Since $h(x)\geq2v(x)$, we have $v(\mu_T)\leq\ell/4$ for every
$T\leq\tau$.

It remains to include the last update. By \eqref{eq:local-covariance},
$\|\Sigma_{\tau+1}\Sigma_\tau^{-1/2}\|_{\op}\leq3$.
Thus the update and the stopped increment bounds give
\begin{align}
    \norm{\mu_{\tau+1}-\mu_\tau}
     & \leq3\alpha|\eps Y_\tau|\norm{Z_\tau}
    \leq24\alpha \sqrt d\,\ell^2.
\end{align}
Since $v$ is $2$-Lipschitz,
\begin{align}
    v(\mu_{\tau+1})
     & \leq\frac\ell4+48\alpha \sqrt d\,\ell^2
    =\frac\ell4+\frac{96}{25\sqrt d\,\ell^6}
    <\frac\ell4+1<\frac\ell2.
\end{align}
This proves \cref{lem:confinement}, including the final mean.

\end{document}
